\chapter{Structure theory and invariant forms}\label{ch:L04}
\begingroup\setlength{\emergencystretch}{3em}

\providecommand{\lfourvac}{\mathbf{1}}
\providecommand{\lfourV}{V}
\providecommand{\lfourm}{\mathfrak{m}}
\providecommand{\lfourrad}{\operatorname{rad}}
\providecommand{\lfourQP}{\mathcal{QP}}
\providecommand{\lfourvc}{\varepsilon}
\providecommand{\lfourH}{\mathrm{H}}

This chapter documents the release of the LeanVOA library named in the front matter. Every Lean quotation below is checked by \texttt{scripts/textbook/sync.py}
to occur verbatim (up to whitespace and comments) in the source files of that release. Checks that were run at an earlier state of the library and not re-run for this edition are marked as such. The chapter covers the thirty-four files assigned to it in
\texttt{Notes/Textbook/assignment.txt}: the files of \leanpath{VOA/\allowbreak Core/} called
\texttt{V1LieAlgebra}, \texttt{LieAlgebraV1}, \texttt{Automorphism}, \texttt{InnerAutomorphism},
\texttt{CommutativeRing}, \texttt{QuotientVOA}, \texttt{VOASubalgebraGenerated}, \texttt{TensorVOA},
\texttt{TensorLieAction}, \texttt{InternalGradingOperator}, \texttt{SimpleImage},
\texttt{SingularVector}, \texttt{ScreeningKernel}, \texttt{BRSTCohomology},
\texttt{WeightCommutatorField}, \texttt{BilinearForm}, \texttt{GradedFormMatrix},
\texttt{QuasiPrimaryDecomp}, and the sixteen files \texttt{InvariantForm*}.

These files are the ``structure theory'' of a vertex operator algebra in the sense in which one
speaks of the structure theory of Lie algebras: the operations that take one algebra and produce
another (subalgebras, quotients, tensor products, kernels of derivations, cohomology), the
symmetries (automorphisms, inner automorphisms, the Lie algebra of currents), the distinguished
vectors (singular vectors) and, above all, the \emph{invariant bilinear form}, which is the tool by
which simplicity is decided. The headline results are:
\begin{enumerate}
\item the weight-one space $\lfourV_1$ of a vertex operator algebra is a Lie algebra under
  $[u,v] = u_{(0)}v$ (Section~\ref{l04:sec:v1});
\item the exponential of a locally nilpotent current zero mode is an automorphism, and
  conjugation of such an exponential by an automorphism is again one (Section~\ref{l04:sec:aut});
\item every ideal is graded and two-sided, there is a maximal proper ideal $\lfourm(\lfourV)$, and
  $\lfourV/\lfourm(\lfourV)$ is a simple vertex operator algebra; a surjective homomorphism onto a
  simple vertex operator algebra has kernel exactly $\lfourm(\lfourV)$
  (Sections~\ref{l04:sec:quot}--\ref{l04:sec:hom});
\item the kernel of the zero mode of a weight-one state (a ``screening operator'') is a vertex
  operator subalgebra (Section~\ref{l04:sec:screening});
\item a singular vector whose generated ideal is proper obstructs simplicity, with two practical
  criteria for properness (Section~\ref{l04:sec:singular});
\item the tensor product of two vertex superalgebras and the BRST cohomology of a vertex
  superalgebra are again vertex superalgebras, under named hypotheses
  (Sections~\ref{l04:sec:tensor} and~\ref{l04:sec:brst});
\item \emph{the theorem of Frenkel--Huang--Lepowsky and Li}: if $\lfourvac\ne0$ and
  $L_1\lfourV_1 = 0$, there is a unique invariant bilinear form with $(\lfourvac|\lfourvac) = 1$
  (Section~\ref{l04:sec:canonical});
\item \emph{Li's radical theorem}: the radical of such a form is the maximal ideal, so $\lfourV$
  is simple if and only if the form is nondegenerate, and the graded dimensions of the simple
  quotient are the ranks of the Gram matrices (Sections~\ref{l04:sec:li} and~\ref{l04:sec:gram}).
\end{enumerate}
Throughout, ``vertex operator algebra'' means the library's \lean{VertexOperatorAlgebra}, which is the
Lepowsky--Li/FLM VOA \emph{of CFT type} ($\lfourV_0 = \C\lfourvac$, $\lfourV_n = 0$ for $n<0$, integer grading,
finite-dimensional pieces; Chapter~\ref{ch:L03}, Definition~\ref{l03:def:VOA}); the headlines above are theorems about VOAs
of CFT type only, a restriction relative to FLM and Lepowsky--Li. The superalgebra headline (item~6) is about other
structures: \lean{SuperVertexAlgebra} (which, like \lean{PreVertexAlgebra}, has no translation-covariance axiom and no
CFT-type condition) and \lean{SuperNearVOA} (which adds a conformal vector and the translation axiom, but not CFT type);
both are described where they arise.
\texttt{\#print axioms} on every headline reports only Lean's three standard axioms; \lean{VOA.pbw} (the PBW theorem of Chapter~\ref{ch:L02}; formerly a project axiom, now a theorem) is used nowhere in this chapter (Section~\ref{l04:sec:audit}).

The chapter also records several findings of the fidelity check that a referee should know before
relying on these files, collected here for convenience and explained where they arise:
\begin{itemize}
\item the normality of the inner automorphism group \lean{Inn} holds \emph{by construction} (it is
  defined as a normal closure), and the predicate \lean{IsExpInner} that defines it is weak enough
  that the identity satisfies it for every current (the warning on p.~\pageref{l04:warn:inn});
\item the hypothesis \lean{homega} of the inner automorphism \lean{inner} is redundant
  (Remark~\ref{l04:rem:homega});
\item the ``tensor product of VOAs'' file constructs a tensor product of vertex
  \emph{superalgebras} without conformal vectors, conditional on two linear-algebra hypotheses
  (the warning on p.~\pageref{l04:warn:tensor});
\item the library's notion of a subalgebra generated by a set always contains $\omega$
  (Definition~\ref{l04:def:generated});
\item Li's general theorem (invariant forms $\leftrightarrow$ $(\lfourV_0/L_1\lfourV_1)^*$) is
  formalised only in the case $L_1\lfourV_1 = 0$, and the symmetry of invariant forms is
  \emph{not} proved (Section~\ref{l04:sec:canonical});
\item four general routes to invariance, recorded in the files, are dead ends; their
  ``refutation'' in the library is conditional on a witness that is not constructed, though one is
  easy to give on paper (Section~\ref{l04:sec:deadends}).
\end{itemize}

\paragraph{Prerequisites.} The definitions of formal distribution, mode, the library's structure
\lean{PreVertexAlgebra} (graded VA1--VA5 data, which lacks translation covariance and is \emph{not} a vertex algebra in
the standard sense; the library's translation-covariant \lean{VertexAlgebra} adds that axiom, Chapter~\ref{ch:L03}), vertex
operator algebra (\lean{VertexOperatorAlgebra}, of CFT type by definition, with its fields \lean{grading}, \lean{Y}, \lean{vacuum}, \lean{conformalVector},
\lean{centralCharge}, \lean{virasoro}, \lean{L0_grading}, \lean{translation}, \lean{vacuum_grade},
\lean{positivity}), the Borcherds identity and skew-symmetry are those of
Chapters~\ref{ch:L01} and~\ref{ch:L03}. Several notions of this chapter (automorphisms, subalgebras, ideals,
quotients, invariant forms) were already used in the Moonshine chapters (for instance the fixed-point
sub-VOA of Chapter~\ref{ch:04}); here they are developed in full. Chapter~\ref{ch:10} uses the invariant-form
results of this chapter to prove that $\Vnat$ is simple.

%======================================================================
\section{Conventions recalled}\label{l04:sec:conv}
%======================================================================

We recall the conventions that every statement of this chapter uses; the reader who has read
Chapter~\ref{ch:L03} may skip to Section~\ref{l04:sec:v1}.

\paragraph{Modes.} A formal distribution $f = \sum_{k\in\Z} f_k z^k$ with coefficients in a space
is stored in Lean as the function $k\mapsto f_k$, i.e.\ \lean{f k} is the coefficient of $z^k$.
The \emph{mode} is
\begin{leancode}
def mode (f : FormalDist 𝒲) (m : ℤ) : 𝒲 := f (-m - 1)
\end{leancode}
File \leanpath{VOA/\allowbreak Foundations/\allowbreak FormalCalculus.lean}. So for the state--field
correspondence $Y(u,z) = \sum_n u_{(n)}z^{-n-1}$ one has \lean{mode (W.Y u) n} $= u_{(n)}$, and the
library abbreviates
\begin{leancode}
def smul (V : PreVertexAlgebra 𝒱) (u : 𝒱) (n : ℤ) (v : 𝒱) : 𝒱 := mode (V.Y u) n v
\end{leancode}
File \leanpath{VOA/\allowbreak Core/\allowbreak VertexAlgebra.lean}. Thus \lean{W.smul u n v} is the
$n$-th product $u_{(n)}v$. The coefficient form \lean{(W.Y v) k u} is $v_{(-k-1)}u$
(\lean{Y_apply_eq_smul} in \leanpath{VOA/\allowbreak Core/\allowbreak InvariantFormSymmetry.lean}).

\paragraph{Virasoro modes.} The Virasoro operators are
\begin{leancode}
def L (W : VertexOperatorAlgebra 𝒱) (n : ℤ) : End ℂ 𝒱 :=
  mode (W.Y W.conformalVector) (n + 1)
\end{leancode}
File \leanpath{VOA/\allowbreak Core/\allowbreak VertexOperatorAlgebra.lean}: $L_n = \omega_{(n+1)}$, so
$L_{-1} = \omega_{(0)}$ is the translation operator $T$, $L_0 = \omega_{(1)}$ is the grading
operator and $L_1 = \omega_{(2)}$ lowers the weight by one.

\paragraph{Weights.} The weight grading is \lean{W.grading : ℤ → Submodule ℂ 𝒱}, an internal
direct sum (\lean{isInternal}) of finite-dimensional pieces (\lean{finrank}); VA1 is the rule
$u\in\lfourV_k$, $v\in\lfourV_\ell\Rightarrow u_{(n)}v\in\lfourV_{k+\ell-n-1}$ (Lean:
\lean{smul_}\allowbreak\lean{mem_}\allowbreak\lean{grading}). The vertex operator algebras of the library are \emph{of CFT type by
definition}: \lean{vacuum_grade} says $\lfourV_0 = \C\lfourvac$ and \lean{positivity} says
$\lfourV_n = 0$ for $n<0$. This is used constantly below, and it is the reason several
hypotheses of the literature (``$\lfourV$ of CFT type'') never appear explicitly. It is also a genuine
\emph{restriction} relative to FLM and Lepowsky--Li, who ask only $\lfourV_n = 0$ for $n\ll0$: VOAs with
$\dim\lfourV_0>1$ or with negative weights are not instances, so every theorem of this chapter is a theorem about
VOAs of CFT type (release audit 2026-10-27, finding m-1). Note that CFT
type as stated allows $\lfourvac = 0$ (then $\lfourV = 0$): the zero space is a VOA in the library's sense (release
audit, finding m-4); this is why the hypothesis $\lfourvac\ne0$ appears in several theorems.

\paragraph{Projections.} \lean{piHom W n : 𝒱 →ₗ[ℂ] 𝒱} is the projection onto $\lfourV_n$ along
the other weight spaces (\leanpath{VOA/\allowbreak Modules/\allowbreak FiniteGenerationProof.lean}).

\paragraph{Lean notation used throughout.} \lean{Submodule ℂ 𝒱} is the type of complex subspaces
of $\lfourV$; \lean{⊥} and \lean{⊤} are $0$ and $\lfourV$; \lean{⊓}, \lean{⨆}, \lean{sSup},
\lean{sInf} are intersection, sum, sum of a set and intersection of a set of subspaces;
\lean{𝒱 ⧸ I} is the quotient space and \lean{I.mkQ} the quotient map; \lean{LinearMap.ker},
\lean{LinearMap.range} are kernel and image; \lean{End ℂ 𝒱} is the algebra of linear
endomorphisms; \lean{𝒱 ≃ₗ[ℂ] 𝒱} the group of linear automorphisms; \lean{↥S} (often implicit)
turns a subspace into a type; \lean{Module.finrank ℂ X} is the dimension of a finite-dimensional
space (and $0$ for an infinite-dimensional one); \lean{∑ᶠ} is a sum over a finite support (it
returns $0$ if the support is infinite, a convention whose harmlessness has to be checked each
time it is used).

%======================================================================
\section{The Lie algebra of currents}\label{l04:sec:v1}
%======================================================================

\subsection{Zero modes and the derivation property}

For a homogeneous state $u\in\lfourV_k$ the \emph{zero mode} is $o(u) = u_{(k-1)}$: by VA1 it
maps $\lfourV_\ell$ into $\lfourV_{k+\ell-(k-1)-1} = \lfourV_\ell$, i.e.\ it preserves every weight
space. For a weight-one state the zero mode is $u_{(0)}$.

\begin{definition}[Zero mode; Gannon §5.2.1]\label{l04:def:zeromode}
For $u\in\lfourV$ and $k\in\Z$, $o_k(u) := u_{(k-1)}\in\End(\lfourV)$.
\end{definition}
\noindent File \leanpath{VOA/\allowbreak Core/\allowbreak LieAlgebraV1.lean}.
\begin{leancode}
def zeroMode (W : VertexOperatorAlgebra 𝒱) (u : 𝒱) (k : ℤ) : End ℂ 𝒱 :=
  mode (W.Y u) (k - 1)

theorem zeroMode_mem_grading (W : VertexOperatorAlgebra 𝒱) {k : ℤ} {u : 𝒱}
    (hu : u ∈ W.grading k) {ℓ : ℤ} {v : 𝒱} (hv : v ∈ W.grading ℓ) :
    W.zeroMode u k v ∈ W.grading ℓ := by
\end{leancode}
\begin{dictionary}
\lean{zeroMode u k} is $u_{(k-1)}$ for \emph{any} $u$, with the weight $k$ supplied by hand,
because a general $u$ is a sum of homogeneous pieces of different weights and ``the'' zero mode
would then not be defined. Only for $u\in\lfourV_k$ (hypothesis \lean{hu} of
\lean{zeroMode_}\allowbreak\lean{mem_}\allowbreak\lean{grading}) is it the zero mode of the literature. Throughout this chapter only
$k = 1$ matters, where \lean{W.zeroMode u 1} $= u_{(0)}$ (indeed $1-1 = 0$).
\end{dictionary}
\begin{physics}
In conformal field theory a weight-$k$ field $u(z) = \sum_m u_m z^{-m-k}$ has modes $u_m =
u_{(m+k-1)}$; the zero mode $u_0 = u_{(k-1)}$ commutes with $L_0$ and so acts within each energy
level. For $k = 1$ the field is a \emph{current} and $u_0 = \oint u(z)\,dz/2\pi i$ is its
\emph{charge}: a conserved quantity that acts on the space of states without changing energies.
\end{physics}
\begin{status}
Definition and an unconditional lemma (VA1); axiom-clean.
\end{status}

\begin{theorem}[Zero modes are derivations; Gannon (5.2.4)]\label{l04:thm:deriv}
For all $u,v,w\in\lfourV$ and $n\in\Z$:
$u_{(0)}(v_{(n)}w) - v_{(n)}(u_{(0)}w) = (u_{(0)}v)_{(n)}w$.
\end{theorem}
\noindent File \leanpath{VOA/\allowbreak Core/\allowbreak LieAlgebraV1.lean}.
\begin{leancode}
theorem zeroMode_isDerivation (W : VertexOperatorAlgebra 𝒱) (u v w : 𝒱) (n : ℤ) :
    W.smul u 0 (W.smul v n w) - W.smul v n (W.smul u 0 w) = W.smul (W.smul u 0 v) n w := by
\end{leancode}
File \leanpath{VOA/\allowbreak Core/\allowbreak Automorphism.lean} restates it as a Leibniz rule:
\begin{leancode}
theorem zeroMode_leibniz (W : VertexOperatorAlgebra 𝒱) (u a b : 𝒱) (n : ℤ) :
    W.smul u 0 (W.smul a n b) = W.smul (W.smul u 0 a) n b + W.smul a n (W.smul u 0 b) := by
\end{leancode}
\begin{dictionary}
Both statements are for \emph{arbitrary} $u$, $v$, $w$, $n$: no weight hypothesis at all. This is
correct: the commutator formula $[u_{(m)},v_{(n)}] = \sum_{j\ge0}\binom mj(u_{(j)}v)_{(m+n-j)}$
at $m = 0$ has the single term $j = 0$. The Lean statement is exactly the textbook one with the
commutator written out. It could not be vacuous: both sides are vectors of $\lfourV$ and there are
no hypotheses.
\end{dictionary}
\begin{physics}
The charge $u_0$ of any field acts on every operator product as a derivation: $[Q,\phi_1(z)\phi_2(w)]
= [Q,\phi_1(z)]\phi_2(w) + \phi_1(z)[Q,\phi_2(w)]$, and $[u_{(0)},Y(v,z)] = Y(u_{(0)}v,z)$ says that
the charge acts on a field by acting on the state that creates it. This is the infinitesimal form of
the statement that conserved charges generate symmetries of the operator algebra.
\end{physics}
\emph{Proof on paper.} The $n$-th product formula $Y(u_{(n)}v,z) = \Res_{z_1}
\big((z_1-z)^n[Y(u,z_1),Y(v,z)]\big)$ (valid in any vertex algebra in the standard sense, and proved in the library for
\lean{VertexOperatorAlgebra} but not for the weaker \lean{PreVertexAlgebra}; Chapter~\ref{ch:L03}) at
$n = 0$ reads $Y(u_{(0)}v,z) = \Res_{z_1}[Y(u,z_1),Y(v,z)] = [u_{(0)},Y(v,z)]$. Taking the coefficient
of $z^{-n-1}$ and applying to $w$ gives the claim. \qed
\begin{status}
Unconditional; \texttt{\#print axioms} = \lean{[propext, Classical.choice, Quot.sound]}.
\end{status}

\subsection{The bracket on the weight-one space}

\begin{definition}[The Lie bracket on $\lfourV_1$]\label{l04:def:bracket}
For $u,v\in\lfourV_1$, $[u,v] := u_{(0)}v$, which lies in $\lfourV_{1+1-0-1} = \lfourV_1$.
\end{definition}
\noindent File \leanpath{VOA/\allowbreak Core/\allowbreak LieAlgebraV1.lean}.
\begin{leancode}
def bracketV1 (W : VertexOperatorAlgebra 𝒱) (u v : W.grading 1) : W.grading 1 :=
  ⟨W.smul (u : 𝒱) 0 (v : 𝒱), by
    have h := W.smul_mem_grading u.2 v.2 0
    rwa [show (1 : ℤ) + 1 - 0 - 1 = 1 by ring] at h⟩
\end{leancode}
\begin{dictionary}
The type \lean{W.grading 1} is the subspace $\lfourV_1$ viewed as a type; an element is a pair of
a vector and a proof of membership, and \lean{(u : 𝒱)} forgets the proof. The anonymous
constructor \lean{⟨x, proof⟩} packages $u_{(0)}v$ with the VA1 proof that it is again in
$\lfourV_1$ (the one-line proof is shown because it is informative: it is exactly the weight count
$1+1-0-1 = 1$). The bracket is literally $u_{(0)}v$, with no sign or normalisation.
\end{dictionary}
\begin{physics}
$\lfourV_1$ is the space of conserved currents $J^a(z)$. Their OPE has the form
$J^a(z)J^b(w)\sim \frac{k^{ab}}{(z-w)^2} + \frac{f^{ab}{}_c J^c(w)}{z-w}$; the simple-pole coefficient is
$J^a_{(0)}J^b = f^{ab}{}_cJ^c$, the structure constants of the global symmetry algebra, and the double
pole $J^a_{(1)}J^b = k^{ab}\lfourvac$ is the level (the current--current two-point function).
\end{physics}
\begin{status}
Definition; axiom-clean.
\end{status}

\begin{theorem}[$\lfourV_1$ is a Lie algebra; Gannon §5.2.1]\label{l04:thm:lie}
The bracket of Definition~\ref{l04:def:bracket} is bilinear, antisymmetric and satisfies the
Jacobi identity $[[u,v],w] = [u,[v,w]] - [v,[u,w]]$. Hence $(\lfourV_1,[\cdot,\cdot])$ is a complex
Lie algebra.
\end{theorem}
\noindent File \leanpath{VOA/\allowbreak Core/\allowbreak LieAlgebraV1.lean}.
\begin{leancode}
theorem bracketV1_jacobi (W : VertexOperatorAlgebra 𝒱) (u v w : W.grading 1) :
    W.bracketV1 (W.bracketV1 u v) w
      = W.bracketV1 u (W.bracketV1 v w) - W.bracketV1 v (W.bracketV1 u w) := by

theorem bracketV1_antisymm (W : VertexOperatorAlgebra 𝒱) (u v : W.grading 1) :
    W.bracketV1 u v = - W.bracketV1 v u :=

noncomputable instance instLieRingGradingOne (W : VertexOperatorAlgebra 𝒱) :
    LieRing (W.grading 1) where
  bracket u v := W.bracketV1 u v

noncomputable instance instLieAlgebraGradingOne (W : VertexOperatorAlgebra 𝒱) :
    LieAlgebra ℂ (W.grading 1) where
  lie_smul c u v := W.bracketV1_smul_right c u v
\end{leancode}
\begin{dictionary}
An \lean{instance} of \lean{LieRing} and \lean{LieAlgebra ℂ} on the type $\lfourV_1$ is Mathlib's way
of saying ``$\lfourV_1$ is a complex Lie algebra''; the fields not shown (\lean{add_lie},
\lean{lie_add}, \lean{lie_self}, \lean{leibniz_lie}) are discharged by the bilinearity lemmas,
antisymmetry and Jacobi. The line \lean{bracket u v := W.bracketV1 u v} pins the Lie bracket
\lean{⁅u, v⁆} of Mathlib to $u_{(0)}v$, so there is no ambiguity about which bracket is meant.
Mathlib's Leibniz form $[u,[v,w]] = [[u,v],w] + [v,[u,w]]$ is equivalent to the displayed Jacobi
identity. The statements are unconditional and exactly the textbook ones.

A historical remark the source records: antisymmetry was at first stated only conditionally
(\lean{bracketV1_}\allowbreak\lean{antisymm_}\allowbreak\lean{of_}\allowbreak\lean{skew}, taking $u_{(0)}v+v_{(0)}u = 0$ as hypothesis) while
skew-symmetry was unproved; now skew-symmetry is a theorem (Chapter~\ref{ch:L03}, Theorem~\ref{l03:thm:skew}) and
\lean{bracketV1_}\allowbreak\lean{antisymm} discharges the hypothesis via \lean{smul_}\allowbreak\lean{zero_}\allowbreak\lean{add_}\allowbreak\lean{swap_}\allowbreak\lean{eq_}\allowbreak\lean{zero}.
\end{dictionary}
\begin{physics}
The currents of a chiral algebra close into a Lie algebra $\mathfrak g = \lfourV_1$: this is the
\emph{continuous global symmetry} of the chiral CFT, and the modes $J^a_n$ form the affine
(Kac--Moody) algebra $\hat{\mathfrak g}$ whose zero modes $J^a_0$ generate $\mathfrak g$ acting on each
energy level. For the Moonshine module $\Vnat_1 = 0$: the theory has no continuous symmetry, only the
finite Monster group.
\end{physics}
\emph{Proof on paper.} Jacobi is Theorem~\ref{l04:thm:deriv} at $n = 0$. Antisymmetry:
skew-symmetry $Y(u,z)v = e^{zT}Y(v,-z)u$ gives, at the coefficient of $z^{-1}$,
$u_{(0)}v = -v_{(0)}u + T(v_{(1)}u) - \tfrac12T^2(v_{(2)}u)+\cdots$. For $u,v\in\lfourV_1$, $v_{(1)}u\in
\lfourV_0 = \C\lfourvac$ and $T\lfourvac = 0$; $v_{(j)}u\in\lfourV_{1-j} = 0$ for $j\ge2$. So
$u_{(0)}v = -v_{(0)}u$. \qed
\begin{status}
Unconditional; \texttt{\#print axioms} on \lean{instLieAlgebraGradingOne} and
\lean{bracketV1_}\allowbreak\lean{antisymm}: \lean{[propext, Classical.choice, Quot.sound]}.
\end{status}

\begin{example}[Affine and lattice examples]\label{l04:ex:affine}
For the universal affine vertex algebra $V^k(\mathfrak g)$ (Chapters~\ref{ch:L08} and~\ref{ch:L23}),
$\lfourV_1 = \mathfrak g$ and $x_{(0)}y = [x,y]$, so Theorem~\ref{l04:thm:lie} returns $\mathfrak g$.
For a lattice vertex algebra $V_L$ with $L$ even and positive definite, $\lfourV_1 = \mathfrak h\oplus
\bigoplus_{\langle\alpha,\alpha\rangle = 2}\C e^\alpha$, which is the simply-laced Lie algebra with root
system $\{\alpha : \langle\alpha,\alpha\rangle = 2\}$ (Frenkel--Kac); for the Leech lattice there
are no roots and $\lfourV_1 = \mathfrak h$ is abelian of dimension $24$.
\end{example}

\subsection{Primary vectors and the Griess product}

\begin{definition}[Primary subspace; Gannon (5.2.3)]\label{l04:def:primary}
$\mathcal P\lfourV_n := \{v\in\lfourV_n : L_kv = 0\text{ for all }k\ge1\}$.
\end{definition}
\noindent File \leanpath{VOA/\allowbreak Core/\allowbreak LieAlgebraV1.lean}.
\begin{leancode}
def primaryV (W : VertexOperatorAlgebra 𝒱) (n : ℤ) : Submodule ℂ 𝒱 where
  carrier := {v | v ∈ W.grading n ∧ ∀ k : ℤ, 1 ≤ k → W.L k v = 0}
\end{leancode}
\begin{dictionary}
Exactly the textbook definition (the omitted fields are the proofs that it is a subspace).
``Primary'' (killed by all $L_k$, $k\ge1$) is stronger than ``quasi-primary'' (killed by $L_1$),
a notion used in Section~\ref{l04:sec:canonical}.
\end{dictionary}
\begin{physics}
Primary states create primary fields: fields that transform under all conformal maps as tensors of
weight $n$, $\phi(z)\mapsto (f'(z))^n\phi(f(z))$. They are the highest-weight vectors of the Virasoro
algebra inside the chiral algebra.
\end{physics}
\begin{status}
Definition; axiom-clean.
\end{status}

\begin{definition}[Griess product]\label{l04:def:griess}
For $u,v\in\lfourV_2$, $u\times v := u_{(1)}v\in\lfourV_{2+2-1-1} = \lfourV_2$.
\end{definition}
\noindent File \leanpath{VOA/\allowbreak Core/\allowbreak LieAlgebraV1.lean}.
\begin{leancode}
def griessMul (W : VertexOperatorAlgebra 𝒱) (u v : W.grading 2) : W.grading 2 :=
  ⟨W.smul (u : 𝒱) 1 (v : 𝒱), by
    have h := W.smul_mem_grading u.2 v.2 1
    rwa [show (2 : ℤ) + 2 - 1 - 1 = 2 by ring] at h⟩

theorem griessMul_comm_of_grading1_eq_bot (W : VertexOperatorAlgebra 𝒱)
    (h1 : W.grading 1 = ⊥) (u v : W.grading 2) :
    W.griessMul u v = W.griessMul v u := by
\end{leancode}
\begin{dictionary}
\lean{h1 : W.grading 1 = ⊥} is $\lfourV_1 = 0$. The theorem says the Griess product is
commutative when there are no currents. This is the correct hypothesis: in general skew-symmetry
gives $u_{(1)}v = v_{(1)}u - T(v_{(2)}u) + \tfrac12T^2(v_{(3)}u)-\cdots$ with $v_{(2)}u\in\lfourV_1$ and
$v_{(3)}u\in\lfourV_0$, and only the $\lfourV_1$ term can survive, so without $\lfourV_1 = 0$ the
product is only commutative modulo $T\lfourV_1$. Associativity is not claimed (the Griess algebra is
not associative).
\end{dictionary}
\begin{physics}
On $\Vnat$, $\lfourV_2$ has dimension $196884$ and the Griess product is the commutative non-associative
algebra whose automorphism group is the Monster; it is the OPE coefficient $T^a(z)T^b(w)\supset
\frac{(a\times b)(w)}{(z-w)^2}$ among the weight-two fields. Commutativity reflects the bosonic
symmetry of the OPE of two weight-two fields once no weight-one field can appear.
\end{physics}
\emph{Proof on paper.} As just explained in the fidelity check, with $\lfourV_1 = 0$ killing the
only correction term. \qed
\begin{status}
Unconditional given \lean{h1}; axiom-clean.
\end{status}

\begin{remark}[Three re-export files]\label{l04:rem:stubs}
\leanpath{VOA/\allowbreak Core/\allowbreak V1LieAlgebra.lean}, \leanpath{VOA/\allowbreak Core/\allowbreak InnerAutomorphism.lean} and
\leanpath{VOA/\allowbreak Core/\allowbreak CommutativeRing.lean} contain no declarations: each is a single
\texttt{import} with a docstring explaining that its content was absorbed into
\texttt{LieAlgebraV1}, \texttt{Automorphism} and \texttt{Commutativity} respectively. The last of these
(the commutative ring structure on a commutative vertex algebra, Gannon Thm.~5.1.4) is documented in
Chapter~\ref{ch:L03}. Nothing in these three files needs refereeing.
\end{remark}

%======================================================================
\section{Automorphisms and inner automorphisms}\label{l04:sec:aut}
%======================================================================

\subsection{The automorphism group}

\begin{definition}[Automorphism; Gannon §5.2.1]\label{l04:def:aut}
An automorphism of $\lfourV$ is a linear bijection $\alpha:\lfourV\to\lfourV$ with
$\alpha(u_{(n)}v) = (\alpha u)_{(n)}(\alpha v)$ for all $u,v,n$, $\alpha\lfourvac = \lfourvac$ and
$\alpha\omega = \omega$. Composition makes the automorphisms a group $\Aut(\lfourV)$.
\end{definition}
\noindent File \leanpath{VOA/\allowbreak Core/\allowbreak Automorphism.lean}.
\begin{leancode}
structure Aut (W : VertexOperatorAlgebra 𝒱) where
  toLinearEquiv : 𝒱 ≃ₗ[ℂ] 𝒱
  map_smul' : ∀ (u : 𝒱) (n : ℤ) (v : 𝒱),
    toLinearEquiv (W.smul u n v) = W.smul (toLinearEquiv u) n (toLinearEquiv v)
  map_vacuum : toLinearEquiv W.vacuum = W.vacuum
  map_omega : toLinearEquiv W.conformalVector = W.conformalVector

instance : Group (Aut W) where
  mul := comp
  one := id W
  inv := inverse

@[simp] theorem mul_apply (α β : Aut W) (v : 𝒱) : (α * β) v = α (β v) := rfl
\end{leancode}
\begin{dictionary}
\lean{𝒱 ≃ₗ[ℂ] 𝒱} is a linear bijection with its inverse. The group law is composition with
$(\alpha\beta)v = \alpha(\beta v)$ (the \lean{rfl} lemma \lean{mul_apply}). Grading preservation is
not an axiom but a theorem (\lean{Aut.map_grading}, Theorem~\ref{l04:thm:autgrading}), because $\alpha$ commutes with $L_0 = \omega_{(1)}$. On paper the
condition $\alpha\lfourvac = \lfourvac$ is even redundant: by CFT type and grading preservation
$\alpha\lfourvac = \lambda\lfourvac$ with $\lambda\ne0$, and applying $\alpha$ to $\lfourvac_{(-1)}\lfourvac = \lfourvac$ gives
$\lambda^2 = \lambda$. The definition
agrees with FLM and Lepowsky--Li. The file also proves \lean{aut_comm_L}: $\alpha L_n = L_n\alpha$.
\end{dictionary}
\begin{physics}
A global (internal) symmetry of the chiral CFT: an invertible map of the space of states preserving
the OPE, the vacuum and the stress tensor. Examples: the $\Z_2$ symmetry $\phi\mapsto-\phi$ of a free
boson, the Weyl group and torus of an affine algebra, the Monster acting on $\Vnat$. Fixing $\omega$
means the symmetry commutes with conformal transformations.
\end{physics}
\begin{status}
Definition and group structure; axiom-clean.
\end{status}

\begin{theorem}[Automorphisms preserve the grading]\label{l04:thm:autgrading}
Let $\alpha\in\Aut(\lfourV)$. Then $\alpha(\lfourV_n)\subseteq\lfourV_n$ for every $n\in\Z$.
\end{theorem}
\noindent File \leanpath{VOA/\allowbreak Modularity/\allowbreak Twisted.lean}.
\begin{leancode}
theorem Aut.map_grading (α : Aut W) {n : ℤ} {v : 𝒱} (hv : v ∈ W.grading n) :
    α v ∈ W.grading n
\end{leancode}

\begin{dictionary}
\lean{Aut W} is the structure of Definition~\ref{l04:def:aut}; \lean{α v} is the application of its underlying linear
equivalence (a coercion to a function); \lean{W.grading n} is the weight space $\lfourV_n$ of the VOA structure
(Chapter~\ref{ch:L03}), a \lean{Submodule ℂ 𝒱}. \textbf{Hypotheses:} the vector $v$ lies in $\lfourV_n$; the integer $n$
is arbitrary. \textbf{Why this is a theorem and not an axiom.} The definition of \lean{Aut} requires
$\alpha(u_{(n)}v) = (\alpha u)_{(n)}\alpha v$, $\alpha\lfourvac = \lfourvac$, $\alpha\omega = \omega$, but \emph{not} that
$\alpha$ respects the grading. FLM and Lepowsky--Li define automorphisms as grading-preserving maps with
$\alpha Y(u,z)\alpha^{-1} = Y(\alpha u,z)$ and $\alpha\omega=\omega$; the statement proved here shows that the library's
(shorter) definition is equivalent to theirs, so nothing is lost. The same-or-stronger-or-weaker check: the
conclusion is equivalent to $\alpha(\lfourV_n)\subseteq\lfourV_n$ for all $n$, hence, applying it to $\alpha^{-1}$, to
equality $\alpha(\lfourV_n) = \lfourV_n$. \textbf{Could the definition have been vacuous?} The identity is an automorphism, so
$\Aut(\lfourV)$ is never empty and the theorem is never about an empty class; and $\lfourV_n$ is the genuine
$n$-eigenspace of $L_0=\omega_{(1)}$, not a convention that might be changed by $\alpha$.
\end{dictionary}

\begin{physics}
An automorphism is an internal symmetry of the chiral algebra, and it fixes the stress tensor. A symmetry that
commutes with $L_0$ cannot mix operators of different scaling dimension: each space of operators of fixed conformal
weight, $\lfourV_n$, is a representation of the symmetry group. For the $\Z_2$ symmetry $\phi\mapsto-\phi$ of a free
boson it acts on the weight-$h$ space by $(-1)^{\#\text{oscillators}}$; for the Monster it is the statement that
$\Vnat_n$ is a finite-dimensional Monster representation for each $n$.
\end{physics}

\textbf{Proof on paper.} Since $\alpha$ is an automorphism, $\alpha(\omega_{(1)}v) = (\alpha\omega)_{(1)}\alpha v =
\omega_{(1)}\alpha v$, i.e.\ $\alpha$ commutes with $L_0=\omega_{(1)}$. Now $\lfourV_n$ is exactly the $n$-eigenspace of
$L_0$: if $x=\sum_kx_k$ is the decomposition of a vector into homogeneous components and $L_0x = nx$, then
$\sum_k(k-n)x_k = 0$, and since the sum $\bigoplus_k\lfourV_k$ is direct, $(k-n)x_k=0$ for each $k$, so $x = x_n$.
Hence for $v\in\lfourV_n$ we have $L_0(\alpha v) = \alpha(L_0v) = n\,\alpha v$, so $\alpha v\in\lfourV_n$. \qed

\begin{status}
Unconditional. \texttt{\#print axioms} \lean{VOA.Aut.map\_grading} $=$
\lean{[propext, Classical.}\allowbreak\lean{choice, Quot.}\allowbreak\lean{sound]} (run for this chapter with
\texttt{lake env lean -DmaxSynthPendingDepth=3}). The PBW theorem \lean{VOA.pbw} is not used.
\end{status}

\subsection{Exponentials of locally nilpotent derivations}

To exponentiate a current zero mode one must make sense of $e^D = \sum_k D^k/k!$ on an
infinite-dimensional space without any topology. The library does this for \emph{locally
nilpotent} $D$: every vector is killed by some power of $D$, so the series is a finite sum on each
vector.

\begin{definition}[Local nilpotency, locally finite exponential]\label{l04:def:exp}
$D\in\End(\lfourV)$ is locally nilpotent if for every $v$ there is $N$ with $D^Nv = 0$. Then
$e^Dv := \sum_{k<N}D^kv/k!$ for any such $N$ (independent of the choice).
\end{definition}
\noindent File \leanpath{VOA/\allowbreak Core/\allowbreak Automorphism.lean}.
\begin{leancode}
noncomputable def expSum (D : End ℂ 𝒱) (N : ℕ) (v : 𝒱) : 𝒱 :=
  ∑ k ∈ Finset.range N, ((k.factorial : ℂ)⁻¹) • ((D ^ k) v)

def IsLocNilp (D : End ℂ 𝒱) : Prop := ∀ v : 𝒱, ∃ N : ℕ, (D ^ N) v = 0

noncomputable def expFun {D : End ℂ 𝒱} (hD : IsLocNilp D) (v : 𝒱) : 𝒱 :=
  expSum D (locNilpIdx hD v) v

noncomputable def expEquiv {D : End ℂ 𝒱} (hD : IsLocNilp D) : 𝒱 ≃ₗ[ℂ] 𝒱 where
  toLinearMap := expLN hD
  invFun := expLN hD.neg
\end{leancode}
\begin{dictionary}
\lean{locNilpIdx hD v} is an index $N$ chosen (by the axiom of choice) with $D^Nv = 0$;
\lean{expSum_indep} proves that any two such indices give the same partial sum, so \lean{expFun} does
not depend on the choice. \lean{expLN} packages \lean{expFun} as a linear map, and \lean{expEquiv}
is the linear bijection with inverse $e^{-D}$ (\lean{hD.neg} is local nilpotency of $-D$; the
inverse property is the convolution identity \lean{expSum_conv}, the finite form of
$e^De^{-D} = 1$). This is the standard algebraic exponential (it is the device used for integrable representations in Kac's \emph{Infinite-dimensional
Lie algebras}).
\end{dictionary}
\begin{physics}
The finite exponential of a nilpotent charge is a \emph{unipotent} symmetry, like $e^{tE}$ for a
raising operator $E$ of $\mathfrak{sl}_2$ acting in a finite-dimensional representation. Each energy
level of a CFT is finite-dimensional and the charge preserves it, which is why local nilpotency is the
natural algebraic condition.
\end{physics}
\begin{status}
Definitions; the linear equivalence is axiom-clean.
\end{status}

\begin{theorem}[The exponential of a derivation is a homomorphism]\label{l04:thm:expsmul}
If $D$ is locally nilpotent and $D(a_{(n)}b) = (Da)_{(n)}b + a_{(n)}(Db)$ for all $a,b$, then
$e^D(a_{(n)}b) = (e^Da)_{(n)}(e^Db)$.
\end{theorem}
\noindent File \leanpath{VOA/\allowbreak Core/\allowbreak Automorphism.lean}.
\begin{leancode}
theorem expFun_smul {W : VertexOperatorAlgebra 𝒱} {D : End ℂ 𝒱} (hD : IsLocNilp D) (n : ℤ)
    (hleib : ∀ a b : 𝒱, D (W.smul a n b) = W.smul (D a) n b + W.smul a n (D b))
    (a b : 𝒱) :
    expFun hD (W.smul a n b) = W.smul (expFun hD a) n (expFun hD b) := by
\end{leancode}
\begin{dictionary}
The Leibniz hypothesis \lean{hleib} is for the fixed $n$ only, and the conclusion is for that $n$:
the statement is per product, which is slightly sharper than needed.
\end{dictionary}
\begin{physics}
A charge that acts as a derivation on all OPEs exponentiates to a finite symmetry that preserves all
OPEs: the passage from the Lie algebra of conserved charges to the symmetry group.
\end{physics}
\emph{Proof on paper.} By induction, $D^k(a_{(n)}b) = \sum_i\binom ki(D^ia)_{(n)}(D^{k-i}b)$
(\lean{iter_leibniz}); dividing by $k!$ and summing over $k$ up to $2N$, where $D^Na = D^Nb = 0$,
the Cauchy product gives $\sum_{i,j<N}\frac{1}{i!j!}(D^ia)_{(n)}(D^jb) = (e^Da)_{(n)}(e^Db)$. \qed
\begin{status}
Unconditional given its hypotheses; axiom-clean.
\end{status}

\subsection{Inner automorphisms}

\begin{definition}[Inner automorphism; Gannon (5.2.4)]\label{l04:def:inner}
For $u\in\lfourV_1$ with $u_{(0)}$ locally nilpotent and $u_{(0)}\omega = 0$, the automorphism
$\exp(u_{(0)})$.
\end{definition}
\noindent File \leanpath{VOA/\allowbreak Core/\allowbreak Automorphism.lean}.
\begin{leancode}
def IsExpInner (u : W.grading 1) (α : Aut W) : Prop :=
  ∀ v : 𝒱, ∃ N : ℕ, α v = ∑ k ∈ Finset.range N,
    ((k.factorial : ℂ)⁻¹) • ((W.zeroMode (u : 𝒱) 1 ^ k) v)

noncomputable def inner (u : W.grading 1)
    (hnilp : VOA.IsLocNilp (W.zeroMode (u : 𝒱) 1))
    (homega : W.zeroMode (u : 𝒱) 1 W.conformalVector = 0) : Aut W where
  toLinearEquiv := VOA.expEquiv hnilp

theorem inner_isExpInner (u : W.grading 1)
    (hnilp : VOA.IsLocNilp (W.zeroMode (u : 𝒱) 1))
    (homega : W.zeroMode (u : 𝒱) 1 W.conformalVector = 0) :
    IsExpInner u (inner u hnilp homega) := by

def innerSet (W : VertexOperatorAlgebra 𝒱) : Set (Aut W) :=
  {α | ∃ u : W.grading 1, IsExpInner u α}

def Inn (W : VertexOperatorAlgebra 𝒱) : Subgroup (Aut W) :=
  Subgroup.normalClosure (innerSet W)

instance Inn_normal (W : VertexOperatorAlgebra 𝒱) : (Inn W).Normal :=
  Subgroup.normalClosure_normal
\end{leancode}
\begin{dictionary}
\lean{inner u hnilp homega} is the automorphism whose underlying linear map is \lean{expEquiv} of
$D = u_{(0)}$; its omitted fields prove that it fixes $\lfourvac$ (since $u_{(0)}\lfourvac = 0$ by
creation), fixes $\omega$ (by \lean{homega}), and intertwines all products
(Theorem~\ref{l04:thm:expsmul} with Leibniz from Theorem~\ref{l04:thm:deriv}). This construction
is faithful to the literature.

\lean{IsExpInner u α} is an interface predicate: \emph{for each $v$ there is some $N$} such that
$\alpha v$ equals the $N$-th partial sum of the exponential series of $u_{(0)}$ on $v$. When
$u_{(0)}$ is locally nilpotent and $N$ is large, the partial sums stabilise at $e^{u_{(0)}}v$; but
the predicate does not require $N$ to be large. See the warning on p.~\pageref{l04:warn:inn}.
\lean{Subgroup.normalClosure S} is the smallest normal subgroup containing the set $S$.
\end{dictionary}
\begin{physics}
Inner automorphisms are the symmetries generated by the currents themselves: $e^{\theta J_0}$ for a
charge $J_0$. Physically, the continuous symmetry group $G$ whose Lie algebra is $\lfourV_1$; its
subgroup generated by such exponentials is normal because conjugating the symmetry generated by a
current $J$ by any symmetry $g$ gives the symmetry generated by the current $gJ$.
\end{physics}
\begin{status}
Definitions; \texttt{\#print axioms} on \lean{inner} and \lean{Inn_normal}:
\lean{[propext, Classical.choice, Quot.sound]}. The construction \lean{inner} is used downstream in
\leanpath{VOA/\allowbreak Moonshine/\allowbreak TrialityRoot.lean}.
\end{status}

\begin{theorem}[Conjugation of inner automorphisms; Gannon (5.2.4)]\label{l04:thm:conj}
If $\beta$ is the exponential of $u_{(0)}$ and $\alpha\in\Aut(\lfourV)$, then $\alpha\beta\alpha^{-1}$
is the exponential of $(\alpha u)_{(0)}$. Moreover $u_{(0)}$ locally nilpotent implies
$(\alpha u)_{(0)}$ locally nilpotent, and $\alpha(e^{u_{(0)}}v) = e^{(\alpha u)_{(0)}}(\alpha v)$.
\end{theorem}
\noindent Files: \lean{isExpInner_conj} is in \leanpath{VOA/Modularity/Twisted.lean};
\lean{aut_conj_expFun} is in \leanpath{VOA/\allowbreak Core/\allowbreak Automorphism.lean}.
\begin{leancode}
theorem VertexOperatorAlgebra.isExpInner_conj {u : W.grading 1} {β : Aut W}
    (hβ : VertexOperatorAlgebra.IsExpInner u β) (α : Aut W) :
    VertexOperatorAlgebra.IsExpInner (⟨α (u : 𝒱), α.map_grading u.2⟩ : W.grading 1)
      (α * β * α⁻¹) := by

theorem aut_conj_expFun (α : Aut W) (u : 𝒱)
    (hnilp : VOA.IsLocNilp (W.zeroMode u 1)) (v : 𝒱) :
    α (VOA.expFun hnilp v) = VOA.expFun (isLocNilp_aut α u hnilp) (α v) := by
\end{leancode}
\begin{dictionary}
The conjugation statement now lives in \leanpath{VOA/\allowbreak Modularity/\allowbreak Twisted.lean} under the full name
\lean{VertexOperatorAlgebra.isExpInner_conj}. An earlier version (in \leanpath{Automorphism.lean}) carried a hypothesis
\lean{hαu : α u ∈ W.grading 1}, needed only to form the element of \lean{W.grading 1}; this book flagged it as redundant
(it always holds by \lean{Aut.map_grading}), and the current statement builds the element directly as
\lean{⟨α u, α.map_grading u.2⟩}, with no hypothesis. \lean{aut_conj_expFun} is the genuine conjugation formula
$\alpha e^{u_{(0)}}\alpha^{-1} = e^{(\alpha u)_{(0)}}$ for locally nilpotent zero modes;
\lean{isExpInner_conj} is the same statement at the level of the interface predicate.
\end{dictionary}
\begin{physics}
Conjugating the symmetry generated by the charge of $J$ by a global symmetry $g$ gives the symmetry
generated by the charge of $gJ$. This is the statement that the inner symmetries form a normal
subgroup: the quotient $\Aut/\mathrm{Inn}$ is the group of ``outer'' (discrete) symmetries, like the
diagram automorphisms of an affine algebra or the Monster for $\Vnat$ (where $\mathrm{Inn}$ is trivial).
\end{physics}
\emph{Proof on paper.} $\alpha(u_{(0)}v) = (\alpha u)_{(0)}(\alpha v)$, hence
$\alpha u_{(0)}^k = (\alpha u)_{(0)}^k\alpha$; apply $\alpha$ to the finite exponential sum. Local
nilpotency transfers because $(\alpha u)_{(0)}^k = \alpha u_{(0)}^k\alpha^{-1}$. \qed
\begin{status}
Unconditional; axiom-clean.
\end{status}

\begin{warning}[title={Trap / dead end: What the normality of \texttt{Inn} does and does not say}]\label{l04:warn:inn}
\textbf{(a) Normality is by construction.} \lean{Inn W} is \emph{defined} as the normal closure of
\lean{innerSet W}, so \lean{Inn_normal} is Mathlib's \lean{Subgroup.}\allowbreak\lean{normalClosure_}\allowbreak\lean{normal}: it holds
for the normal closure of any set whatsoever. The mathematical content of ``$\mathrm{Inn}\trianglelefteq
\Aut$'' --- that the subgroup \emph{generated} by the inner automorphisms is already normal --- is
\lean{isExpInner_conj}, which shows that \lean{innerSet} is stable under conjugation (so that its
normal closure equals the subgroup it generates). The source says exactly this. A referee should
read the headline ``\lean{Inn} is normal'' as resting on \lean{isExpInner_conj}, not on
\lean{Inn_normal}.

\textbf{(b) \lean{IsExpInner} is weak.} Taking $N = 1$ in its definition gives the partial sum
$v$; so \emph{the identity automorphism satisfies \lean{IsExpInner u 1} for every current $u$},
whether or not $u_{(0)}$ is locally nilpotent. We checked this in Lean for this book (scratch file,
not part of the library):
\begin{plaincode}
example (u : W.grading 1) : IsExpInner u (1 : Aut W) := fun v => ⟨1, by simp⟩
\end{plaincode}
This is harmless for the identity (which is inner anyway), but it shows that \lean{innerSet} is
defined by a property weaker than ``$\alpha = e^{u_{(0)}}$'': an automorphism $\alpha$ belongs to it as
soon as, for some $u$, each $\alpha v$ is \emph{some} partial sum of the exponential series of $v$.
Whether \lean{innerSet} can contain automorphisms that are not exponentials of locally nilpotent
current zero modes has not been examined by the library. The safe reading is: the automorphisms
\lean{inner u hnilp homega} are genuine inner automorphisms (\lean{inner_mem_Inn}) and are in
\lean{Inn W}; nothing in this chapter proves that \lean{Inn W} consists only of products of such.

\textbf{(c) Not every current exponentiates this way.} \lean{inner} needs $u_{(0)}$ locally
nilpotent. Root currents $e^\alpha$ of a lattice vertex algebra qualify (in a fixed weight space only
finitely many charges $\beta+k\alpha$ occur), but Cartan currents $h$ do not: $h_{(0)}e^\beta =
\langle h,\beta\rangle e^\beta$ is diagonal. The torus $\exp(\mathfrak h)$ acts by automorphisms, but
they are not constructed by this file.
\end{warning}

\begin{remark}[The hypothesis \texttt{homega} is redundant]\label{l04:rem:homega}
For $u\in\lfourV_1$ in a vertex operator algebra of CFT type, $u_{(0)}\omega = 0$ automatically:
skew-symmetry gives $u_{(0)}\omega = -L_{-1}u + T(L_0u) - \tfrac12T^2(L_1u) + \cdots = -Tu + Tu -
\tfrac12T^2(c\lfourvac) = 0$, since $L_1u\in\lfourV_0 = \C\lfourvac$ and $T\lfourvac = 0$, and higher
terms vanish by weight. The library proves this as \lean{Screening.}\allowbreak\lean{screening_}\allowbreak\lean{conformalVector}
(Section~\ref{l04:sec:screening}), but \lean{inner} still asks for it. We checked in Lean that it is
discharged in one line:
\begin{plaincode}
example (u : W.grading 1) : W.zeroMode (u : 𝒱) 1 W.conformalVector = 0 := by
  rw [← Screening.screening_eq_zeroMode]; exact Screening.screening_conformalVector W u.2
\end{plaincode}
So \lean{inner} is in effect conditional only on local nilpotency.
\end{remark}

%======================================================================
\section{Ideals, quotients and the simple quotient}\label{l04:sec:quot}
%======================================================================

\subsection{Ideals and simplicity}

\begin{definition}[Ideal, simple VOA]\label{l04:def:simple}
An ideal of $\lfourV$ is a subspace $I$ with $u_{(n)}v\in I$ for all $u\in\lfourV$, $v\in I$, $n\in\Z$.
$\lfourV$ is simple if $\lfourV\ne0$ and its only ideals are $0$ and $\lfourV$.
\end{definition}
\noindent File \leanpath{VOA/\allowbreak Core/\allowbreak BilinearForm.lean}.
\begin{leancode}
def IsIdeal (W : VertexOperatorAlgebra 𝒱) (I : Submodule ℂ 𝒱) : Prop :=
  ∀ (u v : 𝒱) (n : ℤ), v ∈ I → W.smul u n v ∈ I

def IsSimple (W : VertexOperatorAlgebra 𝒱) : Prop :=
  (⊤ : Submodule ℂ 𝒱) ≠ ⊥ ∧ ∀ I : Submodule ℂ 𝒱, W.IsIdeal I → I = ⊥ ∨ I = ⊤
\end{leancode}
\begin{dictionary}
\textbf{One-sided by definition.} \lean{IsIdeal} asks only $u_{(n)}I\subseteq I$ ($I$ in the
\emph{right} slot of the product; equivalently, $I$ is a submodule of the adjoint module, a
``left ideal''). The literature (FLM, Lepowsky--Li Def.~3.9.7, FHL) defines an ideal as a subspace
with $u_{(n)}I\subseteq I$ \emph{and} $I_{(n)}u\subseteq I$; Lepowsky--Li observe that for
an $L_{-1}$-stable subspace either condition implies the other by skew-symmetry. The library proves
this (Theorem~\ref{l04:thm:twosided}), so \lean{IsIdeal} is the textbook notion and \lean{IsSimple}
is the textbook simplicity, not a weaker or stronger variant. (Chapter~\ref{ch:10} raised the
question whether \lean{IsSimple} is quantified over left ideals only: it is, and by
Theorem~\ref{l04:thm:twosided} this is the same class of subspaces.)

\textbf{Not graded by definition.} No gradedness is asked; it is automatic
(Theorem~\ref{l04:thm:graded}).

\textbf{Nonzero.} \lean{(⊤ : Submodule ℂ 𝒱) ≠ ⊥} is $\lfourV\ne0$, which excludes the zero
algebra (which has only the ideal $0 = \lfourV$). For a vertex operator algebra this is equivalent
to $\lfourvac\ne0$ (\lean{IsSimple.}\allowbreak\lean{vacuum_}\allowbreak\lean{ne_}\allowbreak\lean{zero} in \leanpath{VOA/\allowbreak Core/\allowbreak SimpleImage.lean}),
since every $u = u_{(-1)}\lfourvac$.

A deprecated duplicate \lean{IsSimple} exists in \leanpath{VOA/\allowbreak Modules/\allowbreak Rationality.lean}
(Chapter~\ref{ch:L12}); every statement of this chapter uses the one above.
\end{dictionary}
\begin{physics}
An ideal is a set of states closed under the action of every field: once such states appear they can
only produce more of themselves, so they form a sector that decouples from correlation functions with
the vacuum. A simple chiral algebra has no such sector: the vacuum module is irreducible and all
states are ``physical''. Quotienting by an ideal is setting null states to zero.
\end{physics}
\begin{status}
Definitions.
\end{status}

\begin{theorem}[Ideals are two-sided]\label{l04:thm:twosided}
If $I$ is an ideal, then $u_{(n)}v\in I$ for all $u\in I$, $v\in\lfourV$, $n\in\Z$.
\end{theorem}
\noindent File \leanpath{VOA/\allowbreak Core/\allowbreak QuotientVOA.lean}.
\begin{leancode}
def AbsorbsLeft (W : VertexOperatorAlgebra 𝒱) (I : Submodule ℂ 𝒱) : Prop :=
  ∀ (u v : 𝒱) (n : ℤ), u ∈ I → W.smul u n v ∈ I

theorem absorbsLeft_of_isIdeal (hI : W.IsIdeal I) : AbsorbsLeft W I := by
\end{leancode}
\begin{dictionary}
\lean{AbsorbsLeft} is the other half of the two-sided condition; the theorem says
\lean{IsIdeal} implies it, unconditionally.
\end{dictionary}
\begin{physics}
Locality: the OPE $u(z)v(w)$ and $v(w)u(z)$ are the same function, so if the states of $I$ are
closed under being acted on by fields, they are also closed when they act on other states. This is
the chiral form of the statement that an operator ideal of a local quantum field theory is
automatically two-sided.
\end{physics}
\emph{Proof on paper.} An ideal is stable under $T = L_{-1} = \omega_{(0)}$, hence under all
$T^j$. Skew-symmetry gives $u_{(n)}v = \sum_{j\ge0}\frac{(-1)^{n+j+1}}{j!}T^j(v_{(n+j)}u)$, a
finite sum of elements $T^j(v_{(n+j)}u)$ with $u\in I$ in the right slot, all in $I$. \qed
\begin{status}
Unconditional; axiom-clean.
\end{status}

\begin{theorem}[Ideals are graded]\label{l04:thm:graded}
Every ideal $I$ satisfies $I = \bigoplus_n(I\cap\lfourV_n)$.
\end{theorem}
\noindent File \leanpath{VOA/\allowbreak Core/\allowbreak QuotientVOA.lean}.
\begin{leancode}
def IsGradedIdeal (W : VertexOperatorAlgebra 𝒱) (I : Submodule ℂ 𝒱) : Prop :=
  I = ⨆ n : ℤ, I ⊓ W.grading n

theorem isGradedIdeal_of_isIdeal (hI : W.IsIdeal I) : IsGradedIdeal W I := by
\end{leancode}
\begin{dictionary}
\lean{⨆ n, I ⊓ W.grading n} is the sum of the homogeneous pieces $I\cap\lfourV_n$; since the
$\lfourV_n$ are independent this sum is direct, so the equation is ``$I$ is graded''. Despite its
name, \lean{IsGradedIdeal} does not include \lean{IsIdeal}; it is a property of subspaces.
\end{dictionary}
\begin{physics}
An ideal is stable under $L_0$, the Hamiltonian; a subspace stable under a diagonalisable
Hamiltonian is spanned by energy eigenstates. Null states can always be taken of definite energy.
\end{physics}
\emph{Proof on paper.} $L_0 = \omega_{(1)}$ maps $I$ into $I$ and acts on $\lfourV_n$ by $n$. If $x =
\sum_{n\in S}x_n\in I$ with $S$ finite, then $(L_0-k)x = \sum_{n\ne k}(n-k)x_n\in I$ for $k\in S$;
by induction on $|S|$ every $x_n\in I$. \qed
\begin{status}
Unconditional; axiom-clean.
\end{status}

\subsection{The quotient vertex operator algebra}

\begin{definition}[Quotient VOA]\label{l04:def:quot}
For an ideal $I$, $\lfourV/I$ with $Y([u],z)[v] = [Y(u,z)v]$, vacuum $[\lfourvac]$, conformal vector
$[\omega]$, central charge $c$ and grading $(\lfourV/I)_n = $ image of $\lfourV_n$.
\end{definition}
\noindent File \leanpath{VOA/\allowbreak Core/\allowbreak QuotientVOA.lean}.
\begin{leancode}
def quotientVOA (W : VertexOperatorAlgebra 𝒱) (I : Submodule ℂ 𝒱)
    (hI : W.IsIdeal I) (hgr : IsGradedIdeal W I) : VertexOperatorAlgebra (𝒱 ⧸ I) where
  grading n := (W.grading n).map I.mkQ
  isInternal := isInternal_map W I hgr
  finrank n := finiteDimensional_map W I n
  Y := Yq hI
  vacuum := I.mkQ W.vacuum

theorem quotientVOA_mkQ_smul (hI : W.IsIdeal I) (hgr : IsGradedIdeal W I)
    (u : 𝒱) (n : ℤ) (v : 𝒱) :
    (quotientVOA W I hI hgr).smul (I.mkQ u) n (I.mkQ v) = I.mkQ (W.smul u n v) :=

def quotientVOA' (W : VertexOperatorAlgebra 𝒱) (I : Submodule ℂ 𝒱) (hI : W.IsIdeal I) :
    VertexOperatorAlgebra (𝒱 ⧸ I) :=
  quotientVOA W I hI (isGradedIdeal_of_isIdeal hI)
\end{leancode}
\begin{dictionary}
\lean{Yq hI} is the induced vertex operator: on classes, \lean{quotientVOA_}\allowbreak\lean{mkQ_}\allowbreak\lean{smul} says
$[u]_{(n)}[v] = [u_{(n)}v]$, which is the defining property. Well-definedness in the left slot uses
Theorem~\ref{l04:thm:twosided}. The omitted fields (conformal vector, central charge, the axioms)
are the images of those of $\lfourV$, and \lean{quotientVOA_}\allowbreak\lean{centralCharge} records that $c$ is
unchanged.

\textbf{Which hypothesis is used where.} \lean{hgr} is used only for \lean{isInternal} (the image
grading is a direct sum only if $I$ is graded) and is automatic (\lean{quotientVOA'}).
\emph{Properness ($\lfourvac\notin I$) is not assumed}: for $I = \lfourV$ the quotient is the zero
space, which satisfies all axioms. So a theorem that concludes ``\lean{quotientVOA W I …} is a VOA''
says nothing about nontriviality; nontriviality has to be proved separately (as in
Theorem~\ref{l04:thm:maxsimple}).
\end{dictionary}
\begin{physics}
The quotient is the chiral algebra in which the states of $I$ are declared null and set to zero.
The central charge is unchanged because the stress tensor is the image of the old one; the OPE of
two surviving fields is the old OPE with the null part discarded.
\end{physics}
\begin{status}
Definition; \texttt{\#print axioms} on \lean{quotientVOA}: \lean{[propext, Classical.choice, Quot.sound]}.
\end{status}

\subsection{The maximal ideal}

\begin{definition}[Maximal proper graded ideal]\label{l04:def:maxideal}
$\lfourm(\lfourV) :=$ the sum of all graded ideals not containing $\lfourvac$; $L(\lfourV) :=
\lfourV/\lfourm(\lfourV)$.
\end{definition}
\noindent File \leanpath{VOA/\allowbreak Core/\allowbreak QuotientVOA.lean}.
\begin{leancode}
def ProperGradedIdeals (W : VertexOperatorAlgebra 𝒱) : Set (Submodule ℂ 𝒱) :=
  {I | W.IsIdeal I ∧ IsGradedIdeal W I ∧ W.vacuum ∉ I}

def maxGradedIdeal (W : VertexOperatorAlgebra 𝒱) : Submodule ℂ 𝒱 :=
  sSup (ProperGradedIdeals W)

def maxQuotientVOA (W : VertexOperatorAlgebra 𝒱) :
    VertexOperatorAlgebra (𝒱 ⧸ maxGradedIdeal W) :=
  quotientVOA W (maxGradedIdeal W) (maxGradedIdeal_isIdeal W) (maxGradedIdeal_isGradedIdeal W)
\end{leancode}
\begin{dictionary}
``Proper'' is encoded as $\lfourvac\notin I$, which for an ideal is equivalent to $I\ne\lfourV$
(\lean{eq_}\allowbreak\lean{top_}\allowbreak\lean{of_}\allowbreak\lean{vacuum_}\allowbreak\lean{mem}: an ideal containing $\lfourvac$ contains every $u = u_{(-1)}\lfourvac$).
\lean{sSup} of the set is the sum of all its members. Since every ideal is graded, ``graded'' is
redundant in \lean{ProperGradedIdeals}, so $\lfourm(\lfourV)$ is the sum of \emph{all} proper ideals
(\lean{le_}\allowbreak\lean{maxGradedIdeal'}). The literature's ``unique maximal proper ideal'' of a VOA of CFT type is
this object.
\end{dictionary}
\begin{physics}
$\lfourm(\lfourV)$ is the space of \emph{all} null states of the theory --- everything that can be
consistently set to zero --- and $L(\lfourV)$ is the physical chiral algebra. For the Virasoro vacuum
module at $c = c_{p,q}$ this passage gives the minimal model $L(c_{p,q},0)$; for $V^k(\mathfrak g)$ at
admissible level it gives $L_k(\mathfrak g)$.
\end{physics}
\begin{status}
Definitions.
\end{status}

\begin{theorem}[The maximal ideal is proper, the quotient is simple]\label{l04:thm:maxsimple}
If $\lfourvac\ne0$: (i) $\lfourvac\notin\lfourm(\lfourV)$; (ii) every ideal $J$ with $\lfourvac\notin J$
lies in $\lfourm(\lfourV)$; (iii) $L(\lfourV) = \lfourV/\lfourm(\lfourV)$ is simple.
\end{theorem}
\noindent File \leanpath{VOA/\allowbreak Core/\allowbreak QuotientVOA.lean}.
\begin{leancode}
theorem vacuum_not_mem_maxGradedIdeal (W : VertexOperatorAlgebra 𝒱) (hvac : W.vacuum ≠ 0) :
    W.vacuum ∉ maxGradedIdeal W := by

theorem le_maxGradedIdeal' {J : Submodule ℂ 𝒱} (hJ : W.IsIdeal J) (hv : W.vacuum ∉ J) :
    J ≤ maxGradedIdeal W :=

theorem isSimple_maxQuotientVOA (W : VertexOperatorAlgebra 𝒱) (hvac : W.vacuum ≠ 0)
    (hgraded : ∀ J : Submodule ℂ (𝒱 ⧸ maxGradedIdeal W), (maxQuotientVOA W).IsIdeal J →
      IsGradedIdeal W (J.comap (maxGradedIdeal W).mkQ)) :
    (maxQuotientVOA W).IsSimple := by

theorem isSimple_maxQuotientVOA' (W : VertexOperatorAlgebra 𝒱) (hvac : W.vacuum ≠ 0) :
    (maxQuotientVOA W).IsSimple := by
\end{leancode}
\begin{dictionary}
\lean{hvac : W.vacuum ≠ 0} is the only hypothesis of the unconditional form
\lean{isSimple_}\allowbreak\lean{maxQuotientVOA'}; it is necessary, since for $\lfourvac = 0$ the algebra is zero and
so is its quotient, which is not simple. The primed theorem is the unprimed one with \lean{hgraded}
discharged by Theorem~\ref{l04:thm:graded}; the unprimed form is a historical artefact whose docstring
still says the gradedness fact is ``not proved in this repo'' --- that sentence is outdated, the
fact is \lean{isGradedIdeal_}\allowbreak\lean{of_}\allowbreak\lean{isIdeal}.

\textbf{Why properness of the sum is not trivial.} A sum of proper ideals need not be proper in a
general algebra. Here it is, because every proper graded ideal meets $\lfourV_0 = \C\lfourvac$ in $0$ and
so lies in $\bigoplus_{n\ne0}\lfourV_n$, and a sum of subspaces of $\bigoplus_{n\ne0}\lfourV_n$
stays there. This is where CFT type ($\lfourV_0 = \C\lfourvac$) is essential.

\textbf{Could the conclusion be vacuous?} \lean{IsSimple} contains $\lfourV/\lfourm\ne0$, which is
(i). It is the honest statement.
\end{dictionary}
\begin{physics}
Removing all null states from the vacuum module leaves an irreducible chiral algebra: a unique
vacuum, no further decoupled sector. The vacuum itself is never null because it has weight $0$ and
the theory has a unique weight-zero state.
\end{physics}
\emph{Proof on paper.} (i), (ii): as in the fidelity check. (iii): an ideal $\bar J$ of the quotient
pulls back to an ideal $J\supseteq\lfourm$ of $\lfourV$; if $\lfourvac\notin J$ then $J\subseteq\lfourm$ by
(ii), so $\bar J = 0$; if $\lfourvac\in J$ then $J = \lfourV$ and $\bar J$ is everything. \qed
\begin{status}
Unconditional (given $\lfourvac\ne0$); \texttt{\#print axioms} on \lean{isSimple_}\allowbreak\lean{maxQuotientVOA'}
and \lean{vacuum_}\allowbreak\lean{not_}\allowbreak\lean{mem_}\allowbreak\lean{maxGradedIdeal}: \lean{[propext, Classical.choice, Quot.sound]}.
\end{status}

%======================================================================
\section{Homomorphisms and the transfer lemma}\label{l04:sec:hom}
%======================================================================

\begin{definition}[Homomorphism]\label{l04:def:hom}
A linear map $\varphi:\lfourV\to T$ between VOAs is a homomorphism if $\varphi\lfourvac_V =
\lfourvac_T$ and $\varphi(u_{(n)}v) = (\varphi u)_{(n)}(\varphi v)$ for all $u,v,n$.
\end{definition}
\noindent File \leanpath{VOA/\allowbreak Core/\allowbreak SimpleImage.lean}.
\begin{leancode}
structure IsVOAHom (V : VertexOperatorAlgebra 𝒱) (T : VertexOperatorAlgebra 𝒲)
    (φ : 𝒱 →ₗ[ℂ] 𝒲) : Prop where
  map_vacuum : φ V.vacuum = T.vacuum
  map_smul : ∀ (u : 𝒱) (n : ℤ) (v : 𝒱), φ (V.smul u n v) = T.smul (φ u) n (φ v)

theorem isIdeal_ker (hφ : IsVOAHom V T φ) : V.IsIdeal (LinearMap.ker φ) := by
\end{leancode}
\begin{dictionary}
A property of a linear map (\lean{Prop}-valued structure). It does \emph{not} require
$\varphi\omega_V = \omega_T$ nor preservation of the grading: these are vertex-algebra homomorphisms
(Kac's notion), not VOA homomorphisms in the sense of FLM. The deviation is deliberate (for a
Sugawara vector landing in a free-field target the conformal vectors differ) and is the
\emph{weaker} hypothesis, so theorems assuming \lean{IsVOAHom} are \emph{stronger} than their
VOA-homomorphism versions. \lean{IsVOAHom.map_L} adds $\varphi\omega_V = \omega_T$ when needed.
\end{dictionary}
\begin{physics}
A map of chiral algebras preserving all OPEs and the vacuum: a realisation of one chiral algebra
inside another, as in a free-field (Wakimoto, Feigin--Fuchs) realisation or a conformal embedding.
The stress tensors may differ (as for a coset or an improved stress tensor), which is why $\omega$ is
not required to map to $\omega$.
\end{physics}
\begin{status}
Definition; \lean{isIdeal_ker} unconditional, axiom-clean.
\end{status}

\begin{theorem}[Transfer lemma]\label{l04:thm:transfer}
Let $\varphi:\lfourV\to T$ be a surjective homomorphism onto a simple VOA $T$. Then $\ker\varphi =
\lfourm(\lfourV)$, so $\varphi$ induces an isomorphism $L(\lfourV)\cong T$, and $\varphi$ kills every
vector of every proper ideal of $\lfourV$.
\end{theorem}
\noindent File \leanpath{VOA/\allowbreak Core/\allowbreak SimpleImage.lean}.
\begin{leancode}
theorem ker_eq_maxGradedIdeal (hφ : IsVOAHom V T φ) (hsurj : Function.Surjective φ)
    (hsimple : T.IsSimple) : LinearMap.ker φ = QuotVOA.maxGradedIdeal V :=

theorem apply_eq_zero_of_mem_properIdeal (hφ : IsVOAHom V T φ)
    (hsurj : Function.Surjective φ) (hsimple : T.IsSimple)
    {J : Submodule ℂ 𝒱} (hJ : V.IsIdeal J) (hJv : V.vacuum ∉ J)
    {x : 𝒱} (hx : x ∈ J) : φ x = 0 :=

noncomputable def maxQuotientEquiv (hφ : IsVOAHom V T φ) (hsurj : Function.Surjective φ)
    (hsimple : T.IsSimple) : (𝒱 ⧸ QuotVOA.maxGradedIdeal V) ≃ₗ[ℂ] 𝒲 :=
\end{leancode}
\begin{dictionary}
Hypotheses: $\varphi$ a homomorphism (Definition~\ref{l04:def:hom}), surjective, $T$ simple. No
properness or ``$\varphi\ne0$'' hypothesis is needed: simplicity of $T$ gives $\lfourvac_T\ne0$, hence
$\lfourvac_V\ne0$ and $\lfourvac_V\notin\ker\varphi$. The isomorphism \lean{maxQuotientEquiv} is a
linear isomorphism; \lean{isVOAHom_}\allowbreak\lean{maxQuotientHom} shows the underlying map is a homomorphism from
\lean{maxQuotientVOA V}. Two non-vacuity witnesses in the same file
(\lean{maxGradedIdeal_}\allowbreak\lean{eq_}\allowbreak\lean{bot_}\allowbreak\lean{of_}\allowbreak\lean{isSimple}, from $\varphi = \mathrm{id}$, and
\lean{ker_}\allowbreak\lean{mkQ_}\allowbreak\lean{max_}\allowbreak\lean{eq_}\allowbreak\lean{maxGradedIdeal}, from $\varphi$ the projection) re-derive known facts through
the lemma, which would fail if an inclusion had been stated backwards.
\end{dictionary}
\begin{physics}
If a chiral algebra is realised inside another and the image is irreducible, then the realisation
kills exactly the null states: every null vector (for instance the Malikov--Feigin--Fuchs singular
vector of an affine algebra at admissible level) vanishes automatically in the realisation, without
computing it. This is how free-field realisations identify simple quotients.
\end{physics}
\emph{Proof on paper.} $\ker\varphi$ is an ideal not containing $\lfourvac$, so $\ker\varphi\subseteq
\lfourm(\lfourV)$. Conversely $\varphi(\lfourm(\lfourV))$ is an ideal of $T$ (by surjectivity) not containing
$\lfourvac_T$ (if $\lfourvac_T = \varphi(m)$ with $m\in\lfourm$, then $\lfourvac-m\in\ker\varphi\subseteq\lfourm$,
so $\lfourvac\in\lfourm$, contradiction), hence $0$ by simplicity. \qed
\begin{status}
Unconditional given its hypotheses; \texttt{\#print axioms} on \lean{ker_}\allowbreak\lean{eq_}\allowbreak\lean{maxGradedIdeal},
\lean{apply_}\allowbreak\lean{eq_}\allowbreak\lean{zero_}\allowbreak\lean{of_}\allowbreak\lean{mem_}\allowbreak\lean{properIdeal}, \lean{maxQuotientEquiv}: \lean{[propext, Classical.choice, Quot.sound]}.
\end{status}

%======================================================================
\section{Subalgebras generated by a set}\label{l04:sec:gen}
%======================================================================

We use the structure
\lean{VOASubalgebra W} (Definition~\ref{l04:def:subalg} below): a subspace containing $\lfourvac$ and $\omega$, closed under all products and
under the projections \lean{piHom W n}; \lean{VOASubalgebra.}\allowbreak\lean{toVOA} makes it a VOA with the same
central charge.

\begin{definition}[Vertex operator subalgebra]\label{l04:def:subalg}
A \emph{vertex operator subalgebra} of a VOA $\lfourV$ is a linear subspace $S\subseteq\lfourV$ containing $\lfourvac$ and
$\omega$, closed under all products $u_{(n)}v$ ($u,v\in S$, $n\in\Z$), and \emph{graded}: $S = \bigoplus_n(S\cap\lfourV_n)$,
equivalently every homogeneous component of every element of $S$ lies in $S$.
\end{definition}
\noindent File \leanpath{VOA/\allowbreak Modularity/\allowbreak Twisted.lean}.
\begin{leancode}
structure VOASubalgebra (W : VertexOperatorAlgebra 𝒱) where
  carrier : Submodule ℂ 𝒱
  vacuum_mem : W.vacuum ∈ carrier
  conformal_mem : W.conformalVector ∈ carrier
  smul_mem : ∀ u ∈ carrier, ∀ (n : ℤ), ∀ v ∈ carrier, W.smul u n v ∈ carrier
  graded : ∀ (n : ℤ), ∀ v ∈ carrier, piHom W n v ∈ carrier
\end{leancode}
(\lean{piHom W n} is the projection $\lfourV\to\lfourV_n$ along the direct sum decomposition,
\leanpath{VOA/\allowbreak Modules/\allowbreak FiniteGenerationProof.lean}.)

\begin{dictionary}
\lean{carrier : Submodule ℂ 𝒱} is the subspace $S$; \lean{smul_mem} is closure under all integer products
(closure under all $u_{(n)}$, $n\in\Z$, is the full closure under the vertex operators); \lean{graded} says that the projection onto each
$\lfourV_n$ keeps one inside $S$, which is the same as $S = \bigoplus(S\cap\lfourV_n)$. \textbf{Containing $\omega$ is part of
the definition:} these are subalgebras \emph{with the same conformal vector} (``full'' subalgebras), as is needed for
orbifolds and for generated subalgebras (Definition~\ref{l04:def:generated}); the commutant or coset subalgebras, which
carry a different conformal vector, are not of this kind. The notion is not vacuous: $S=\lfourV$ is one example, every generated subalgebra
(Definition~\ref{l04:def:generated}) is another, and the fixed points of a group of automorphisms are a third (the
orbifold, treated with the Moonshine construction).
\end{dictionary}

\begin{physics}
A vertex operator subalgebra is a set of local operators closed under the operator product expansion and
containing the same vacuum and the same stress tensor: a sub-chiral algebra of the same theory. Examples are the
subalgebra generated by a subset of the currents together with $T$, and the fixed-point algebra of an orbifold. Read
in the reverse direction the large algebra is an extension of the small one, as a lattice VOA is an extension of the one
attached to a sublattice.
\end{physics}

\begin{status}
Definition. Axiom-clean (\texttt{\#print axioms} \lean{VOA.VOASubalgebra} $=$
\lean{[propext, Classical.}\allowbreak\lean{choice, Quot.}\allowbreak\lean{sound]}).
\end{status}

\begin{theorem}[A subalgebra is a VOA]\label{l04:thm:subalg}
A vertex operator subalgebra $S\subseteq\lfourV$, with the restricted grading $S_n = S\cap\lfourV_n$, the restricted vertex
operator $Y_S(u,z) = Y(u,z)|_S$, the same vacuum $\lfourvac$, the same conformal vector $\omega$ and the same central
charge, is a vertex operator algebra.
\end{theorem}
\noindent File \leanpath{VOA/\allowbreak Modularity/\allowbreak Twisted.lean}.
\begin{leancode}
noncomputable def toVOA : VertexOperatorAlgebra S.carrier where
  grading := S.gradingS
  isInternal := S.gradingS_isInternal
  finrank := S.gradingS_finrank
  Y := S.YS
  vacuum := ⟨W.vacuum, S.vacuum_mem⟩

theorem toVOA_centralCharge : S.toVOA.centralCharge = W.centralCharge
\end{leancode}
(The remaining fields of \lean{toVOA} are proofs of the axioms and are omitted; \lean{S.gradingS n} is $S\cap\lfourV_n$ regarded
as a subspace of $S$, and \lean{S.YS u} is the restriction of $Y(u,z)$ to $S$, both in the same file.)

\begin{dictionary}
The new VOA lives on the \emph{type} \lean{S.carrier}, the subspace regarded as a vector space in its own right
(printed \lean{↥S.carrier} in some contexts), so the vacuum has to be given as a pair
\lean{⟨W.vacuum, S.vacuum_mem⟩} (an element together with the proof that it lies in $S$). \lean{toVOA} is a
\emph{definition} producing a VOA, not a statement; its fidelity rests on its data being the restricted ones, which is
visible in the displayed lines: grading $S\cap\lfourV_n$, vertex operator $Y|_S$, vacuum $\lfourvac$. (Among the omitted fields,
the conformal vector is \lean{⟨W.conformalVector, S.conformal_mem⟩} and the central charge is \lean{W.centralCharge}.) \lean{toVOA_centralCharge} proves that the central charge is unchanged; its proof is
\lean{rfl}, because the definition sets it so. \textbf{Hypotheses:} none besides $S$ being a \lean{VOASubalgebra}; in particular
the finite-dimensionality and positivity of the graded pieces must come from $\lfourV$, and do. \textbf{Could the
construction be vacuous or cheating?} The structure \lean{VertexOperatorAlgebra} has all the axioms of
Definition~\ref{l03:def:VOA}; they are discharged field by field in the source (not assumed). In particular, CFT type
holds because $S_0\subseteq\lfourV_0 = \C\lfourvac\subseteq S$.
\end{dictionary}

\begin{physics}
A subalgebra that carries the same stress tensor is itself a chiral algebra with the same central charge: the
operator product expansion of the sub-chiral algebra closes, the same $T(z)$ generates its conformal transformations,
so the conformal anomaly $c$ (Casimir energy $-c/24$ on the cylinder) is the same. This is how one passes from a CFT to
a sub-CFT: the weights stay the same because the spectrum of $L_0$ on $S$ is the restriction of the one on $\lfourV$, so the
discrete, finite-degeneracy spectrum is inherited.
\end{physics}

\textbf{Proof on paper.} Every axiom is inherited. The grading $S=\bigoplus_n(S\cap\lfourV_n)$ is direct and has finite-dimensional
pieces because $S\cap\lfourV_n\subseteq\lfourV_n$ and $\lfourV_n$ is finite-dimensional, and $S\cap\lfourV_n = 0$ for $n<0$;
\lean{graded} gives that the sum is all of $S$. The operators $Y(u,z)$, $u\in S$, preserve $S$ because $S$ is closed under all
products. Vacuum axiom, creation, regularity, locality and the Jacobi/Borcherds relations are identities between
operators on $\lfourV$, so they hold for the restrictions; the Virasoro relations hold because $\omega\in S$ and the modes
$L_n = \omega_{(n+1)}$ restrict; $L_0$ acts on $S\cap\lfourV_n$ by $n$; translation $Y(L_{-1}v,z) = \partial_zY(v,z)$ restricts;
and CFT type holds because $S_0\subseteq\lfourV_0 = \C\lfourvac$ contains $\lfourvac$. \qed

\begin{status}
Unconditional. \texttt{\#print axioms} of \lean{VOA.VOASubalgebra.toVOA} and \lean{VOA.VOASubalgebra.toVOA\_centralCharge}
are each \lean{[propext, Classical.}\allowbreak\lean{choice, Quot.}\allowbreak\lean{sound]} (run for this chapter).
\end{status}

\begin{definition}[Generated subalgebra]\label{l04:def:generated}
For $S\subseteq\lfourV$, $\langle S\rangle$ is the intersection of all vertex operator subalgebras
(in the sense just recalled) containing $S$.
\end{definition}
\noindent File \leanpath{VOA/\allowbreak Core/\allowbreak VOASubalgebraGenerated.lean}.
\begin{leancode}
def generated (S : Set 𝒱) : VOASubalgebra W :=
  sInf {A : VOASubalgebra W | S ⊆ A.carrier}

theorem generated_le {S : Set 𝒱} (A : VOASubalgebra W) (hA : S ⊆ A.carrier) :
    (generated W S).carrier ≤ A.carrier :=

theorem mem_generated {S : Set 𝒱} {motive : 𝒱 → Prop}
    (base : ∀ x ∈ S, motive x)
    (zero : motive 0)
    (add : ∀ x y : 𝒱, motive x → motive y → motive (x + y))
    (smul : ∀ (c : ℂ) (x : 𝒱), motive x → motive (c • x))
    (vac : motive W.vacuum)
    (conf : motive W.conformalVector)
    (prod : ∀ u : 𝒱, motive u → ∀ (n : ℤ) (v : 𝒱), motive v → motive (W.smul u n v))
    (proj : ∀ (n : ℤ) (v : 𝒱), motive v → motive (piHom W n v))
    {x : 𝒱} (hx : x ∈ (generated W S).carrier) : motive x := by

noncomputable def generatedVOA (S : Set 𝒱) :
    VertexOperatorAlgebra (generated W S).carrier :=
  (generated W S).toVOA
\end{leancode}
\begin{dictionary}
\lean{sInf} uses the intersection structure \lean{instInfSetVOASubalgebra} proved in the same file.
\lean{generated_le} is minimality; \lean{mem_generated} is the induction principle (every property
true on $S$, $0$, $\lfourvac$, $\omega$ and stable under sums, scalars, products and projections holds on
$\langle S\rangle$).

\textbf{Deviation from the literature.} The usual ``vertex subalgebra generated by $S$'' (Kac,
Lepowsky--Li §3.9) is the smallest subspace containing $S$ and $\lfourvac$ and closed under products; it
need not contain $\omega$. Here \lean{VOASubalgebra} always contains $\omega$, so \lean{generated W S}
is the vertex subalgebra generated by $S\cup\{\omega\}$, closed in addition under the projections.
The projection clause adds nothing: the vertex subalgebra generated by $S\cup\{\omega\}$ is stable
under $L_0 = \omega_{(1)}$, hence graded by the argument of Theorem~\ref{l04:thm:graded}. But the
$\omega$ clause is a real difference: for instance the subalgebra generated by a Heisenberg current $h$
inside a lattice VOA contains the full stress tensor $\omega$, not only the Heisenberg one. When a
statement needs a subalgebra with a \emph{different} conformal vector (a coset, a commutant) this
notion is the wrong one.
\end{dictionary}
\begin{physics}
The sub-chiral algebra generated by a set of fields \emph{together with the stress tensor}: all normal
ordered products and derivatives of the given fields and of $T(z)$. Always containing $T(z)$ means
the subalgebra has the same central charge as the ambient theory.
\end{physics}
\begin{status}
Definitions and an induction principle; axiom-clean.
\end{status}

%======================================================================
\section{Screening operators and their kernels}\label{l04:sec:screening}
%======================================================================

\subsection{Derivations with subalgebra kernels}

\begin{definition}[VOA derivation]\label{l04:def:voader}
$D\in\End(\lfourV)$ is a VOA derivation if $D(u_{(n)}v) = (Du)_{(n)}v + u_{(n)}(Dv)$ for all $u,n,v$,
$D\lfourvac = 0$, $D\omega = 0$, and $D$ commutes with every projection onto $\lfourV_n$.
\end{definition}
\noindent File \leanpath{VOA/\allowbreak Core/\allowbreak ScreeningKernel.lean}.
\begin{leancode}
structure IsVOADerivation (W : VertexOperatorAlgebra 𝒱) (D : End ℂ 𝒱) : Prop where
  leibniz : ∀ (u : 𝒱) (n : ℤ) (v : 𝒱),
      D (W.smul u n v) = W.smul (D u) n v + W.smul u n (D v)
  vacuum : D W.vacuum = 0
  conformal : D W.conformalVector = 0
  piHom_comm : ∀ (n : ℤ) (v : 𝒱), D (piHom W n v) = piHom W n (D v)

def kerSubalgebra (W : VertexOperatorAlgebra 𝒱) (D : End ℂ 𝒱)
    (hD : IsVOADerivation W D) : VOASubalgebra W where
  carrier := LinearMap.ker D
\end{leancode}
\begin{dictionary}
The four clauses are exactly what makes $\ker D$ a \lean{VOASubalgebra}: Leibniz gives closure
under products, the next two give $\lfourvac,\omega\in\ker D$, and \lean{piHom_comm} gives gradedness.
\lean{kerSubalgebra} is then a definition whose remaining fields are short proofs.
\end{dictionary}
\begin{physics}
The states annihilated by a symmetry generator that commutes with the stress tensor form a closed
sub-chiral algebra: the ``invariant'' operators.
\end{physics}
\begin{status}
Definitions; axiom-clean.
\end{status}

\subsection{Screening operators}

\begin{definition}[Screening operator]\label{l04:def:screening}
For $a\in\lfourV$, $Q_a := a_{(0)}$.
\end{definition}
\noindent File \leanpath{VOA/\allowbreak Core/\allowbreak ScreeningKernel.lean}.
\begin{leancode}
def screening (W : VertexOperatorAlgebra 𝒱) (a : 𝒱) : End ℂ 𝒱 := mode (W.Y a) 0
\end{leancode}
\begin{dictionary}
Literally $a_{(0)}$, the zero mode \lean{W.zeroMode a 1} of a weight-one state
(\lean{screening_}\allowbreak\lean{eq_}\allowbreak\lean{zeroMode}); on $\lfourV_1$ it is the adjoint action $[a,\cdot]$
(\lean{screening_}\allowbreak\lean{eq_}\allowbreak\lean{bracketV1}).

\textbf{Relation to screening operators in the literature.} In the Coulomb gas (Feigin--Fuchs,
Dotsenko--Fateev) a screening operator is $\oint e^{\alpha\phi}(z)dz$, the zero mode of a field
$e^{\alpha\phi}$ of weight one \emph{for a deformed (background-charge) stress tensor} $T^\lambda$, and it
usually acts between \emph{different} Fock modules. The library's notion is the special case in
which $a$ is a weight-one state of the vertex operator algebra itself, for its own $\omega$. The source
records honestly that the concrete triplet screening operator of
\leanpath{VOA/\allowbreak Examples/\allowbreak TripletScreening.lean} fits this definition only when
$\lambda = 0$ and $\langle\alpha,\alpha\rangle = 2$, because the background-charge VOA structure on
the lattice module is not in the repository.
\end{dictionary}
\begin{physics}
A screening charge is the contour integral of a weight-one field. Because it has weight one, the
integral is conformally invariant and commutes with the Virasoro algebra; the operators it
annihilates (its ``commutant'') form a chiral algebra --- the Virasoro minimal models, the $W$-algebras
and the triplet algebras are all defined as such commutants inside free-field theories.
\end{physics}
\begin{status}
Definition.
\end{status}

\begin{theorem}[A weight-one state gives a VOA derivation]\label{l04:thm:scrder}
For every $a$, $Q_a$ is a derivation of every product and $Q_a\lfourvac = 0$. If $a\in\lfourV_1$, then
moreover $Q_a$ preserves each $\lfourV_n$ and $Q_a\omega = 0$; hence $Q_a$ is a VOA derivation.
\end{theorem}
\noindent File \leanpath{VOA/\allowbreak Core/\allowbreak ScreeningKernel.lean}.
\begin{leancode}
theorem screening_leibniz (W : VertexOperatorAlgebra 𝒱) (a : 𝒱) (u : 𝒱) (n : ℤ) (v : 𝒱) :
    screening W a (W.smul u n v)
      = W.smul (screening W a u) n v + W.smul u n (screening W a v) := by

theorem screening_conformalVector (W : VertexOperatorAlgebra 𝒱) {a : 𝒱}
    (ha : a ∈ W.grading 1) : screening W a W.conformalVector = 0 := by

theorem screening_isDerivation (W : VertexOperatorAlgebra 𝒱) {a : 𝒱}
    (ha : a ∈ W.grading 1) : IsVOADerivation W (screening W a) where
\end{leancode}
\begin{dictionary}
\lean{ha : a ∈ W.grading 1} is the only hypothesis, and it is needed only for grading preservation
and $Q_a\omega = 0$. \textbf{No primarity of $a$ is assumed}, and none is needed:
$a_{(0)}\omega = 0$ holds for every $a\in\lfourV_1$ in a VOA of CFT type (Remark~\ref{l04:rem:homega}). A
reader who expects ``screening currents must be primary'' should note that in a VOA of CFT type a
weight-one state is automatically annihilated by $L_n$, $n\ge2$, and $L_1a\in\C\lfourvac$; it is this
that makes the computation work.
\end{dictionary}
\begin{physics}
$[Q_a,L_n] = 0$: the charge of a weight-one current commutes with conformal transformations, so it
maps primaries to primaries of the same weight. $Q_a\omega = 0$ is the statement that the stress tensor
is neutral under the charge.
\end{physics}
\emph{Proof on paper.} Leibniz is Theorem~\ref{l04:thm:deriv}; $Q_a\lfourvac = 0$ is creation;
grading preservation is VA1 with weight $1$; $Q_a\omega = 0$ is the skew-symmetry computation of
Remark~\ref{l04:rem:homega}. \qed
\begin{status}
Unconditional; \texttt{\#print axioms} on \lean{screening_}\allowbreak\lean{conformalVector}: \lean{[propext, Classical.choice, Quot.sound]}.
\end{status}

\begin{definition}[Screening kernels]\label{l04:def:scrkernel}
For $a\in\lfourV_1$, $\ker Q_a$ as a vertex operator subalgebra; for a family $(a_i)$ of weight-one
states, $\bigcap_i\ker Q_{a_i}$.
\end{definition}
\noindent File \leanpath{VOA/\allowbreak Core/\allowbreak ScreeningKernel.lean}.
\begin{leancode}
def screeningKernel (W : VertexOperatorAlgebra 𝒱) (a : 𝒱) (ha : a ∈ W.grading 1) :
    VOASubalgebra W :=
  kerSubalgebra W (screening W a) (screening_isDerivation W ha)

def screeningKernelInter {ι : Type*} (W : VertexOperatorAlgebra 𝒱) (a : ι → 𝒱)
    (ha : ∀ i, a i ∈ W.grading 1) : VOASubalgebra W :=
  sInf (Set.range fun i => screeningKernel W (a i) (ha i))

theorem mem_screeningKernelInter {ι : Type*} {W : VertexOperatorAlgebra 𝒱} {a : ι → 𝒱}
    {ha : ∀ i, a i ∈ W.grading 1} {x : 𝒱} :
    x ∈ (screeningKernelInter W a ha).carrier ↔ ∀ i, screening W (a i) x = 0 := by
\end{leancode}
\begin{dictionary}
The family may be indexed by any type \lean{ι} (finite or not). \lean{screeningKernelVOA} and
\lean{screeningKernelInterVOA} are the corresponding VOAs (via \lean{toVOA}), with the same central
charge. Membership is literally ``killed by every $Q_{a_i}$''.
\end{dictionary}
\begin{physics}
The commutant of a set of screening charges: the chiral algebra of operators that commute with all of
them. This is the definition of a $W$-algebra by screenings, and of the minimal models as the
commutant of the two Dotsenko--Fateev screenings $Q_\pm$ (in the background-charge setting; see the
caveat above).
\end{physics}
\begin{status}
Definitions; axiom-clean.
\end{status}

\begin{theorem}[Screenings with total-derivative contact commute]\label{l04:thm:scrcommute}
If $a_{(0)}b = \lambda\,Tc$ for some $c$ and $\lambda\in\C$, then $Q_aQ_b = Q_bQ_a$.
\end{theorem}
\noindent File \leanpath{VOA/\allowbreak Core/\allowbreak ScreeningKernel.lean}.
\begin{leancode}
theorem screening_commute_of_zeroProduct_eq_translation
    (W : VertexOperatorAlgebra 𝒱) (a b c : 𝒱) (scalar : ℂ)
    (hcontact : W.smul a 0 b = scalar • W.toPreVertexAlgebra.T c) :
    Commute (screening W a) (screening W b) := by
\end{leancode}
\begin{dictionary}
\lean{Commute A B} is $AB = BA$. No weight hypothesis is needed.
\end{dictionary}
\begin{physics}
Two screening charges commute if the simple pole of their OPE is a total derivative, because the
contour integral of a derivative vanishes: $[Q_a,Q_b] = \oint(a_{(0)}b)(w)dw = \lambda\oint\partial c = 0$.
This is the standard check that the two Coulomb-gas screenings commute.
\end{physics}
\emph{Proof on paper.} $[Q_a,Q_b] = (a_{(0)}b)_{(0)} = \lambda (Tc)_{(0)}$ and $(Tc)_{(n)} = -nc_{(n-1)}$
vanishes at $n = 0$. \qed
\begin{status}
Unconditional; axiom-clean.
\end{status}

\begin{theorem}[Transfer lemma for a screening-kernel image]\label{l04:thm:scrtransfer}
Let $\varphi:\lfourV\to W$ be a homomorphism whose image is exactly $\ker Q_a$ for some $a\in W_1$.
If the VOA $\ker Q_a$ is simple, then $\ker\varphi = \lfourm(\lfourV)$.
\end{theorem}
\noindent File \leanpath{VOA/\allowbreak Core/\allowbreak SimpleImage.lean}.
\begin{leancode}
theorem ker_eq_maxGradedIdeal_of_screeningKernel {V : VertexOperatorAlgebra 𝒱}
    {W : VertexOperatorAlgebra 𝒲} {a : 𝒲} (ha : a ∈ W.grading 1) {φ : 𝒱 →ₗ[ℂ] 𝒲}
    (hφ : IsVOAHom V W φ) (hin : ∀ x, screening W a (φ x) = 0)
    (honto : ∀ y, screening W a y = 0 → ∃ x, φ x = y)
    (hsimple : (screeningKernelVOA W a ha).IsSimple) :
    LinearMap.ker φ = QuotVOA.maxGradedIdeal V := by
\end{leancode}
\begin{dictionary}
\lean{hin} and \lean{honto} together say $\operatorname{im}\varphi = \ker Q_a$. \lean{hsimple} is
simplicity of the kernel VOA, which in applications is the hard input. The proof corestricts
$\varphi$ to the subalgebra (\lean{isVOAHom_}\allowbreak\lean{codRestrict}) and applies Theorem~\ref{l04:thm:transfer}.
\end{dictionary}
\begin{physics}
If a universal chiral algebra maps onto the commutant of a screening charge and that commutant is
irreducible, the map identifies the commutant with the simple quotient. This is the logic of
Adamovi\'c's realisations of simple affine algebras, used in Chapter~\ref{ch:L19}.
\end{physics}
\begin{status}
Conditional on \lean{hsimple}; axiom-clean.
\end{status}

\begin{remark}[The super screening kernel]\label{l04:rem:super}
The same file defines \lean{superScreening} on a super near-VOA and
\lean{superScreeningKernel W a h} for \lean{h : IsSuperScreening W a}, a structure whose fields
\lean{leibniz} and \lean{conformal} are the super Leibniz rule and $a_{(0)}\omega = 0$ \emph{as
hypotheses} (the super Borcherds machinery needed to prove them is not in the library). So the super
statement is genuinely conditional on these two fields. Super vertex algebras are the subject of
Chapter~\ref{ch:L17}.
\end{remark}

%======================================================================
\section{Singular vectors}\label{l04:sec:singular}
%======================================================================

A singular vector is the algebraic shadow of a null state: a vector of positive weight that the
lowering operators kill. The file \leanpath{VOA/\allowbreak Core/\allowbreak SingularVector.lean} is the
interface between explicit singular-vector computations (Chapters~\ref{ch:L06}, \ref{ch:L08}) and the
structural conclusion ``$\lfourV$ is not simple''. It distinguishes two notions, and the distinction is the
first thing a referee should understand.

\begin{definition}[Singular and strongly singular vectors]\label{l04:def:singular}
Let $k>0$. A vector $s\ne0$ in $\lfourV_k$ is \emph{(Virasoro) singular of weight $k$} if $L_ns = 0$
for all $n>0$. It is \emph{strongly singular of weight $k$} if $u_{(n)}s = 0$ for every homogeneous
$u\in\lfourV_m$ and every $n\ge m$.
\end{definition}
\noindent File \leanpath{VOA/\allowbreak Core/\allowbreak SingularVector.lean}.
\begin{leancode}
structure IsSingularOfWeight (W : VertexOperatorAlgebra 𝒱) (k : ℤ) (s : 𝒱) : Prop where
  ne_zero : s ≠ 0
  weight_pos : 0 < k
  homogeneous : s ∈ W.grading k
  annihilated : ∀ n : ℤ, 0 < n → W.L n s = 0

structure IsStronglySingularOfWeight (W : VertexOperatorAlgebra 𝒱) (k : ℤ) (s : 𝒱) : Prop where
  ne_zero : s ≠ 0
  weight_pos : 0 < k
  homogeneous : s ∈ W.grading k
  annihilated : ∀ (m : ℤ) (u : 𝒱), u ∈ W.grading m → ∀ n : ℤ, m ≤ n → W.smul u n s = 0
\end{leancode}
\begin{dictionary}
\textbf{Index convention.} $u\in\lfourV_m$ maps $\lfourV_k$ to $\lfourV_{m+k-n-1}$, so $n\ge m$ are
exactly the modes that strictly \emph{lower} the weight (the target weight is $\le k-1$). For $u =
\omega\in\lfourV_2$, $\omega_{(n)} = L_{n-1}$ with $n\ge2$, i.e.\ $L_j$ with $j\ge1$: strong singularity
implies Virasoro singularity (\lean{IsStronglySingularOfWeight.}\allowbreak\lean{toIsSingularOfWeight}).

\textbf{Nonvanishing and positive weight are required} (\lean{ne_zero}, \lean{weight_pos}). The
question raised in the instructions for this book --- whether lowest-weight-type data require
nonvanishing --- is answered here in the affirmative: $s = 0$ is excluded, so ``$\lfourV$ has a
singular vector'' is not trivially true. Positive weight excludes $s = \lfourvac$, which is killed by
all lowering modes.

\textbf{Comparison with the literature.} ``Singular vector'' in the Virasoro and affine literature
means a vector annihilated by the positive modes of the \emph{generators}; in a VOA generated by
those fields this is equivalent to strong singularity (positive modes of generators generate, by the
Borcherds identity, all weight-lowering modes). The library does not prove this equivalence --- the
source lists it as missing --- so the two notions are kept apart, and only strong singularity, or
Virasoro singularity with an extra pairing condition, is used to obstruct simplicity.
\end{dictionary}
\begin{physics}
A singular vector is a null state: a descendant that is also a highest-weight state. It has zero
norm and zero overlap with every state of the module. In the Virasoro vacuum module at $c = c_{p,q}$
there is one at level $(p-1)(q-1)$; setting it to zero is what makes the minimal model have finitely
many primaries.
\end{physics}
\begin{status}
Definitions.
\end{status}

\begin{warning}[title={Trap / dead end: Virasoro singularity alone does not make the generated ideal proper}]\label{l04:warn:primary}
The tempting argument ``everything in the ideal $(s)$ has weight $\ge\operatorname{wt}s>0$, and
$\lfourvac$ has weight $0$'' is false, because modes $u_{(n)}$ with $n\ge\operatorname{wt}u$ lower the
weight. The library's counterexample: in a simple affine VOA $L_k(\mathfrak g)$, $k\ne0$, every
$x\in\lfourV_1 = \mathfrak g$ is primary ($L_nx = 0$ for $n>0$), yet $y_{(1)}x = k(x|y)\lfourvac$ reaches the
vacuum in one step, so $(x) = \lfourV$. What fails is strong singularity: $y_{(1)}x\ne0$ with
$\operatorname{wt}y = 1\le1$. The weight argument is proved only one layer deep
(\lean{smul_mem_wtGE}: $u_{(n)}s$ has weight $\ge k$ for strongly singular $s$); turning it into
properness of $(s)$ would need $(s) = \operatorname{span}\{u_{(n)}s\}$ (Lepowsky--Li Prop.~4.5.6),
which is not in the library.
\end{warning}

\begin{definition}[Ideal generated by a vector]\label{l04:def:idealgen}
$(s) :=$ the intersection of all ideals containing $s$.
\end{definition}
\noindent File \leanpath{VOA/\allowbreak Core/\allowbreak SingularVector.lean}.
\begin{leancode}
def idealGen (W : VertexOperatorAlgebra 𝒱) (s : 𝒱) : Submodule ℂ 𝒱 :=
  sInf {I : Submodule ℂ 𝒱 | s ∈ I ∧ W.IsIdeal I}

theorem idealGen_isIdeal (W : VertexOperatorAlgebra 𝒱) (s : 𝒱) : W.IsIdeal (idealGen W s) := by
\end{leancode}
\begin{dictionary}
The family is nonempty ($\lfourV$ is in it) and an intersection of ideals is an ideal, so \lean{idealGen}
is the smallest ideal containing $s$. Standard.
\end{dictionary}
\begin{physics}
The null submodule generated by a null state: all states obtained from it by the action of fields.
\end{physics}
\begin{status}
Definition; axiom-clean.
\end{status}

\subsection{Two routes to properness}

Properness is obtained through invariant forms (Section~\ref{l04:sec:forms}, which the reader may
consult first): the radical of a normalised fully invariant form is a proper ideal, so any singular
vector in a radical generates a proper ideal. The library needs the radical in both slots.

\begin{definition}[Right radical]\label{l04:def:rightrad}
For a bilinear form $B$, $\lfourrad_R B = \{b : B(a,b) = 0\ \forall a\}$.
\end{definition}
\noindent File \leanpath{VOA/\allowbreak Core/\allowbreak SingularVector.lean}.
\begin{leancode}
def rightRadical (B : InvariantForm W) : Submodule ℂ 𝒱 := LinearMap.ker B.form.flip

theorem rightRadical_isIdeal {B : InvariantForm W} (hB : B.IsFullyInvariant) :
    W.IsIdeal (rightRadical B) := by
\end{leancode}
\begin{dictionary}
\lean{B.form.flip} is $(a,b)\mapsto B(b,a)$; its kernel is the right radical. \lean{InvariantForm} and
\lean{IsFullyInvariant} are Definitions~\ref{l04:def:invform} and~\ref{l04:def:fullinv} below. That the
right radical is an ideal is less immediate than for the left radical, because full invariance moves
modes from the left slot to the right; the proof is a downward induction on the weight of the acting
state.
\end{dictionary}
\begin{physics}
States orthogonal to everything from the other side of the two-point function; for a symmetric
form (which is the physical case) the two radicals coincide, and Theorem~\ref{l04:thm:lirad} shows
they coincide in general here.
\end{physics}
\begin{status}
Unconditional given full invariance; axiom-clean.
\end{status}

\begin{theorem}[Singular vectors obstruct simplicity]\label{l04:thm:notsimple}
Assume $\lfourvac\ne0$ and $L_1\lfourV_1 = 0$. (a) If $s$ is strongly singular of weight $k$, then $(s)$ is
a proper nonzero ideal, so $\lfourV$ is not simple. (b) If $s$ is Virasoro singular of weight $k$ and
$s_{(2k-1)}x = 0$ for all $x\in\lfourV_k$, the same conclusion holds.
\end{theorem}
\noindent File \leanpath{VOA/\allowbreak Core/\allowbreak SingularVector.lean}.
\begin{leancode}
theorem not_isSimple_of_stronglySingular_FHL (hne : W.vacuum ≠ 0)
    (hqp : ∀ v, v ∈ W.grading 1 → W.L 1 v = 0) {k : ℤ} {s : 𝒱}
    (hs : IsStronglySingularOfWeight W k s) : ¬ W.IsSimple :=

theorem not_isSimple_of_shapovalov_FHL (hne : W.vacuum ≠ 0)
    (hqp : ∀ v, v ∈ W.grading 1 → W.L 1 v = 0) {k : ℤ} {s : 𝒱} (hs : IsSingularOfWeight W k s)
    (hpair : ∀ x ∈ W.grading k, W.smul s (2 * k - 1) x = 0) : ¬ W.IsSimple :=
\end{leancode}
\begin{dictionary}
\lean{hne}, \lean{hqp} are the hypotheses of the canonical invariant form
(Theorem~\ref{l04:thm:canonical}); they enter only to produce a fully invariant normalised form, and
the versions without the suffix \lean{_FHL} take such a form \lean{B} instead. \lean{hpair} is the
\emph{Shapovalov condition}: $s_{(2k-1)}$ maps $\lfourV_k$ to $\lfourV_{k+k-(2k-1)-1} = \lfourV_0$, so
$s_{(2k-1)}x$ is the ``vacuum component'' of the pairing of $s$ with $x$, and the condition says $s$ is
in the radical of the level-$k$ Gram form. It is a finite linear condition on one weight space, which
is what an explicit computation can check. In (b) only $L_1s = 0$ is really used (to collapse the
invariance sum); the rest of Virasoro singularity is carried along.

The conclusion \lean{¬ W.IsSimple} is genuinely informative: it is false for simple VOAs, and the
hypotheses are satisfiable (Example~\ref{l04:ex:c0}).
\end{dictionary}
\begin{physics}
A null state with vanishing two-point function against its own energy level generates a decoupled
sector, so the theory built on the vacuum module is reducible; the physical theory is the quotient.
Condition (b) is how degenerate values of $c$ are found in practice: the Kac determinant at level $k$
vanishes.
\end{physics}
\emph{Proof on paper.} Let $B$ be the canonical form. (a) For $x\in\lfourV_p$ write $x = x_{(-1)}\lfourvac$
and expand $B(x,s) = B(x_{(-1)}\lfourvac,s)$ by full invariance: every term is $\pm\frac1{j!}
B(\lfourvac,(L_1^jx)_{(2p-1-j)}s)$, and $2p-1-j\ge p-j = \operatorname{wt}(L_1^jx)$ when $p\ge1$, so each
term vanishes by strong singularity; for $p = 0$, $B(\lfourvac,s) = 0$ by weight orthogonality. So $s$ is
in the right radical, a proper ideal (it misses $\lfourvac$ since $B(\lfourvac,\lfourvac) = 1$); $(s)$ is
contained in it, and $(s)\ni s\ne0$. (b) With $L_1s = 0$ the invariance sum for $B(s,x)$ collapses to
$(-1)^kB(\lfourvac,s_{(2k-1)}x)$ for $x\in\lfourV_k$, which is $0$ by \lean{hpair}; for $x$ in other weights
$B(s,x) = 0$ by orthogonality. So $s$ is in the left radical. \qed
\begin{status}
Conditional only on the displayed hypotheses; \texttt{\#print axioms} on both \lean{_FHL} theorems:
\lean{[propext, Classical.choice, Quot.sound]}.
\end{status}

\begin{example}[Central charge zero]\label{l04:ex:c0}
In the universal Virasoro VOA $V(c,0)$ (Chapter~\ref{ch:L06}), $\lfourV_1 = 0$ (so \lean{hqp} holds
trivially), $\lfourV_2 = \C\omega$, and $L_1\omega = 0$, $L_2\omega = \tfrac c2\lfourvac$. At $c = 0$,
$s = \omega$ is Virasoro singular of weight $2$, and the Shapovalov condition is
$\omega_{(3)}\omega = L_2\omega = 0$. Theorem~\ref{l04:thm:notsimple}(b) says $V(0,0)$ is not simple; its
simple quotient is $L(0,0) = \C\lfourvac$, the trivial theory.
\end{example}

\begin{definition}[Quotient by a singular vector]\label{l04:def:quotsing}
$\lfourV/(s)$ with the quotient structure.
\end{definition}
\noindent File \leanpath{VOA/\allowbreak Core/\allowbreak SingularVector.lean}.
\begin{leancode}
noncomputable def quotientBySingular (W : VertexOperatorAlgebra 𝒱) (s : 𝒱) :
    VertexOperatorAlgebra (𝒱 ⧸ idealGen W s) :=
  QuotVOA.quotientVOA' W (idealGen W s) (idealGen_isIdeal W s)

theorem le_maxGradedIdeal_of_proper {s : 𝒱} (hv : W.vacuum ∉ idealGen W s) :
    idealGen W s ≤ QuotVOA.maxGradedIdeal W :=
\end{leancode}
\begin{dictionary}
Defined for every $s$ (for $(s) = \lfourV$ it is the zero VOA; \lean{quotientBySingular_}\allowbreak\lean{vacuum_}\allowbreak\lean{ne_}\allowbreak\lean{zero}
gives nonvanishing when $(s)$ is proper). \lean{le_}\allowbreak\lean{maxGradedIdeal_}\allowbreak\lean{of_}\allowbreak\lean{proper} says that $\lfourV/(s)$
surjects onto $L(\lfourV)$. The converse --- that $(s) = \lfourm(\lfourV)$ for the Virasoro minimal-model
vector or for the affine admissible vector (Feigin--Fuchs, Kac--Kazhdan) --- is a representation-theoretic
input that the file explicitly does \emph{not} assume or prove.
\end{dictionary}
\begin{physics}
Setting the null vector to zero; whether this already gives the physical theory (no further null
states) is a separate theorem about the vacuum module.
\end{physics}
\begin{status}
Definition and unconditional lemma; axiom-clean.
\end{status}

%======================================================================
\section{Tensor products}\label{l04:sec:tensor}
%======================================================================

The file \leanpath{VOA/\allowbreak Core/\allowbreak TensorVOA.lean} constructs the tensor product of two
\emph{vertex superalgebras} in the sense of \lean{SuperVertexAlgebra}
(\leanpath{VOA/\allowbreak Core/\allowbreak SuperVertexAlgebra.lean}, Chapter~\ref{ch:L17}). We recall
what that structure is, in words: a space with a $\C$-valued weight grading (internal, finite-dimensional
pieces) and an independent $\Z/2$ parity grading, a linear state--field map $Y$, a vacuum of weight $0$
and even parity, and the axioms VA1 (weights add), VA1$'$ (parities add), VA2 ($Y(\lfourvac,z) = 1$),
VA3 (creation), VA4 (locality up to the Koszul sign $(-1)^{|u||v|}$) and VA5 (truncation). It has
\emph{no conformal vector}, no translation operator as a field, and no bounded-below condition.

\begin{definition}[Tensor product vertex operator]\label{l04:def:tenY}
For vertex superalgebras $V$ on $A$, $W$ on $B$:
$Y(u\otimes v,z)(u'\otimes v') = (-1)^{|v||u'|}\,Y(u,z)u'\otimes Y(v,z)v'$, with weight grading
$(A\otimes B)_h = \bigoplus_{h_1+h_2 = h}V_{h_1}\otimes W_{h_2}$, parity grading by addition, and
vacuum $\lfourvac\otimes\lfourvac$.
\end{definition}
\noindent File \leanpath{VOA/\allowbreak Core/\allowbreak TensorVOA.lean}.
\begin{leancode}
noncomputable def rawCoeff (V : SuperVertexAlgebra A) (W : SuperVertexAlgebra B) (a : A) (b : B)
    (s : ℤ) : A ⊗[ℂ] B →ₗ[ℂ] A ⊗[ℂ] B :=
  TensorProduct.lift (rawBil V W a b s)

noncomputable def tenY (V : SuperVertexAlgebra A) (W : SuperVertexAlgebra B) :
    A ⊗[ℂ] B →ₗ[ℂ] FormalDist (End ℂ (A ⊗[ℂ] B)) :=
  (2⁻¹ : ℂ) •
    ((postCompHom (1 + kappaTen V W 1)).comp (rawY V W)
      + (postCompHom (1 - kappaTen V W 1)).comp
          ((rawY V W).comp
            (TensorProduct.map (LinearMap.id : A →ₗ[ℂ] A) ((parGr W).parityInvol))))

theorem tenY_tmul_parity (V : SuperVertexAlgebra A) (W : SuperVertexAlgebra B) (a : A)
    {q : ZMod 2} {b : B} (hb : b ∈ W.parity q) (s : ℤ) (z : A ⊗[ℂ] B) :
    tenY V W (a ⊗ₜ[ℂ] b) s z = rawCoeff V W a b s (kappaTen V W q z) := by

def tenGrading (V : SuperVertexAlgebra A) (W : SuperVertexAlgebra B) (h : ℂ) :
    Submodule ℂ (A ⊗[ℂ] B) :=
  ⨆ h₁ : ℂ, tenSub (V.grading h₁) (W.grading (h - h₁))
\end{leancode}
\begin{dictionary}
\lean{A ⊗[ℂ] B} is the algebraic tensor product, \lean{a ⊗ₜ b} a pure tensor. \lean{rawCoeff a b s}
is the sign-free coefficient of $z^s$ in $Y(a,z)\otimes Y(b,z)$, i.e.\ $x\otimes y\mapsto\sum_m
(a$-coefficient of $z^m)x\otimes(b$-coefficient of $z^{s-m})y$ (\lean{rawCoeff_tmul}; finite by
truncation). \lean{kappaTen V W q} is $(-1)^{qF}\otimes1$, the parity operator to the power $q$ on the
first factor. The definition of \lean{tenY} is a device to keep $Y$ linear: for $b$ of parity $q$,
\lean{tenY_}\allowbreak\lean{tmul_}\allowbreak\lean{parity} shows it equals $\mathtt{rawCoeff}\circ((-1)^{qF}\otimes1)$, which on
$u'\otimes v'$ gives exactly $(-1)^{q|u'|}Y(a,z)u'\otimes Y(b,z)v'$, the Koszul formula of the
definition. \lean{tenSub P Q} is the span of $P\otimes Q$ inside $A\otimes B$, so \lean{tenGrading} is
the convolution grading. These agree with Kac (\emph{Vertex algebras for beginners}, §4.3) and FHL §3.
\end{dictionary}
\begin{physics}
Two decoupled chiral CFTs stacked: the states are products, the fields are products of fields, and
energies add. The sign is the usual fermionic statistics: moving a fermionic field of the second
theory past a fermionic state of the first costs a sign. For purely bosonic theories the sign is
absent.
\end{physics}
\begin{status}
Definitions; axiom-clean.
\end{status}

\begin{theorem}[The tensor product is a vertex superalgebra]\label{l04:thm:tensor}
If the parity involutions of $V$ and $W$ preserve the weight spaces, and the convolution weight
spaces are independent and finite-dimensional, then Definition~\ref{l04:def:tenY} is a vertex
superalgebra.
\end{theorem}
\noindent File \leanpath{VOA/\allowbreak Core/\allowbreak TensorVOA.lean}.
\begin{leancode}
def ParityCompatible (V : SuperVertexAlgebra A) : Prop :=
  ∀ (h : ℂ) (u : A), u ∈ V.grading h → (parGr V).parityInvol u ∈ V.grading h

structure TensorGradingHyp (V : SuperVertexAlgebra A) (W : SuperVertexAlgebra B) : Prop where
  iSupIndep : iSupIndep (tenGrading V W)
  finrank : ∀ h : ℂ, FiniteDimensional ℂ (tenGrading V W h)

noncomputable def tensorSVA (V : SuperVertexAlgebra A) (W : SuperVertexAlgebra B)
    (hV : ParityCompatible V) (hW : ParityCompatible W) (hg : TensorGradingHyp V W) :
    SuperVertexAlgebra (A ⊗[ℂ] B) where
  grading := tenGrading V W
\end{leancode}
\begin{dictionary}
\textbf{The hypotheses.} \lean{ParityCompatible}: the parity operator $(-1)^F$ maps each weight space to
itself (the algebra is bigraded). It holds in every example, but the \lean{SuperVertexAlgebra}
axioms do not force it, so it is a genuine hypothesis. \lean{TensorGradingHyp}: \lean{iSupIndep} says
the convolution weight spaces are independent (their sum is direct), \lean{finrank} that each is
finite-dimensional. The first is true over $\C$ by flatness (the file explains how it would be proved
but does not prove it). The second is \emph{false in general}: for complex weights, $h = h_1+h_2$ can
have infinitely many solutions with $V_{h_1}\ne0\ne W_{h_2}$; it holds when both gradings are bounded
below and discrete, which is the situation of VOAs of CFT type. (The file's docstring speaks of
``three'' linear-algebra facts; the structure has two fields, the third --- internality of the parity
grading --- being proved as \lean{tenParity_}\allowbreak\lean{isInternal}.)

\textbf{What is proved.} All of VA1--VA5 and VA1$'$: locality comes from Dong's lemma applied to the
mutually commuting families $a_{(n)}\otimes1$ and $(-1)^{qF}\otimes b_{(n)}$.
\end{dictionary}
\begin{physics}
Stacking two decoupled theories gives a consistent theory: locality of the product fields follows
from locality in each factor and the fact that fields of different factors (anti)commute. The
finiteness hypothesis is the requirement that each energy level of the combined theory be
finite-dimensional, which fails only for pathological spectra.
\end{physics}
\emph{Proof on paper.} Linearity is built in. The weight and parity rules follow from those of the
factors. The vacuum axioms follow from $Y(\lfourvac,z) = 1$ and creation in each factor. Locality: the
field of $a\otimes b$ is the product of $\ell_a(z) = Y(a,z)\otimes1$ and $r_b(z) = (-1)^{qF}\otimes
Y(b,z)$; these two families commute exactly up to the Koszul sign, so their normal-ordered product is
their plain product, and Dong's lemma (Chapter~\ref{ch:L03}, Theorem~\ref{l03:thm:dong}) gives mutual locality of all such
products. \qed
\begin{status}
Conditional on \lean{ParityCompatible} (twice) and \lean{TensorGradingHyp}; \texttt{\#print axioms}
on \lean{tensorSVA}: \lean{[propext, Classical.choice, Quot.sound]}. The two clauses of
\lean{TensorGradingHyp} are carried as open hypotheses downstream (in
\leanpath{VOA/\allowbreak Examples/\allowbreak DSReductionSl2.lean}, Chapter~\ref{ch:L19}).
\end{status}

\begin{warning}[title={Trap / dead end: There is no tensor product of vertex operator algebras in the library}]\label{l04:warn:tensor}
Despite its name, \leanpath{VOA/\allowbreak Core/\allowbreak TensorVOA.lean} does not construct a tensor
product of \lean{VertexOperatorAlgebra}s. It builds a \lean{SuperVertexAlgebra}: no conformal vector
$\omega\otimes\lfourvac+\lfourvac\otimes\omega$, no proof that central charges add, no CFT-type statement.
The source says so in its closing note (``The conformal layer \ldots\ is not built here''). A reader
who needs $c(V\otimes W) = c(V)+c(W)$ will not find it in this chapter's files.
\end{warning}

\begin{definition}[Tensor product of two Lie actions]\label{l04:def:tenlie}
For Lie algebra representations $R:L\to\End V$, $S:L\to\End W$: $(R\otimes S)(x) = R(x)\otimes1 +
1\otimes S(x)$.
\end{definition}
\noindent File \leanpath{VOA/\allowbreak Core/\allowbreak TensorLieAction.lean}.
\begin{leancode}
def tensorLieAction (R : L →ₗ⁅ℂ⁆ End ℂ V) (S : L →ₗ⁅ℂ⁆ End ℂ W) :
    L →ₗ⁅ℂ⁆ End ℂ (V ⊗[ℂ] W) where
  toFun x := (R x).rTensor W + (S x).lTensor V

theorem tensorLieAction_eigen (R : L →ₗ⁅ℂ⁆ End ℂ V) (S : L →ₗ⁅ℂ⁆ End ℂ W)
    (x : L) (v : V) (w : W) (a b : ℂ) (ha : R x v=a • v) (hb : S x w=b • w) :
    tensorLieAction R S x (v ⊗ₜ[ℂ] w)=(a+b) • (v ⊗ₜ[ℂ] w) := by

theorem complex_tmul_ne_zero {v : V} {w : W} (hv : v ≠ 0) (hw : w ≠ 0) :
    v ⊗ₜ[ℂ] w ≠ 0 := by
\end{leancode}
\begin{dictionary}
\lean{L →ₗ⁅ℂ⁆ End ℂ V} is a Lie algebra homomorphism, i.e.\ a representation. \lean{rTensor W}
and \lean{lTensor V} are $f\otimes1$ and $1\otimes f$. The file takes the representations as explicit
data (rather than type-class instances) so that several representations on one space can coexist.
Eigenvalues add, and pure tensors of nonzero vectors are nonzero. All standard.
\end{dictionary}
\begin{physics}
The total charge of a two-particle (or two-sector) state is the sum of the charges: the coproduct
$\Delta(x) = x\otimes1+1\otimes x$ of a Lie algebra symmetry. This is how $L_0$ and current zero modes act
on a stacked theory.
\end{physics}
\begin{status}
Unconditional; axiom-clean.
\end{status}

%======================================================================
\section{Weight operators and weight-homogeneous fields}\label{l04:sec:weightop}
%======================================================================

Two small files provide operator-first grading tools used by free-field constructions.

\begin{definition}[The weight operator of an internal grading]\label{l04:def:intweight}
Given an internal grading $A = \bigoplus_iG_i$ and scalars $w_i$, the operator acting by $w_i$ on $G_i$;
for a $\Z$-grading with $w_n = n$, the canonical ``$L_0$''.
\end{definition}
\noindent File \leanpath{VOA/\allowbreak Core/\allowbreak InternalGradingOperator.lean}.
\begin{leancode}
noncomputable def internalWeightOperatorZ (G : ℤ → Submodule ℂ A)
    (hG : DirectSum.IsInternal G) : End ℂ A :=
  internalWeightOperator G hG (fun n => (n : ℂ))

theorem internalWeightOperatorZ_comm_of_maps
    (G : ℤ → Submodule ℂ A) (hG : DirectSum.IsInternal G)
    (f : End ℂ A) (d : ℤ)
    (hf : ∀ n x, x ∈ G n → f x ∈ G (n + d)) :
    internalWeightOperatorZ G hG * f - f * internalWeightOperatorZ G hG =
      (d : ℂ) • f := by
\end{leancode}
\begin{dictionary}
\lean{DirectSum.IsInternal G} says $A$ is the internal direct sum of the $G_n$. The theorem: an operator
shifting degrees by $d$ satisfies $[W,f] = d f$. Standard.
\end{dictionary}
\begin{physics}
An operator that raises the energy by $d$ has commutator $d$ with the Hamiltonian; this is how one
reads off the weight of a mode, $[L_0,u_{(n)}] = (\operatorname{wt}u-n-1)u_{(n)}$.
\end{physics}
\begin{status}
Unconditional; axiom-clean.
\end{status}

\begin{definition}[Weight-commutator field]\label{l04:def:wcf}
A formal field $a(z) = \sum_ja_jz^j$ of $\End A$ is homogeneous of weight $D$ for $W\in\End A$ if
$[W,a_j] = (D+j)a_j$ for all $j$.
\end{definition}
\noindent File \leanpath{VOA/\allowbreak Core/\allowbreak WeightCommutatorField.lean}.
\begin{leancode}
def WeightCommutatorField (W : End ℂ A) (a : FormalDist (End ℂ A)) (D : ℤ) : Prop :=
  ∀ j : ℤ, W * a j - a j * W = ((D + j : ℤ) : ℂ) • a j

theorem nprod {W : End ℂ A} {a b : FormalDist (End ℂ A)}
    (haf : IsField a) (hbf : IsField b) {Da Db : ℤ}
    (ha : WeightCommutatorField W a Da) (hb : WeightCommutatorField W b Db) :
    WeightCommutatorField W (VOA.nprod haf hbf) (Da + Db) := by
\end{leancode}
\begin{dictionary}
\lean{a j} is the coefficient of $z^j$. For $a = Y(u,z)$ with $u$ of weight $D$, the coefficient of
$z^j$ is $u_{(-j-1)}$, of degree $D+j$: so the definition is ``$a$ is the field of a weight-$D$ state'' in
operator form, without any vertex algebra. \lean{VOA.nprod} is the normal-ordered product of two
fields (Chapter~\ref{ch:L03}); weights add. The file also proves closure under sums, scalars and
derivative (weight $D+1$).
\end{dictionary}
\begin{physics}
The normal-ordered product $:\!a(z)b(z)\!:$ of fields of dimensions $D_a$ and $D_b$ has dimension
$D_a+D_b$, and $\partial a$ has dimension $D_a+1$: the dimension counting of composite operators.
\end{physics}
\begin{status}
Unconditional; axiom-clean.
\end{status}

%======================================================================
\section{BRST cohomology}\label{l04:sec:brst}
%======================================================================

\begin{definition}[BRST complex]\label{l04:def:brst}
A BRST complex on a vertex superalgebra $V$ is a $\Z$-grading $V = \bigoplus_gV^{(g)}$ (ghost number)
together with $Q\in\End V$ that preserves weights, is odd, raises ghost number by one, satisfies
$Q^2 = 0$ and $Q(u_{(n)}v) = (Qu)_{(n)}v + ((-1)^Fu)_{(n)}(Qv)$ for all $u,v,n$.
\end{definition}
\noindent File \leanpath{VOA/\allowbreak Core/\allowbreak BRSTCohomology.lean}.
\begin{leancode}
structure BRSTComplex (V : SuperVertexAlgebra 𝒱) where
  ghost : ℤ → Submodule ℂ 𝒱
  ghost_isInternal : DirectSum.IsInternal ghost
  Q : Module.End ℂ 𝒱
  Q_wt : ∀ h : ℂ, ∀ v ∈ V.grading h, Q v ∈ V.grading h
  Q_odd : ∀ p : ZMod 2, ∀ v ∈ V.parity p, Q v ∈ V.parity (p + 1)
  Q_ghost : ∀ g : ℤ, ∀ v ∈ ghost g, Q v ∈ ghost (g + 1)
  Q_sq : Q * Q = 0
  Q_deriv : ∀ (u v : 𝒱) (n : ℤ),
    Q (mode (V.Y u) n v) = mode (V.Y (Q u)) n v + mode (V.Y (sgn V u)) n (Q v)
\end{leancode}
\begin{dictionary}
\lean{sgn V} is the parity operator $(-1)^F$. For $u$ of parity $p$ the derivation rule is the usual
$Q(u_{(n)}v) = (Qu)_{(n)}v + (-1)^pu_{(n)}(Qv)$ (\lean{Q_deriv_homog}); writing the sign as an operator
makes the rule bilinear for all $u$. The definition is abstract: it does not ask that $Q$ be the zero
mode of a BRST current $J_{\mathrm{BRST}}$, nor that the ghost number be additive under products
(that is the separate predicate \lean{GhostAdditive}). Every BRST differential of the literature
(Feigin--Frenkel, Kac--Roan--Wakimoto) has these properties, so the notion is at least as general as
needed. \textbf{No concrete instance is constructed in this chapter's files}; the Drinfeld--Sokolov
file of Chapter~\ref{ch:L19} carries one as a hypothesis.
\end{dictionary}
\begin{physics}
The BRST charge of a gauge-fixed theory: a nilpotent fermionic symmetry raising the ghost number. Its
cohomology is the space of physical states and operators --- gauge invariant modulo pure gauge. In
quantum Hamiltonian (Drinfeld--Sokolov) reduction, $Q$ imposes the constraint $J^{e}(z) = 1$ on an affine
current, and the cohomology is the $W$-algebra.
\end{physics}
\begin{status}
Definition.
\end{status}

\begin{theorem}[BRST cohomology is a vertex superalgebra]\label{l04:thm:brst}
$\lfourH(V,Q) := \ker Q/\operatorname{im}Q$ carries a vertex superalgebra structure with
$Y([u],z)[v] = [Y(u,z)v]$, vacuum $[\lfourvac]$ and the induced weight and parity gradings; the ghost
grading descends. If the ghost number is additive, $\lfourH^0$ is closed under all products. If $Q = 0$,
$\lfourH(V,0)\cong V$.
\end{theorem}
\noindent File \leanpath{VOA/\allowbreak Core/\allowbreak BRSTCohomology.lean}.
\begin{leancode}
abbrev Coh : Type u := C.kerQ ⧸ C.imQ'

noncomputable def cohomology : SuperVertexAlgebra C.Coh where
  grading := C.gradingH
  isInternal := C.gradingH_isInternal
  finrank := C.gradingH_finiteDimensional
  parity := C.parityH
  parity_isInternal := C.parityH_isInternal
  Y := C.Ycohom
  vacuum := C.prj C.vacK

def GhostAdditive : Prop :=
  ∀ (g g' : ℤ), ∀ u ∈ C.ghost g, ∀ v ∈ C.ghost g', ∀ n : ℤ,
    mode (V.Y u) n v ∈ C.ghost (g + g')

theorem H0_mode (hadd : C.GhostAdditive) {ubar vbar : C.Coh}
    (hu : ubar ∈ C.H0) (hv : vbar ∈ C.H0) (n : ℤ) :
    mode (C.Ycohom ubar) n vbar ∈ C.H0 := by

noncomputable def cohEquivOfQZero (hQ : C.Q = 0) : 𝒱 ≃ₗ[ℂ] C.Coh :=
  LinearEquiv.ofBijective (C.toCoh hQ) ⟨C.toCoh_injective hQ, C.toCoh_surjective hQ⟩
\end{leancode}
\begin{dictionary}
\lean{C.kerQ} is $\ker Q$; \lean{C.imQ'} is $\operatorname{im}Q$ viewed inside $\ker Q$ (it is contained
there since $Q^2 = 0$); \lean{Coh} is the quotient space. \lean{Ycohom} is the induced state--field
map; well-definedness rests on $\ker Q$ being closed under products and $\operatorname{im}Q$ being a
two-sided ideal \emph{inside} $\ker Q$ (\lean{mode_}\allowbreak\lean{mem_}\allowbreak\lean{imQ_}\allowbreak\lean{left}, \lean{mode_}\allowbreak\lean{mem_}\allowbreak\lean{imQ_}\allowbreak\lean{right}).
\lean{gradingH h} is the image of $V_h\cap\ker Q$. \lean{H0} is the ghost-number-zero part. All fifteen
fields of \lean{SuperVertexAlgebra} are proved; nothing is assumed beyond the \lean{BRSTComplex}.
\textbf{Nontriviality is not part of the theorem}: if $\lfourvac\in\operatorname{im}Q$ the cohomology is
zero, which is a valid vertex superalgebra. \lean{cohEquivOfQZero} (with \lean{toCoh_vacuum},
\lean{toCoh_mode}) is the non-vacuity sanity check that for $Q = 0$ nothing is lost.
\end{dictionary}
\begin{physics}
Physical operators form a closed operator algebra: the OPE of two BRST-closed operators is closed, and
exact operators decouple (their OPE with closed ones is exact). The ghost-number-zero part is the
algebra of physical fields --- for Drinfeld--Sokolov reduction, the $W$-algebra $W^k(\mathfrak g,f)$.
\end{physics}
\emph{Proof on paper.} Let $a,x\in\ker Q$. Then $Q(a_{(n)}x) = 0$ by the derivation rule. If $a = Qb$
and $Qx = 0$, $(Qb)_{(n)}x = Q(b_{(n)}x)$, exact. If $Qa = 0$ and $x = Qb$, apply the rule to $(-1)^Fa$
(which is still closed, since $Q(-1)^F = -(-1)^FQ$): $Q(((-1)^Fa)_{(n)}b) = a_{(n)}Qb$ up to the sign,
exact. So $\operatorname{im}Q$ is a two-sided ideal of the vertex algebra $\ker Q$ and the axioms
descend. The three gradings descend because $Q$ shifts each by a bijection of the index set. \qed
\begin{status}
\lean{cohomology} unconditional (given a \lean{BRSTComplex}); \lean{H0_mode} conditional on
\lean{GhostAdditive}. \texttt{\#print axioms} on \lean{cohomology}, \lean{H0_mode},
\lean{cohEquivOfQZero}: \lean{[propext, Classical.choice, Quot.sound]}.
\end{status}

\begin{theorem}[Nontrivial vacuum; conformal structure descends]\label{l04:thm:brstconf}
(a) If $\lfourvac$ has ghost number $0$, $V^{(-1)} = 0$ and $\lfourvac\ne0$, then $\lfourvac\notin
\operatorname{im}Q$, so the cohomology vacuum is nonzero. (b) If $V$ is a super near-VOA and
$Q\omega = 0$, then $\lfourH(V,Q)$ is a super near-VOA with conformal vector $[\omega]$ and the same
central charge.
\end{theorem}
\noindent File \leanpath{VOA/\allowbreak Core/\allowbreak BRSTCohomology.lean}.
\begin{leancode}
theorem vacuum_not_mem_imQ (h0 : V.vacuum ∈ C.ghost 0) (hneg : C.ghost (-1) = ⊥)
    (hne : V.vacuum ≠ 0) : V.vacuum ∉ C.imQ := by

noncomputable def cohomologyNearVOA : SuperNearVOA C.Coh where
  toSuperVertexAlgebra := C.cohomology
  conformalVector := C.prj (confK C hω)
  conformal_mem := Submodule.mem_map_of_mem W.conformal_mem
  conformal_parity := Submodule.mem_map_of_mem W.conformal_parity
  centralCharge := W.centralCharge
\end{leancode}
\begin{dictionary}
(a) \lean{hneg : C.ghost (-1) = ⊥} says there are no states of ghost number $-1$. It is a
\emph{sufficient} condition: the ghost-number-zero part of $\operatorname{im}Q$ is $Q(V^{(-1)})$. A
referee should be aware that in the standard Kac--Roan--Wakimoto complex the charged ghost system has
states of negative ghost number, so this particular criterion is not the one that would establish a
nonzero vacuum there (in that case one uses the vanishing of the cohomology outside degree $0$, which is
not in the library). (b) \lean{SuperNearVOA} (\leanpath{VOA/\allowbreak Core/\allowbreak SuperNearVOA.lean},
Chapter~\ref{ch:L17}) is a \lean{SuperVertexAlgebra} with a conformal vector $\omega$ of weight $2$, even,
whose modes satisfy the Virasoro relations, with $L_0$ acting by the weight and $L_{-1}$ as translation;
\lean{hω} is $Q\omega = 0$. The omitted fields (Virasoro relations, $L_0$-grading, translation) are
transported from $\ker Q$.
\end{dictionary}
\begin{physics}
(a) The vacuum is a physical state, not pure gauge. (b) If the stress tensor is BRST invariant, the
physical theory is conformal, with the total central charge (matter plus ghosts) of the complex. In
Drinfeld--Sokolov reduction one must first improve the Sugawara stress tensor so that it is BRST
closed; this is the role of \lean{hω}.
\end{physics}
\begin{status}
(a) conditional on \lean{h0}, \lean{hneg}, \lean{hne}; (b) conditional on \lean{hω}. Both axiom-clean.
\end{status}

%======================================================================
\section{Invariant bilinear forms}\label{l04:sec:forms}
%======================================================================

\subsection{Motivation: the two-point function}

In a unitary CFT the space of states carries an inner product for which $L_n^\dagger = L_{-n}$ and,
more generally, the adjoint of a field is obtained by the inversion $z\mapsto1/z$: for a quasi-primary
field of weight $h$, $\phi(z)^\dagger = \bar z^{-2h}\phi^\dagger(1/\bar z)$. Removing complex conjugation,
one obtains a \emph{bilinear} pairing on the chiral algebra, the BPZ (Belavin--Polyakov--Zamolodchikov)
pairing; restricted to quasi-primaries it is the Zamolodchikov metric, the matrix of two-point
functions $\langle\phi_a(z)\phi_b(w)\rangle = G_{ab}(z-w)^{-2h}$. Frenkel--Huang--Lepowsky (FHL) turned this
into the following algebraic definition. The map $w\mapsto1/w$ has derivative $-w^{-2}$; on a state $u$
the corresponding operator is $e^{zL_1}(-z^{-2})^{L_0}$, where the factor $(-z^{-2})^{L_0}$ is the
Jacobian for a quasi-primary and $e^{zL_1}$ corrects for the failure of $u$ to be quasi-primary.

\begin{definition}[Invariant form, the library's two layers]\label{l04:def:invform}
An \emph{invariant form} in the library's sense is a bilinear form $(\cdot|\cdot)$ on $\lfourV$ with
$(L_na|b) = (a|L_{-n}b)$ for all $n$. It is \emph{normalised} if $(\lfourvac|\lfourvac) = 1$, and
\emph{Gannon-normalised} if $(\lfourvac|\lfourvac) = -1$. Its \emph{left radical} is $\{a : (a|b) = 0\
\forall b\}$; it is \emph{symmetric} if $(a|b) = (b|a)$.
\end{definition}
\noindent File \leanpath{VOA/\allowbreak Core/\allowbreak BilinearForm.lean}.
\begin{leancode}
structure InvariantForm (W : VertexOperatorAlgebra 𝒱) where
  form : LinearMap.BilinForm ℂ 𝒱
  virasoro_adj : ∀ (n : ℤ) (a b : 𝒱), form (W.L n a) b = form a (W.L (-n) b)

def IsSymmetric (B : InvariantForm W) : Prop := ∀ a b : 𝒱, B.form a b = B.form b a

def NormalizedAt (B : InvariantForm W) (κ : ℂ) : Prop :=
  B.form W.vacuum W.vacuum = κ

def Normalized (B : InvariantForm W) : Prop := B.NormalizedAt 1

def GannonNormalized (B : InvariantForm W) : Prop := B.NormalizedAt (-1)

def leftRadical (B : InvariantForm W) : Submodule ℂ 𝒱 := LinearMap.ker B.form
\end{leancode}
\begin{dictionary}
\lean{LinearMap.BilinForm ℂ 𝒱} is the type of bilinear forms $\lfourV\times\lfourV\to\C$, curried as
$\lfourV\to(\lfourV\to\C)$; its kernel is the set of $a$ with $(a|\cdot) = 0$. The full name of the
structure is \lean{VOA.}\allowbreak\lean{VertexOperatorAlgebra.}\allowbreak\lean{InvariantForm} (distinct from the Lie-algebra
\lean{VOA.}\allowbreak\lean{InvariantForm} of \leanpath{VOA/\allowbreak Foundations/\allowbreak AffineLie.lean}).

\textbf{Important: \lean{InvariantForm} alone is weaker than FHL invariance.} It records only the
$u = \omega$ slice of FHL invariance (Virasoro adjointness). The full FHL condition is the separate
predicate \lean{IsFullyInvariant} (Definition~\ref{l04:def:fullinv}), and every substantive theorem
assumes it. A referee must check, for each statement mentioning an \lean{InvariantForm}, whether
\lean{IsFullyInvariant} is also assumed. Section~\ref{l04:sec:deadends} shows that the gap between the
two is real.

\textbf{Normalisation.} The public convention is $(\lfourvac|\lfourvac) = +1$; Gannon's book uses $-1$
(5.1.12a), and the library keeps it as \lean{GannonNormalized}.
\end{dictionary}
\begin{physics}
The bilinear BPZ pairing: $(a|b)$ is the coefficient of the two-point function of the fields of $a$ and
$b$ with $a$ inserted at infinity. Virasoro adjointness $L_n^\dagger = L_{-n}$ is the statement that the
stress tensor is self-conjugate.
\end{physics}
\begin{status}
Definitions.
\end{status}

\begin{definition}[Full (FHL) invariance]\label{l04:def:fullinv}
$(Y(u,z)a|b) = (a|Y(e^{zL_1}(-z^{-2})^{L_0}u,z^{-1})b)$ for all $u,a,b$; in modes, for $u\in\lfourV_{\mathrm{wt}}$,
\begin{equation}\label{l04:eq:fhl}
(u_{(n)}a|b) = (-1)^{\mathrm{wt}}\sum_{j=0}^{\mathrm{wt}}\frac1{j!}\big(a\,\big|\,(L_1^ju)_{(2\mathrm{wt}-n-2-j)}b\big).
\end{equation}
\end{definition}
\noindent File \leanpath{VOA/\allowbreak Core/\allowbreak BilinearForm.lean}.
\begin{leancode}
abbrev IsFullyInvariant (B : InvariantForm W) : Prop :=
  ∀ (wt : ℤ) (u : 𝒱), u ∈ W.grading wt → ∀ (n : ℤ) (a b : 𝒱),
    B.form (W.smul u n a) b
      = (-1 : ℂ) ^ wt • ∑ j ∈ Finset.range (wt.toNat + 1), ((j.factorial : ℂ)⁻¹) •
          B.form a (W.smul (((W.L 1) ^ j) u) (2 * wt - n - 2 - (j : ℤ)) b)
\end{leancode}
\begin{dictionary}
\textbf{Derivation of the mode form.} $(-z^{-2})^{L_0}u = (-1)^{\mathrm{wt}}z^{-2\mathrm{wt}}u$;
$e^{zL_1}$ contributes $\sum_jz^jL_1^j/j!$; $Y(v,z^{-1}) = \sum_kv_{(k)}z^{k+1}$. The total power of $z$ is
$-2\mathrm{wt}+j+k+1$, which equals $-n-1$ exactly when $k = 2\mathrm{wt}-n-2-j$. The sum stops at
$j = \mathrm{wt}$ because $L_1^ju\in\lfourV_{\mathrm{wt}-j} = 0$ for $j>\mathrm{wt}$ (CFT type;
\lean{L_}\allowbreak\lean{one_}\allowbreak\lean{pow_}\allowbreak\lean{eq_}\allowbreak\lean{zero}). So the Lean term is exactly \eqref{l04:eq:fhl}, with \lean{(-1 : ℂ) ^ wt}
$= (-1)^{\mathrm{wt}}$ (a \lean{zpow}, harmless since only $\mathrm{wt}\ge0$ has $u\ne0$),
\lean{Finset.range (wt.toNat + 1)} $= \{0,\dots,\mathrm{wt}\}$ and \lean{((W.L 1) ^ j) u} $= L_1^ju$.

\textbf{Quantifiers.} For each weight and each homogeneous $u$; general $u$ follows by linearity. For
$\mathrm{wt}<0$ the hypothesis forces $u = 0$ and both sides vanish.

\textbf{Same as FHL.} This is FHL (5.2.1) / Gannon (5.1.11a) verbatim. \textbf{A documentation slip:}
the docstring above \lean{IsFullyInvariant} writes the right side as ``$((L_1^ju)_{(\cdot)}a|b)$''; the
Lean term has \lean{B.form a (W.smul … b)}, i.e.\ $(a|(L_1^ju)_{(\cdot)}b)$, which is the correct formula.
The statement, not the comment, is what is used. The \lean{virasoro_adj} field is the slice $u = \omega$
(where $L_1\omega = 0$, \lean{L1_}\allowbreak\lean{conformal_}\allowbreak\lean{zero}).
\end{dictionary}
\begin{physics}
The adjoint of the field of $u$ is the field of the BPZ conjugate state at the inverted point. For a
quasi-primary current $J$ of weight $1$ it reads $(J_na|b) = -(a|J_{-n}b)$: the bilinear form makes $J_n$
anti-adjoint to $J_{-n}$. (The positive-definite Hermitian inner product of a unitary theory is obtained
by composing with an antilinear involution that sends $J$ to $-\bar J$; this is why signs like
$(h|h) = -k$ below are not a contradiction with unitarity.)
\end{physics}
\begin{status}
Definition.
\end{status}

\subsection{Elementary consequences}

\begin{theorem}[Orthogonality, radical, uniqueness, nondegeneracy, evaluation on $\lfourV_1$]\label{l04:thm:formbasics}
Let $B$ be an invariant form. (i) $(\lfourV_m|\lfourV_n) = 0$ for $m\ne n$. Now let $B$ be fully invariant.
(ii) Its left radical is an ideal. (iii) If $B_1,B_2$ are fully invariant and normalised, $B_1 = B_2$.
(iv) If $\lfourV$ is simple and $B$ is fully invariant and normalised, its radical is $0$. (v) If $B$ is
fully invariant and normalised, $u_{(1)}v = -(u|v)\lfourvac$ for $u,v\in\lfourV_1$.
\end{theorem}
\noindent File \leanpath{VOA/\allowbreak Core/\allowbreak BilinearForm.lean}.
\begin{leancode}
theorem orthogonal_of_ne (B : InvariantForm W) {m n : ℤ} (hmn : m ≠ n) {a b : 𝒱}
    (ha : a ∈ W.grading m) (hb : b ∈ W.grading n) : B.form a b = 0 := by

theorem radical_isIdeal {B : InvariantForm W} (hB : B.IsFullyInvariant) :
    W.IsIdeal (leftRadical B) := by

theorem invariantForm_unique {B₁ B₂ : InvariantForm W}
    (h₁ : B₁.IsFullyInvariant) (h₂ : B₂.IsFullyInvariant)
    (hn₁ : B₁.Normalized) (hn₂ : B₂.Normalized) : B₁.form = B₂.form := by

theorem invariantForm_nondegenerate {B : InvariantForm W} (hsimple : W.IsSimple)
    (hB : B.IsFullyInvariant) (hn : B.Normalized) : leftRadical B = ⊥ := by

theorem eval_oneMode_grading1 {B : InvariantForm W} (hB : B.IsFullyInvariant)
    (hn : B.Normalized) {u v : 𝒱} (hu : u ∈ W.grading 1) (hv : v ∈ W.grading 1) :
    W.smul u 1 v = -B.form u v • W.vacuum := by
\end{leancode}
\begin{dictionary}
All are the textbook statements (Gannon 5.1.11b, 5.1.12b), with these precisions. (i) needs only
\lean{virasoro_adj} at $n = 0$. (iii) is uniqueness \emph{with normalisation}, i.e.\ equality, not
``up to a scalar''; it needs no simplicity and no $L_1\lfourV_1 = 0$ --- if two such forms exist they
agree. (iv) derives nondegeneracy from simplicity; the converse is Theorem~\ref{l04:thm:lirad}. (v) has a
minus sign because of the public normalisation $(\lfourvac|\lfourvac) = 1$ and the $(-1)^{\mathrm{wt}}$ of
FHL; in Gannon's normalisation it reads $u_{(1)}v = (u|v)\lfourvac$.
\end{dictionary}
\begin{physics}
(i) Energy eigenstates of different energy are orthogonal. (ii) Null states (zero overlap with
everything) form a decoupled sector. (iii) Two-point functions are fixed by the OPE and the vacuum
normalisation. (iv) An irreducible theory has a nondegenerate Zamolodchikov metric. (v) The double pole of
two currents is their two-point function: for a free boson $h(z)h(w)\sim k/(z-w)^2$ one has $h_{(1)}h =
k\lfourvac$ and so $(h|h) = -k$.
\end{physics}
\emph{Proof on paper.} (i) $m(a|b) = (L_0a|b) = (a|L_0b) = n(a|b)$. (ii) If $(a|\cdot) = 0$, then
$(u_{(n)}a|b)$ is by \eqref{l04:eq:fhl} a combination of $(a|\cdots) = 0$. (iii) $D = B_1-B_2$ is fully
invariant with $(\lfourvac|\lfourvac)_D = 0$; by (i) and $\lfourV_0 = \C\lfourvac$, $\lfourvac\in\lfourrad D$; the
radical is an ideal containing $\lfourvac$, hence $\lfourV$, so $D = 0$. (iv) The radical is an ideal; it is
not $\lfourV$ since $(\lfourvac|\lfourvac) = 1$. (v) Write $u = u_{(-1)}\lfourvac$ and apply \eqref{l04:eq:fhl} with
$\mathrm{wt} = 1$, $n = -1$: $(u|v) = -\big((\lfourvac|u_{(1)}v) + (\lfourvac|(L_1u)_{(0)}v)\big)$; the second
term vanishes since $L_1u\in\C\lfourvac$ and $\lfourvac_{(0)} = 0$; and $u_{(1)}v = \mu\lfourvac$ gives
$(\lfourvac|u_{(1)}v) = \mu$. \qed
\begin{status}
Unconditional given the displayed hypotheses; \texttt{\#print axioms} on \lean{radical_isIdeal} and
\lean{invariantForm_}\allowbreak\lean{unique}: \lean{[propext, Classical.choice, Quot.sound]}.
\end{status}

\begin{theorem}[Symmetry of the $(1)$-pairing on $\lfourV_1$]\label{l04:thm:v1symm}
For $u,v\in\lfourV_1$: $u_{(1)}v = v_{(1)}u$ and $u_{(0)}v+v_{(0)}u = 0$.
\end{theorem}
\noindent File \leanpath{VOA/\allowbreak Core/\allowbreak InvariantFormSymmetry.lean}.
\begin{leancode}
theorem oneMode_symm_grading1 (W : VertexOperatorAlgebra 𝒱) {u v : 𝒱}
    (hu : u ∈ W.grading 1) (hv : v ∈ W.grading 1) :
    W.smul u 1 v = W.smul v 1 u := by

theorem zeroMode_antisymm_grading1 (W : VertexOperatorAlgebra 𝒱) {u v : 𝒱}
    (hu : u ∈ W.grading 1) (hv : v ∈ W.grading 1) :
    W.smul u 0 v + W.smul v 0 u = 0 := by
\end{leancode}
\begin{dictionary}
Unconditional consequences of skew-symmetry; with Theorem~\ref{l04:thm:formbasics}(v) they make every
fully invariant normalised form symmetric \emph{on $\lfourV_1$}. Symmetry of the whole form is a theorem
of FHL (Prop.~5.3.6) and Li, but it is \textbf{not proved in the library} for any form, including the
canonical one; the file's closing comment states the gap. Downstream statements never assume symmetry
without saying so.
\end{dictionary}
\begin{physics}
The level matrix $k^{ab}$ in $J^a(z)J^b(w)\sim k^{ab}/(z-w)^2$ is symmetric, because the double pole of a
bosonic OPE is symmetric under exchange.
\end{physics}
\emph{Proof on paper.} Skew-symmetry at $n = 1$ and $n = 0$; the correction terms $T^j(v_{(1+j)}u)$ lie in
$T^j\lfourV_{-j}$, which is $0$ for $j\ge1$ (or $T\C\lfourvac = 0$). \qed
\begin{status}
Unconditional; axiom-clean.
\end{status}

%======================================================================
\section{The canonical invariant form (FHL, Li)}\label{l04:sec:canonical}
%======================================================================

\subsection{The construction}

Li's idea: if a normalised fully invariant form exists, then putting $a = \lfourvac$ in
\eqref{l04:eq:fhl} shows $(x|b)$ is forced to be the vacuum coefficient of a certain mode product. One
therefore \emph{defines} the form by that formula and proves invariance.

\begin{definition}[Vacuum coordinate and the canonical pairing]\label{l04:def:canbilin}
Let $\lfourvc:\lfourV\to\C$ be the coefficient of $\lfourvac$ in the weight-zero component. For homogeneous
$a\in\lfourV_p$ put
\begin{equation*}
\mathrm{adjPair}_p(a,b) = (-1)^p\sum_{j=0}^{p}\frac1{j!}\,\lfourvc\big((L_1^ja)_{(2p-1-j)}b\big),
\end{equation*}
and extend to $\mathrm{canBilin}(a,b) = -\sum_p\mathrm{adjPair}_p(a_p,b)$ over the homogeneous components
$a_p$ of $a$.
\end{definition}
\noindent File \leanpath{VOA/\allowbreak Core/\allowbreak InvariantFormCanonical.lean}.
\begin{leancode}
noncomputable def vacCoord (W : VertexOperatorAlgebra V) (hne : W.vacuum ≠ 0) : V →ₗ[ℂ] ℂ :=
  ((LinearEquiv.toSpanNonzeroSingleton ℂ V W.vacuum hne).symm.toLinearMap).comp
    ((grade0ToSpan W).toLinearMap.comp (grade0 W))

noncomputable def adjPair (W : VertexOperatorAlgebra V) (hne : W.vacuum ≠ 0) (wt : ℤ)
    (a b : V) : ℂ :=
  (-1 : ℂ) ^ wt • ∑ j ∈ Finset.range (wt.toNat + 1), ((j.factorial : ℂ)⁻¹) •
    vacCoord W hne (W.smul (((W.L 1) ^ j) a) (2 * wt - 1 - (j : ℤ)) b)

noncomputable def canBilin (W : VertexOperatorAlgebra V) (hne : W.vacuum ≠ 0) :
    V →ₗ[ℂ] (V →ₗ[ℂ] ℂ) :=
  letI : DirectSum.Decomposition W.grading := W.isInternal.chooseDecomposition
  (DirectSum.toModule ℂ ℤ (V →ₗ[ℂ] ℂ) (bSlotOfGrade W hne)).comp
    (DirectSum.decomposeLinearEquiv W.grading).toLinearMap
\end{leancode}
\begin{dictionary}
\lean{grade0 W} is the projection to $\lfourV_0$, \lean{grade0ToSpan} rewrites $\lfourV_0 = \C\lfourvac$, and
\lean{toSpanNonzeroSingleton … hne} identifies $\C\lfourvac$ with $\C$ via $c\lfourvac\mapsto c$, which needs
$\lfourvac\ne0$ (\lean{hne}). So \lean{vacCoord} is $\lfourvc$. The mode index $2p-1-j$ is
$2\mathrm{wt}-n-2-j$ at $n = -1$: $\mathrm{adjPair}_p(a,b)$ is the value of the right side of
\eqref{l04:eq:fhl} at $u = a$, $n = -1$, read against $\lfourvac$, i.e.\ ``$(a_{(-1)}\lfourvac|b)$ if the form
existed''. \lean{canBilin} glues these over the decomposition of $a$ (\lean{bSlotOfGrade} is the
weight-$p$ piece $-\mathrm{adjPair}_p$, \lean{canBilin_homog}). The leading minus makes
$\mathrm{canBilin}(\lfourvac,\lfourvac) = -1$ (\lean{canBilin_vacuum}), Gannon's normalisation; the public
form is its negative.
\end{dictionary}
\begin{physics}
$(a|b)$ is read off as the vacuum component of the BPZ-conjugate of $a$ acting on $b$: ``$\langle0|
a^\dagger b\rangle$''. The chiral two-point function is computed by bringing one operator to infinity and
taking the vacuum expectation value.
\end{physics}
\begin{status}
Definitions; axiom-clean.
\end{status}

\begin{theorem}[Why $L_1\lfourV_1 = 0$ is needed]\label{l04:thm:virslice}
For every $b$: $\mathrm{canBilin}(L_{-1}\lfourvac,b) = 0$ and
$\mathrm{canBilin}(\lfourvac,L_1b) = -\lfourvc(L_1b)$. Hence the $n = -1$ Virasoro adjointness of canBilin at
$a = \lfourvac$ holds if and only if $\lfourvc(L_1b) = 0$ for all $b$, i.e.\ iff $L_1\lfourV_1 = 0$.
\end{theorem}
\noindent File \leanpath{VOA/\allowbreak Core/\allowbreak InvariantFormCanonicalFHL.lean}.
\begin{leancode}
theorem virSlice_vacuum (hne : W.vacuum ≠ 0) (b : V) :
    canBilin W hne (W.L (-1) W.vacuum) b = 0
      ∧ canBilin W hne W.vacuum (W.L 1 b) = - vacCoord W hne (W.L 1 b) := by
\end{leancode}
\begin{dictionary}
The Lean statement is the two displayed equalities; the ``iff'' is the one-line consequence (since
$L_1b$ has weight-zero component $L_1b_1$, the condition is $L_1\lfourV_1 = 0$, proved as
\lean{vacCoord_}\allowbreak\lean{L1_}\allowbreak\lean{eq_}\allowbreak\lean{zero} in \leanpath{VOA/\allowbreak Core/\allowbreak InvariantFormCanonicalVir.lean}). It is the
library's certificate that the hypothesis \lean{hqp} of Theorem~\ref{l04:thm:canonical} cannot be
dropped for this construction. More is true on paper: \emph{any} normalised form satisfying the
$n = 1$ Virasoro adjointness has $(\lfourvac|L_1v) = (L_{-1}\lfourvac|v) = 0$ for $v\in\lfourV_1$, and
$L_1v\in\C\lfourvac$, so $L_1\lfourV_1 = 0$ is necessary for the existence of a normalised invariant form at
all. That converse is not stated in the library.
\end{dictionary}
\begin{physics}
$L_1\lfourV_1 = 0$ says the currents are quasi-primary. When it fails --- e.g.\ the free boson with
background charge, $T = \tfrac12(\partial\phi)^2+\lambda\partial^2\phi$, where $L_1(\partial\phi)\propto\lambda\lfourvac$ ---
the vacuum is not self-conjugate: the Coulomb-gas two-point function $\langle V_\alpha V_\beta\rangle$ is
nonzero only for $\alpha+\beta = 2\alpha_0$, and the BPZ pairing pairs the theory with a twisted version
of itself. No normalised invariant form exists then.
\end{physics}
\begin{status}
Unconditional; axiom-clean.
\end{status}

\subsection{The proof of invariance}

The proof that canBilin is fully invariant under $L_1\lfourV_1 = 0$ occupies most of the
\texttt{InvariantForm*} files. Mathematically it has three steps.

\paragraph{Step 1: reduction to the diagonal.} By weight orthogonality and bilinearity it suffices to
prove \eqref{l04:eq:fhl} for homogeneous $u\in\lfourV_{\mathrm{wt}}$, $a\in\lfourV_p$, $b\in\lfourV_q$ at the one mode
$n = p+\mathrm{wt}-q-1$ for which $u_{(n)}a\in\lfourV_q$ (all other modes give $0 = 0$). This is the
predicate \lean{CanDiagonal}.

\paragraph{Step 2: the quasi-primary decomposition.} Under $L_1\lfourV_1 = 0$, every $a\in\lfourV_p$, $p\ge1$, is
$a_0+L_{-1}c$ with $L_1a_0 = 0$ and $c\in\lfourV_{p-1}$ (Theorem~\ref{l04:thm:qp}). Both sides of
\eqref{l04:eq:fhl} obey the same $L_{-1}$-covariance recursion in $a$ (from $[L_{-1},u_{(n)}] =
-nu_{(n-1)}$ and the Virasoro adjointness $\mathrm{canBilin}(L_{-1}a,b) = \mathrm{canBilin}(a,L_1b)$, the
latter proved under \lean{hqp} as \lean{canBilin_}\allowbreak\lean{Lneg1_}\allowbreak\lean{adj}), so it suffices to treat quasi-primary $a$,
and similarly quasi-primary $u$.

\paragraph{Step 3: the quasi-primary matching identity (Li, Prop.~2.6).} For quasi-primary $u\in\lfourV_w$
and $a\in\lfourV_p$, $\mathrm{canBilin}(u_{(n)}a,b) = (-1)^w\mathrm{canBilin}(a,u_{(2w-n-2)}b)$, proved by a
direct computation with the vacuum coordinate: skew-symmetry shows $\lfourvc(u_{(n)}v) = \pm\lfourvc(v_{(n)}u)$
(the $T$-tower dies under $\lfourvc$), and a raising identity $(2p-M-2)\lfourvc(a_{(M+1)}X) =
-\lfourvc(a_{(M)}L_1X)$ for quasi-primary $a$ kills low modes.

\begin{theorem}[Quasi-primary decomposition]\label{l04:thm:qp}
If $L_1\lfourV_1 = 0$ and $p\ge1$, every $a\in\lfourV_p$ can be written $a = a_0+L_{-1}c$ with $a_0\in\lfourV_p$,
$L_1a_0 = 0$, $c\in\lfourV_{p-1}$.
\end{theorem}
\noindent File \leanpath{VOA/\allowbreak Core/\allowbreak QuasiPrimaryDecomp.lean}.
\begin{leancode}
theorem qpDecomp (W : VertexOperatorAlgebra V)
    (hqp : ∀ v, v ∈ W.grading 1 → W.L 1 v = 0)
    {p : ℤ} (hp : 1 ≤ p) {a : V} (ha : a ∈ W.grading p) :
    ∃ a₀ c : V, a₀ ∈ W.grading p ∧ W.L 1 a₀ = 0 ∧ c ∈ W.grading (p - 1) ∧ a = a₀ + W.L (-1) c := by
\end{leancode}
\begin{dictionary}
The statement asserts existence of a decomposition (the sum $\lfourQP_p+L_{-1}\lfourV_{p-1}$ is everything),
not directness; directness also holds for $p\ge2$ but is not needed. The hypothesis \lean{hqp} is used
exactly at $p = 1$, where $L_{-1}\lfourV_0 = L_{-1}\C\lfourvac = 0$ forces $\lfourV_1 = \lfourQP_1$. The docstring's
``$\lfourV_p = \mathcal{P}\lfourV_p\oplus L_{-1}\lfourV_{p-1}$'' should be read with $\mathcal P$ meaning
quasi-primary, not primary.
\end{dictionary}
\begin{physics}
Every state is a quasi-primary or a derivative: under the global conformal group $SL_2$ (generated by
$L_{-1},L_0,L_1$) each energy level decomposes into lowest-weight vectors (quasi-primaries) and their
$L_{-1}$-descendants. Quasi-primary fields are those that transform simply under global conformal
transformations, $\phi(z)\mapsto(cz+d)^{-2h}\phi(\frac{az+b}{cz+d})$.
\end{physics}
\emph{Proof on paper.} For quasi-primary $v\in\lfourV_m$, $L_1L_{-1}^{j+1}v = (j+1)(2m+j)L_{-1}^jv$ (the
$\mathfrak{sl}_2$ tower). By induction on $q$, every $b\in\lfourV_q$ is a finite sum $\sum_kL_{-1}^kw_k$ with
$w_k\in\lfourV_{q-k}$ quasi-primary (\lean{tower_}\allowbreak\lean{sum_}\allowbreak\lean{exists}): given $b$, $L_1b\in\lfourV_{q-1}$ has such an
expansion, and one inverts $L_1L_{-1}$ term by term (the coefficient $(k+1)(2q-k)$ vanishes only where the
base vector is $0$) to find $c$ with $L_1(b-L_{-1}c) = 0$. \qed
\begin{status}
Conditional on \lean{hqp}; \texttt{\#print axioms}: \lean{[propext, Classical.choice, Quot.sound]}.
\end{status}

\begin{theorem}[The diagonal identity, unconditionally under $L_1\lfourV_1 = 0$]\label{l04:thm:diag}
If $\lfourvac\ne0$ and $L_1\lfourV_1 = 0$, canBilin satisfies \eqref{l04:eq:fhl} at every matching mode.
\end{theorem}
\noindent Files \leanpath{VOA/\allowbreak Core/\allowbreak InvariantFormCanonicalFHL.lean}, \leanpath{VOA/\allowbreak Core/\allowbreak InvariantFormFHL2.lean}, \leanpath{VOA/\allowbreak Core/\allowbreak InvariantFormQPMatchGen.lean}.
\begin{leancode}
def CanDiagonal (hne : W.vacuum ≠ 0) : Prop :=
  ∀ (wt p q : ℤ) (u a b : V), u ∈ W.grading wt → a ∈ W.grading p → b ∈ W.grading q →
    canBilin W hne (W.smul u (p + wt - q - 1) a) b = canRHS hne u a b wt (p + wt - q - 1)

theorem canBilin_qp_qp_adj (hne : W.vacuum ≠ 0) (hqp : ∀ v, v ∈ W.grading 1 → W.L 1 v = 0)
    {w : ℤ} {u : V} (hu : u ∈ W.grading w) (hqpu : W.L 1 u = 0)
    {p : ℤ} {a : V} (ha : a ∈ W.grading p) (hqpa : W.L 1 a = 0) (n : ℤ) (b : V) :
    canBilin W hne (W.smul u n a) b = (-1:ℂ)^w • canBilin W hne a (W.smul u (2*w - n - 2) b) := by

theorem canDiagonal_unconditional (hne : W.vacuum ≠ 0)
    (hqp : ∀ v, v ∈ W.grading 1 → W.L 1 v = 0) : CanDiagonal hne :=
  canDiagonal_holds hne hqp (fun {_p} hp {_a} ha => qpDecomp W hqp hp ha)
    (fun {_w} {_u} hu hqpu {_p} {_a} ha hqpa n b => canBilin_qp_qp_adj hne hqp hu hqpu ha hqpa n b)
\end{leancode}
\begin{dictionary}
\lean{canRHS hne u a b wt n} is the right side of \eqref{l04:eq:fhl} for canBilin (same file). The
one-line proof of \lean{canDiagonal_}\allowbreak\lean{unconditional} is shown because it is informative: it is the
assembly lemma \lean{canDiagonal_}\allowbreak\lean{holds} fed with the two genuine inputs, Theorem~\ref{l04:thm:qp} and the
quasi-primary matching identity \lean{canBilin_}\allowbreak\lean{qp_}\allowbreak\lean{qp_}\allowbreak\lean{adj} (Step 3, for quasi-primary $u$ and $a$, all
modes). ``Unconditional'' in the name means ``no hypothesis beyond \lean{hne}, \lean{hqp}''.
\end{dictionary}
\begin{physics}
The two-point function of two quasi-primaries of weights $w$ and $p$ transforms covariantly under
inversion; everything else follows by global conformal invariance, because descendants are derivatives
of quasi-primaries.
\end{physics}
\begin{status}
Conditional only on \lean{hne}, \lean{hqp}; \texttt{\#print axioms} on \lean{canDiagonal_}\allowbreak\lean{unconditional} and
\lean{canBilin_}\allowbreak\lean{qp_}\allowbreak\lean{qp_}\allowbreak\lean{adj}: \lean{[propext, Classical.choice, Quot.sound]}.
\end{status}

\begin{theorem}[FHL, Li: existence and uniqueness of the invariant form]\label{l04:thm:canonical}
Let $\lfourV$ be a VOA (of CFT type) with $\lfourvac\ne0$ and $L_1\lfourV_1 = 0$. There is a fully invariant form
with $(\lfourvac|\lfourvac) = 1$; every fully invariant normalised form equals it; it is nondegenerate if
$\lfourV$ is simple; and $u_{(1)}v = -(u|v)\lfourvac$ on $\lfourV_1$.
\end{theorem}
\noindent File \leanpath{VOA/\allowbreak Core/\allowbreak InvariantFormFHL2.lean}.
\begin{leancode}
noncomputable def canonicalInvariantForm (hne : W.vacuum ≠ 0)
    (hqp : ∀ v, v ∈ W.grading 1 → W.L 1 v = 0) : InvariantForm W :=
  canFormFHL hne (canDiagonal_unconditional hne hqp)

theorem canonicalInvariantForm_isFullyInvariant (hne : W.vacuum ≠ 0)
    (hqp : ∀ v, v ∈ W.grading 1 → W.L 1 v = 0) :
    (canonicalInvariantForm hne hqp).IsFullyInvariant :=

theorem canonicalInvariantForm_normalized (hne : W.vacuum ≠ 0)
    (hqp : ∀ v, v ∈ W.grading 1 → W.L 1 v = 0) :
    (canonicalInvariantForm hne hqp).Normalized :=

theorem invariantForm_unique_FHL (hne : W.vacuum ≠ 0)
    (hqp : ∀ v, v ∈ W.grading 1 → W.L 1 v = 0) {B : InvariantForm W}
    (hB : B.IsFullyInvariant) (hn : B.Normalized) :
    B.form = (canonicalInvariantForm hne hqp).form :=

theorem invariantForm_nondegenerate_FHL (hne : W.vacuum ≠ 0)
    (hqp : ∀ v, v ∈ W.grading 1 → W.L 1 v = 0) (hsimple : W.IsSimple) :
    InvariantForm.leftRadical (canonicalInvariantForm hne hqp) = ⊥ :=

theorem eval_oneMode_grading1_FHL (hne : W.vacuum ≠ 0)
    (hqp : ∀ v, v ∈ W.grading 1 → W.L 1 v = 0) {u v : V}
    (hu : u ∈ W.grading 1) (hv : v ∈ W.grading 1) :
    W.smul u 1 v = -(canonicalInvariantForm hne hqp).form u v • W.vacuum :=
\end{leancode}
\begin{dictionary}
\textbf{Hypotheses.} \lean{hne}: $\lfourvac\ne0$ (equivalently $\lfourV\ne0$). \lean{hqp}: $L_1\lfourV_1 = 0$.
CFT type is built into \lean{VertexOperatorAlgebra}. \lean{canFormFHL} is the negative of the
Gannon-normalised form \lean{canForm} built from canBilin and the Virasoro adjointness it inherits from
\lean{CanDiagonal} (the $u = \omega$ slice, using $L_1\omega = 0$).

\textbf{Is the definition meaningful?} \lean{canonicalInvariantForm} is a \lean{def}, so its existence
says nothing; the content is in \lean{_}\allowbreak\lean{isFullyInvariant} and \lean{_normalized}, which hold with no
further hypotheses. Normalisation excludes the zero form.

\textbf{Same, stronger or weaker than Li?} Li (JPAA 96, 1994) proves, for a vertex operator algebra
in his setting, that the space of invariant bilinear forms is linearly isomorphic to
$(\lfourV_0/L_1\lfourV_1)^*$, and that every invariant form is symmetric. The library proves the case
$\lfourV_0 = \C\lfourvac$, $L_1\lfourV_1 = 0$, in which $(\lfourV_0/L_1\lfourV_1)^* = \C$: existence and uniqueness up to
the normalisation. \textbf{It does not prove} (a) the general isomorphism with $(\lfourV_0/L_1\lfourV_1)^*$
(when $L_1\lfourV_1 = \C\lfourvac$ the space of forms is $0$; the library only has the necessity witness of
Theorem~\ref{l04:thm:virslice}), nor (b) symmetry of the form. So the Lean result is \emph{weaker} than
Li's theorem in these two respects and \emph{the same} in the case it treats. Neither omission affects
the uses in this book: $\Vnat$ has $\lfourV_1 = 0$, and nothing downstream depends on symmetry being
proved; where it is needed (Chapter~\ref{ch:L07}, \lean{form_globalTowers_eq_zero} and
\lean{form_globalTowers_eq_zero_of_ne} in \leanpath{VOA/\allowbreak Core/\allowbreak GlobalConformalTowers.lean}) it is an
explicit hypothesis \lean{hB : B.IsSymmetric}.
\end{dictionary}
\begin{physics}
A chiral CFT whose currents are quasi-primary has a canonical two-point function, unique once the
vacuum is normalised: the BPZ metric is determined by the OPE. For $\Vnat$ ($\lfourV_1 = 0$) the condition is
vacuous; for the free boson $h$ with $T = \tfrac12h^2$ it holds, and fails as soon as a background charge
is turned on.
\end{physics}
\emph{Proof on paper.} Existence: Steps 1--3 above (Theorems~\ref{l04:thm:qp}, \ref{l04:thm:diag}); the
negative of canBilin has $(\lfourvac|\lfourvac) = 1$. Uniqueness, nondegeneracy, evaluation:
Theorem~\ref{l04:thm:formbasics}(iii)--(v). \qed
\begin{status}
Conditional only on \lean{hne}, \lean{hqp}, which are the theorem's genuine hypotheses (and \lean{hsimple}
for nondegeneracy). \texttt{\#print axioms} on all six declarations: \lean{[propext, Classical.choice, Quot.sound]}.
\end{status}

\begin{warning}[title={Trap / dead end: Stale status lines in the intermediate files}]\label{l04:warn:stale}
The docstrings of \leanpath{VOA/\allowbreak Core/\allowbreak InvariantFormCanonical.lean},
\texttt{…CanonicalFHL.lean}, \texttt{…CanonicalVir.lean} and \texttt{…CanonicalPStep.lean} still carry a
status ``(half-filled circle) bottoms out at the unproved residual'' \lean{CanVirasoroAdj}, \lean{CanDiagonal} or
\lean{CanPStep} ``($\to$ Li 1994)''. These were the open residuals of the construction while it was being
built; they are now all discharged by \lean{canDiagonal_}\allowbreak\lean{unconditional} (and
\lean{canVirasoroAdj_}\allowbreak\lean{of_}\allowbreak\lean{diagonal}), as the header of \texttt{InvariantFormFHL2.lean} states. The
residual predicates remain as definitions; no theorem of this chapter depends on them as hypotheses. A
referee reading those files should not conclude that Theorem~\ref{l04:thm:canonical} is conditional.
\end{warning}

%======================================================================
\section{Li's radical theorem: simplicity is nondegeneracy}\label{l04:sec:li}
%======================================================================

\begin{theorem}[Li, Cor.~3.2]\label{l04:thm:lirad}
Let $B$ be a fully invariant normalised form on $\lfourV$. Then
\begin{enumerate}
\item $\lfourrad_LB = \lfourm(\lfourV) = \lfourrad_RB$;
\item $\lfourV$ is simple if and only if $\lfourrad_LB = 0$;
\item for homogeneous $x\in\lfourV_n$: $x\in\lfourm(\lfourV)\iff(x|y) = 0$ for all $y\in\lfourV_n$; equivalently, the
weight-$n$ part of $\lfourm(\lfourV)$ is the radical of the restriction of $B$ to $\lfourV_n\times\lfourV_n$.
\end{enumerate}
\end{theorem}
\noindent File \leanpath{VOA/\allowbreak Core/\allowbreak InvariantFormRadical.lean}.
\begin{leancode}
theorem leftRadical_eq_maxGradedIdeal {B : InvariantForm W} (hB : B.IsFullyInvariant)
    (hn : B.Normalized) :
    InvariantForm.leftRadical B = QuotVOA.maxGradedIdeal W :=

theorem leftRadical_eq_rightRadical {B : InvariantForm W} (hB : B.IsFullyInvariant)
    (hn : B.Normalized) :
    InvariantForm.leftRadical B = SingVec.rightRadical B := by

theorem isSimple_iff_leftRadical_eq_bot {B : InvariantForm W} (hB : B.IsFullyInvariant)
    (hn : B.Normalized) :
    W.IsSimple ↔ InvariantForm.leftRadical B = ⊥ := by

theorem mem_maxGradedIdeal_iff_grading {B : InvariantForm W} (hB : B.IsFullyInvariant)
    (hn : B.Normalized) {n : ℤ} {x : 𝒱} (hx : x ∈ W.grading n) :
    x ∈ QuotVOA.maxGradedIdeal W ↔ ∀ y ∈ W.grading n, B.form x y = 0 := by

def gradedForm (n : ℤ) : LinearMap.BilinForm ℂ (W.grading n) :=
  B.form.compl₁₂ (W.grading n).subtype (W.grading n).subtype

theorem ker_gradedForm {B : InvariantForm W} (hB : B.IsFullyInvariant) (hn : B.Normalized)
    (n : ℤ) :
    LinearMap.ker (gradedForm B n) = (QuotVOA.maxGradedIdeal W).comap (W.grading n).subtype := by
\end{leancode}
\begin{dictionary}
\textbf{Hypotheses.} Only full invariance and normalisation of the given $B$. In particular no
$L_1\lfourV_1 = 0$ is needed --- if some such $B$ exists the theorem applies; $\lfourvac\ne0$ follows from
normalisation (\lean{vacuum_}\allowbreak\lean{ne_}\allowbreak\lean{zero_}\allowbreak\lean{of_}\allowbreak\lean{normalized}); CFT type is built in. \textbf{Notation.}
\lean{gradedForm B n} is $B$ restricted to $\lfourV_n$ (\lean{compl₁₂} composes both slots with the inclusion
$\lfourV_n\hookrightarrow\lfourV$); \lean{comap} pulls the subspace $\lfourm(\lfourV)$ back along that inclusion. So
\lean{ker_gradedForm} is (3) as an equality of subspaces of $\lfourV_n$; \lean{maxGradedIdeal_}\allowbreak\lean{inf_}\allowbreak\lean{grading}
and \lean{maxGradedIdeal_}\allowbreak\lean{inf_}\allowbreak\lean{grading_}\allowbreak\lean{eq_}\allowbreak\lean{bot_}\allowbreak\lean{iff} restate it as $\lfourm\cap\lfourV_n = $ image of the radical, and
$\lfourm\cap\lfourV_n = 0\iff$ the level-$n$ form is nondegenerate. The \lean{_canonical} corollaries
(\lean{leftRadical_}\allowbreak\lean{canonical_}\allowbreak\lean{eq_}\allowbreak\lean{maxGradedIdeal}, \lean{isSimple_}\allowbreak\lean{iff_}\allowbreak\lean{canonical_}\allowbreak\lean{nondegenerate}, …) instantiate
$B$ at Theorem~\ref{l04:thm:canonical}, at the cost of \lean{hne}, \lean{hqp}.

\textbf{Strength.} (2) is an equivalence, strictly more than Theorem~\ref{l04:thm:formbasics}(iv). (1)
includes the equality of the two radicals, a shadow of the (unproved) symmetry of $B$. (3) is the
statement that makes simplicity \emph{computable}: membership in the abstract supremum $\lfourm(\lfourV)$ is
decided by a finite Gram matrix. \textbf{Vacuity.} The theorem is used to prove $\Vnat$ simple
(Chapter~\ref{ch:10}) and is false for the square-zero impostor of that chapter, whose analogous form has
the twisted sector in its radical.
\end{dictionary}
\begin{physics}
A chiral CFT is irreducible exactly when its Zamolodchikov metric is nondegenerate: a state with zero
two-point function against every state is a null state, and null states are exactly what generates a
decoupled ideal. Moreover null states can be found level by level, because the metric is block diagonal
in the energy: the null states at level $n$ are the kernel of the level-$n$ Gram matrix. This is the Kac
determinant method of Virasoro representation theory, here for an arbitrary chiral algebra.
\end{physics}
\emph{Proof on paper.} $\lfourrad_LB$ is an ideal (Theorem~\ref{l04:thm:formbasics}(ii)) not containing
$\lfourvac$, so $\lfourrad_LB\subseteq\lfourm$. Conversely let $I$ be a proper ideal. First $(\lfourvac|z) = 0$ for $z\in I$:
the homogeneous components of $z$ lie in $I$ (Theorem~\ref{l04:thm:graded}); those of nonzero weight are
orthogonal to $\lfourvac$, and $I\cap\lfourV_0 = 0$ by properness. Next, for homogeneous $x\in I\cap\lfourV_p$, write
$x = x_{(-1)}\lfourvac$ and expand: $(x|b) = (-1)^p\sum_j\frac1{j!}(\lfourvac|(L_1^jx)_{(2p-1-j)}b)$. Each $L_1^jx\in I$,
so $(L_1^jx)_{(m)}b\in I$ by two-sidedness (Theorem~\ref{l04:thm:twosided}), and each term vanishes. So
$I\subseteq\lfourrad_LB$. The right radical is handled the same way with the slots exchanged. (2) is (1) with
Theorem~\ref{l04:thm:maxsimple}. (3): pairings of $x\in\lfourV_n$ with other weights vanish by orthogonality. \qed
\begin{status}
Conditional only on \lean{hB}, \lean{hn}; \texttt{\#print axioms} on \lean{leftRadical_}\allowbreak\lean{eq_}\allowbreak\lean{maxGradedIdeal},
\lean{isSimple_}\allowbreak\lean{iff_}\allowbreak\lean{leftRadical_}\allowbreak\lean{eq_}\allowbreak\lean{bot}, \lean{mem_}\allowbreak\lean{maxGradedIdeal_}\allowbreak\lean{iff_}\allowbreak\lean{grading},
\lean{leftRadical_}\allowbreak\lean{eq_}\allowbreak\lean{rightRadical}, \lean{isSimple_}\allowbreak\lean{iff_}\allowbreak\lean{canonical_}\allowbreak\lean{nondegenerate}:
\lean{[propext, Classical.choice, Quot.sound]}. The source records an independent validation at
$c = -22/5$ (Lee--Yang): the level-4 radical is one-dimensional, matching $\dim V(-\frac{22}5,0)_4 = 2$ and
$\dim L(-\frac{22}5,0)_4 = 1$; that check uses \texttt{Examples} files and is not part of this file.
\end{status}

%======================================================================
\section{Gram matrices: rank computes the simple quotient}\label{l04:sec:gram}
%======================================================================

\begin{definition}[Level-$n$ Gram matrix]\label{l04:def:gram}
For a basis $(e_i)$ of $\lfourV_n$, $G_n = (B(e_i,e_j))_{i,j}$; for the canonical form, $\mathrm{canGram}$.
\end{definition}
\noindent File \leanpath{VOA/\allowbreak Core/\allowbreak GradedFormMatrix.lean}.
\begin{leancode}
noncomputable def gramMatrix (B : InvariantForm W) (n : ℤ) (e : Basis ι ℂ (W.grading n)) :
    Matrix ι ι ℂ :=
  LinearMap.BilinForm.toMatrix e (gradedForm B n)

noncomputable def canGram : Matrix ι ι ℂ := gramMatrix (canonicalInvariantForm hne hqp) n e
\end{leancode}
\begin{dictionary}
\lean{Basis ι ℂ (W.grading n)} is a basis of $\lfourV_n$ indexed by a finite type \lean{ι}
(\lean{[Fintype ι]}); \lean{toMatrix e} is the matrix of a bilinear form in that basis. \lean{canGram} takes
\lean{hne}, \lean{hqp}, \lean{n}, \lean{e} as section variables.
\end{dictionary}
\begin{physics}
The matrix of inner products of the states at energy level $n$ --- the Kac (Shapovalov) matrix of the
vacuum module, computed with the BPZ form.
\end{physics}
\begin{status}
Definitions.
\end{status}

\begin{theorem}[Rank, determinant and the simple quotient]\label{l04:thm:gram}
Let $B$ be fully invariant and normalised. Then $\det G_n\ne0\iff\lfourm\cap\lfourV_n = 0$, and
$\operatorname{rank}G_n+\dim(\lfourm\cap\lfourV_n) = \dim\lfourV_n$. For the canonical form,
$\dim L(\lfourV)_n = \operatorname{rank}\mathrm{canGram}_n$; if $\det\mathrm{canGram}_n = 0$ then $\dim L(\lfourV)_n<
\dim\lfourV_n$.
\end{theorem}
\noindent File \leanpath{VOA/\allowbreak Core/\allowbreak GradedFormMatrix.lean}.
\begin{leancode}
theorem gramMatrix_det_ne_zero_iff (hB : B.IsFullyInvariant) (hn : B.Normalized) :
    (gramMatrix B n e).det ≠ 0 ↔ maxGradedIdeal W ⊓ W.grading n = ⊥ := by

theorem rank_gramMatrix_add_finrank (hB : B.IsFullyInvariant) (hn : B.Normalized) :
    (gramMatrix B n e).rank + finrank ℂ (maxGradedIdeal W ⊓ W.grading n : Submodule ℂ 𝒱)
      = finrank ℂ (W.grading n) := by

theorem finrank_maxQuotient_grading (hne : W.vacuum ≠ 0)
    (hqp : ∀ v, v ∈ W.grading 1 → W.L 1 v = 0) (n : ℤ) (e : Basis ι ℂ (W.grading n)) :
    finrank ℂ ((QuotVOA.maxQuotientVOA W).grading n) = (canGram hne hqp n e).rank := by

theorem finrank_maxQuotient_grading_lt_of_det_eq_zero (hne : W.vacuum ≠ 0)
    (hqp : ∀ v, v ∈ W.grading 1 → W.L 1 v = 0) (n : ℤ) (e : Basis ι ℂ (W.grading n))
    (hdet : (canGram hne hqp n e).det = 0) :
    finrank ℂ ((QuotVOA.maxQuotientVOA W).grading n) < finrank ℂ (W.grading n) := by
\end{leancode}
\begin{dictionary}
\lean{finrank ℂ X} is $\dim X$; \lean{(QuotVOA.maxQuotientVOA W).grading n} is $L(\lfourV)_n$, which by
definition (\lean{maxQuotientVOA_}\allowbreak\lean{grading_}\allowbreak\lean{eq}, \lean{rfl}) is the image of $\lfourV_n$ in the quotient.
\lean{Matrix.rank} is the rank over $\C$. Finite-dimensionality of $\lfourV_n$ is a structure field, so no
finiteness hypothesis appears; the basis \lean{e} is supplied by the caller. \textbf{Rank, not
determinant.} The source stresses (and we repeat) that $\dim L(\lfourV)_n$ is the rank of the Gram matrix of
\emph{this} form; a matrix of a different form with proportional determinant has the same zero set but
can have a different rank. In particular the ``top-mode'' Shapovalov form \lean{RigidBridge.}\allowbreak\lean{shapForm} of
Chapter~\ref{ch:L07} agrees with \lean{gradedForm} only up to the sign $(-1)^n$ and only on quasi-primary
vectors (\lean{gradedForm_}\allowbreak\lean{eq_}\allowbreak\lean{signed_}\allowbreak\lean{shapForm_}\allowbreak\lean{of_}\allowbreak\lean{quasiPrimary}); at level $4$ of a Virasoro VOA no basis is
quasi-primary, and the two Gram determinants differ by the constant $21$.
\end{dictionary}
\begin{physics}
The number of physical (non-null) states at level $n$ is the rank of the Kac matrix; a zero of the Kac
determinant signals that a null state appears at that level and the character of the physical theory drops
there. This is how the characters of minimal models are computed.
\end{physics}
\emph{Proof on paper.} Theorem~\ref{l04:thm:lirad}(3) identifies $\lfourm\cap\lfourV_n$ with the kernel of the
restricted form, whose dimension is $\dim\lfourV_n-\operatorname{rank}G_n$; the quotient map
$\lfourV_n\to L(\lfourV)_n$ is surjective with that kernel. A zero determinant means a nonzero kernel. \qed
\begin{status}
Conditional only on the displayed hypotheses; \texttt{\#print axioms} on all four: \lean{[propext, Classical.choice, Quot.sound]}.
\end{status}

\begin{example}[Level two of the Virasoro vacuum module]\label{l04:ex:vir2}
In $V(c,0)$, $\lfourV_2 = \C\omega$ with $\omega = L_{-2}\lfourvac$, and $(\omega|\omega) = (\lfourvac|L_2L_{-2}\lfourvac) =
(\lfourvac|(4L_0+\tfrac c2)\lfourvac) = \tfrac c2$. So $\mathrm{canGram}_2 = (c/2)$: for $c\ne0$, $\dim L(c,0)_2 = 1$; for
$c = 0$, $\dim L(0,0)_2 = 0$, consistent with Example~\ref{l04:ex:c0}. At level $4$ the Gram matrix in the basis
$L_{-4}\lfourvac$, $L_{-2}^2\lfourvac$ has determinant proportional to $c^2(5c+22)$, whose root $c = -\frac{22}5$ is the
Lee--Yang model (Chapter~\ref{ch:L06}).
\end{example}

%======================================================================
\section{Dead ends: four general routes to full invariance}\label{l04:sec:deadends}
%======================================================================

Before the canonical construction succeeded, the project tried to prove that \emph{every}
\lean{InvariantForm} --- every Virasoro-adjoint bilinear form --- is fully invariant, reducing the claim
in four successive files to sharper and sharper residual predicates:
\begin{center}\small
\begin{tabular}{>{\raggedright\arraybackslash}p{0.27\textwidth}>{\raggedright\arraybackslash}p{0.33\textwidth}>{\raggedright\arraybackslash}p{0.3\textwidth}}
\toprule
file & reduction theorem & residual hypothesis\\
\midrule
\texttt{InvariantFormFHL} & \lean{isFullyInvariant_}\allowbreak\lean{of_}\allowbreak\lean{propagation} & \lean{hstep} (weight propagation)\\
\texttt{InvariantFormDiagonal} & \lean{isFullyInvariant_}\allowbreak\lean{of_}\allowbreak\lean{diagonal} & \lean{DiagonalFHL B}\\
\texttt{InvariantFormPStep} & \lean{isFullyInvariant_}\allowbreak\lean{of_}\allowbreak\lean{pstep} & \lean{PStep B}\\
\texttt{InvariantFormFHLClosure} & \lean{isFullyInvariant_}\allowbreak\lean{of_}\allowbreak\lean{inputs} & \lean{GoodBound B}, \lean{CreationVanish B}\\
\bottomrule
\end{tabular}
\end{center}
{\raggedright\noindent Files \leanpath{VOA/\allowbreak Core/\allowbreak InvariantFormDiagonal.lean}, \leanpath{VOA/\allowbreak Core/\allowbreak InvariantFormPStep.lean},
\leanpath{VOA/\allowbreak Core/\allowbreak InvariantFormFHLClosure.lean}.\par}
\begin{leancode}
def DiagonalFHL (B : InvariantForm W) : Prop :=
  ∀ (wt p q : ℤ) (u a b : 𝒱), u ∈ W.grading wt → a ∈ W.grading p → b ∈ W.grading q →
    B.form (W.smul u (p + wt - q - 1) a) b = fhlRHS B u a b wt (p + wt - q - 1)

theorem isFullyInvariant_of_diagonal (B : InvariantForm W) (hdiag : DiagonalFHL B) :
    B.IsFullyInvariant :=

def PStep (B : InvariantForm W) : Prop :=
  ∀ (wt q : ℤ), 0 ≤ wt → 0 ≤ q → (∀ q' : ℤ, q' < q → DiagAt B wt q') → DiagAt B wt q

theorem isFullyInvariant_of_inputs (B : InvariantForm W) (hG : GoodBound B) (hC : CreationVanish B) :
    B.IsFullyInvariant :=
\end{leancode}
All four reductions are proved and axiom-clean, but the universal statements they reduce to are false.

\begin{theorem}[Conditional delimitation]\label{l04:thm:delim}
Assume $L_1\lfourV_1 = 0$. A fully invariant form is symmetric on $\lfourV_1$. Hence, if some invariant form
$B_0$ satisfies $B_0(u_0,v_0)\ne B_0(v_0,u_0)$ for some $u_0,v_0\in\lfourV_1$, then not every invariant form is
fully invariant, and none of the four residual predicates holds universally.
\end{theorem}
\noindent File \leanpath{VOA/\allowbreak Core/\allowbreak InvariantFormDelimitation.lean}.
\begin{leancode}
theorem fullyInvariant_forces_v1_symm (B : InvariantForm W)
    (hqp : ∀ v, v ∈ W.grading 1 → W.L 1 v = 0)
    (hFI : B.IsFullyInvariant) {u v : 𝒱} (hu : u ∈ W.grading 1) (hv : v ∈ W.grading 1) :
    B.form u v = B.form v u := by

theorem general_fullyInvariant_not_a_theorem
    (hqp : ∀ v, v ∈ W.grading 1 → W.L 1 v = 0)
    {B₀ : InvariantForm W} {u₀ v₀ : 𝒱} (hu₀ : u₀ ∈ W.grading 1) (hv₀ : v₀ ∈ W.grading 1)
    (hasym : B₀.form u₀ v₀ ≠ B₀.form v₀ u₀) :
    ¬ (∀ B : InvariantForm W, B.IsFullyInvariant) := fun hAll =>
\end{leancode}
\begin{dictionary}
The four theorems \lean{general_}\allowbreak\lean{fullyInvariant_}\allowbreak\lean{not_}\allowbreak\lean{a_}\allowbreak\lean{theorem}, \lean{general_}\allowbreak\lean{diagonalFHL_}\allowbreak\lean{not_}\allowbreak\lean{a_}\allowbreak\lean{theorem},
\lean{general_}\allowbreak\lean{pstep_}\allowbreak\lean{not_}\allowbreak\lean{a_}\allowbreak\lean{theorem}, \lean{general_}\allowbreak\lean{fhlClosure_}\allowbreak\lean{not_}\allowbreak\lean{a_}\allowbreak\lean{theorem} all take the witness
\lean{B₀}, \lean{u₀}, \lean{v₀}, \lean{hasym} as hypotheses. \textbf{The library does not construct
such a witness}; the files say so. The library map's phrase ``four general invariant-form routes refuted''
must therefore be read as ``refuted conditionally''. The companion \leanpath{VOA/\allowbreak Core/\allowbreak InvariantFormTranspose.lean}
shows that the transpose $B^{\mathsf T}(a,b) = B(b,a)$ of an invariant form is again one
(\lean{InvariantForm.}\allowbreak\lean{transpose}) and that the skew-symmetry identity \lean{form_smul_skew} holds for both,
so the general toolkit cannot distinguish $B$ from $B^{\mathsf T}$
(\lean{no_}\allowbreak\lean{structural_}\allowbreak\lean{symmetry_}\allowbreak\lean{of_}\allowbreak\lean{antisymmetrize_}\allowbreak\lean{ne}).
\end{dictionary}
\begin{physics}
Virasoro covariance alone does not fix the two-point functions of a theory with several currents: one can
rotate the currents before pairing. Only covariance under the whole chiral algebra does.
\end{physics}
\emph{Proof on paper.} By \eqref{l04:eq:fhl} at $\mathrm{wt} = 1$, $n = -1$, $a = \lfourvac$, with $L_1u = 0$:
$(u|v) = -(\lfourvac|u_{(1)}v)$; by Theorem~\ref{l04:thm:v1symm}, $u_{(1)}v = v_{(1)}u$. \qed
\begin{status}
Unconditional (\lean{fullyInvariant_}\allowbreak\lean{forces_}\allowbreak\lean{v1_}\allowbreak\lean{symm}) and conditional on an unconstructed witness (the four
\lean{…_not_a_theorem}); axiom-clean.
\end{status}

\begin{warning}[title={Trap / dead end: A witness exists on paper; do not re-attempt the general route}]\label{l04:warn:deadend}
The missing witness is easy to give mathematically (this argument is not formalised). Let $\lfourV$ be the
rank-two Heisenberg VOA with currents $h^1,h^2$, $T = \frac12((h^1)^2+(h^2)^2)$, so $L_1\lfourV_1 = 0$, and let $B$ be
its canonical form, with $B(h^i,h^j) = -\delta^{ij}$. The rotation $g$ by $90^\circ$ in the $(h^1,h^2)$-plane extends
to an automorphism of $\lfourV$; it commutes with every $L_n$. Then $B_0(a,b) := B(a,gb)$ is Virasoro adjoint:
$B_0(L_na,b) = B(a,L_{-n}gb) = B(a,gL_{-n}b) = B_0(a,L_{-n}b)$. With $gh^1 = h^2$, $gh^2 = -h^1$ one gets
$B_0(h^1,h^2) = B(h^1,-h^1) = 1$ but $B_0(h^2,h^1) = B(h^2,h^2) = -1$, so $B_0$ is not symmetric on $\lfourV_1$. By
Theorem~\ref{l04:thm:delim}, the universal general-$B$ statement is false for this $\lfourV$. The moral recorded by
the project: \emph{full invariance cannot be derived from Virasoro adjointness}; the canonical-form route of
Section~\ref{l04:sec:canonical} is the only one, and the four general reductions should not be revived.
\end{warning}

%======================================================================
\section{Machine verification for this chapter}\label{l04:sec:audit}
%======================================================================

For this book (when the chapter was first written, at an earlier commit; not re-run for this edition) we ran \texttt{\#print axioms} on seventy-one declarations of the chapter's files, from a
scratch file importing the relevant modules of the built tree, with
\begin{plaincode}
export PATH="$HOME/.elan/bin:$PATH"
lake env lean -DmaxSynthPendingDepth=3 Axioms.lean
\end{plaincode}
(the flag avoids the audit trap of Section~\ref{ch15:sec:trap}: a bare
\texttt{lake env lean} drops \texttt{maxSynthPendingDepth}). Every one reported exactly
\begin{plaincode}
depends on axioms: [propext, Classical.choice, Quot.sound]
\end{plaincode}
{\raggedright including \lean{instLieAlgebraGradingOne}, \lean{inner}, \lean{isExpInner_conj}, \lean{expFun_smul},
\lean{quotientVOA}, \lean{isGradedIdeal_}\allowbreak\lean{of_}\allowbreak\lean{isIdeal}, \lean{isSimple_}\allowbreak\lean{maxQuotientVOA'},
\lean{ker_}\allowbreak\lean{eq_}\allowbreak\lean{maxGradedIdeal}, \lean{screeningKernel}, \lean{screening_}\allowbreak\lean{conformalVector},
\lean{ker_}\allowbreak\lean{eq_}\allowbreak\lean{maxGradedIdeal_}\allowbreak\lean{of_}\allowbreak\lean{screeningKernel}, \lean{generated}, \lean{tensorSVA},
\lean{not_}\allowbreak\lean{isSimple_}\allowbreak\lean{of_}\allowbreak\lean{stronglySingular_}\allowbreak\lean{FHL}, \lean{not_}\allowbreak\lean{isSimple_}\allowbreak\lean{of_}\allowbreak\lean{shapovalov_}\allowbreak\lean{FHL},
\lean{BRSTComplex.}\allowbreak\lean{cohomology}, \lean{cohomologyNearVOA}, \lean{canonicalInvariantForm_}\allowbreak\lean{isFullyInvariant},
\lean{invariantForm_}\allowbreak\lean{unique_}\allowbreak\lean{FHL}, \lean{qpDecomp}, \lean{canBilin_}\allowbreak\lean{qp_}\allowbreak\lean{qp_}\allowbreak\lean{adj},
\lean{leftRadical_}\allowbreak\lean{eq_}\allowbreak\lean{maxGradedIdeal}, \lean{isSimple_}\allowbreak\lean{iff_}\allowbreak\lean{leftRadical_}\allowbreak\lean{eq_}\allowbreak\lean{bot},
\lean{finrank_}\allowbreak\lean{maxQuotient_}\allowbreak\lean{grading}, \lean{general_}\allowbreak\lean{fullyInvariant_}\allowbreak\lean{not_}\allowbreak\lean{a_}\allowbreak\lean{theorem}. The PBW theorem \lean{VOA.pbw} appears nowhere. We did \emph{not} re-run the aggregate \texttt{lake build}. The two Lean
checks quoted in the warning on p.~\pageref{l04:warn:inn} and Remark~\ref{l04:rem:homega} were compiled the same way in a
separate scratch file and are not part of the library.\par}

\paragraph{Release audit 2026-10-27.} \texttt{Notes/\allowbreak Textbook/\allowbreak referee/\allowbreak RELEASE\_AUDIT\_FOUNDATIONS.md}.
Its findings bearing on this chapter: \lean{VertexOperatorAlgebra} is the LL/FLM VOA of CFT type, so every result here
(the Lie algebra $\lfourV_1$, quotients, simplicity, the FHL--Li invariant form) is a result about VOAs of CFT type; the
zero space is a VOA, so the nonvanishing hypotheses $\lfourvac\ne0$ in this chapter are needed; the Borcherds identity
used in the proofs is a theorem for \lean{VertexOperatorAlgebra}, not for \lean{PreVertexAlgebra}.

%======================================================================
\section{Referee's checklist}\label{l04:sec:checklist}
%======================================================================

All entries are \texttt{\#print axioms}-clean (standard axioms only, no \lean{VOA.pbw}).

\begingroup\small
\begin{longtable}{>{\raggedright\arraybackslash}p{0.28\textwidth}>{\raggedright\arraybackslash}p{0.3\textwidth}>{\raggedright\arraybackslash}p{0.2\textwidth}>{\raggedright\arraybackslash}p{0.1\textwidth}}
\toprule
Lean declaration & Intended meaning & Hypotheses / status & Check\\
\midrule\endhead
\lean{zeroMode_}\allowbreak\lean{isDerivation} & $u_{(0)}$ is a derivation of every product & none; unconditional & Thm.~\ref{l04:thm:deriv}\\
\lean{instLieAlgebraGradingOne} & $\lfourV_1$ is a Lie algebra, $[u,v] = u_{(0)}v$ & none; unconditional & Thm.~\ref{l04:thm:lie}\\
\lean{griessMul_}\allowbreak\lean{comm_}\allowbreak\lean{of_}\allowbreak\lean{grading1_}\allowbreak\lean{eq_}\allowbreak\lean{bot} & Griess product commutative & $\lfourV_1 = 0$ & Def.~\ref{l04:def:griess}\\
\lean{Aut}, \lean{Group (Aut W)} & automorphism group & definition & Def.~\ref{l04:def:aut}\\
\lean{Aut.map_grading} & $\alpha(\lfourV_n)\subseteq\lfourV_n$ & none; unconditional & Thm.~\ref{l04:thm:autgrading}\\
\lean{VOASubalgebra}, \lean{VOASubalgebra.}\allowbreak\lean{toVOA}, \lean{toVOA_centralCharge} & graded subspace containing $\lfourvac,\omega$, closed under products; it is a VOA of the same central charge & definition; $\omega\in S$ required; unconditional & Def.~\ref{l04:def:subalg}, Thm.~\ref{l04:thm:subalg}\\
\lean{expFun_smul} & $e^D$ of a derivation is a homomorphism & $D$ loc.\ nilpotent, Leibniz & Thm.~\ref{l04:thm:expsmul}\\
\lean{inner} & $\exp(u_{(0)})\in\Aut$ & $u_{(0)}$ loc.\ nilp.; \lean{homega} redundant & Def.~\ref{l04:def:inner}, Rem.~\ref{l04:rem:homega}\\
\lean{isExpInner_conj}, \lean{aut_conj_expFun} & $\alpha e^{u_{(0)}}\alpha^{-1} = e^{(\alpha u)_{(0)}}$ & none (former redundant \lean{hαu} removed) & Thm.~\ref{l04:thm:conj}\\
\lean{Inn}, \lean{Inn_normal} & inner automorphisms normal & \textbf{normal by construction}; weak predicate & p.~\pageref{l04:warn:inn}\\
\lean{IsIdeal}, \lean{IsSimple} & ideal, simple VOA & left ideals = two-sided & Def.~\ref{l04:def:simple}\\
\lean{absorbsLeft_}\allowbreak\lean{of_}\allowbreak\lean{isIdeal} & ideals are two-sided & none & Thm.~\ref{l04:thm:twosided}\\
\lean{isGradedIdeal_}\allowbreak\lean{of_}\allowbreak\lean{isIdeal} & ideals are graded & none & Thm.~\ref{l04:thm:graded}\\
\lean{quotientVOA}, \lean{quotientVOA'} & $\lfourV/I$ is a VOA & $I$ ideal; properness not assumed & Def.~\ref{l04:def:quot}\\
\lean{maxGradedIdeal}, \lean{isSimple_}\allowbreak\lean{maxQuotientVOA'} & $\lfourV/\lfourm$ is simple & $\lfourvac\ne0$ & Thm.~\ref{l04:thm:maxsimple}\\
\lean{IsVOAHom} & homomorphism & $\omega$ not required to be preserved & Def.~\ref{l04:def:hom}\\
\lean{IsVOAHom.}\allowbreak\lean{ker_}\allowbreak\lean{eq_}\allowbreak\lean{maxGradedIdeal} & surjection onto simple has kernel $\lfourm$ & surjective, target simple & Thm.~\ref{l04:thm:transfer}\\
\lean{VOASubalgebra.}\allowbreak\lean{generated} & subalgebra generated by $S$ & \textbf{always contains $\omega$} & Def.~\ref{l04:def:generated}\\
\lean{screening_}\allowbreak\lean{isDerivation}, \lean{screeningKernel} & $\ker a_{(0)}$ is a subVOA & $a\in\lfourV_1$ (no primarity) & Thm.~\ref{l04:thm:scrder}, Def.~\ref{l04:def:scrkernel}\\
\lean{Screening.}\allowbreak\lean{ker_}\allowbreak\lean{eq_}\allowbreak\lean{maxGradedIdeal_}\allowbreak\lean{of_}\allowbreak\lean{screeningKernel} & screening-kernel image is $L(\lfourV)$ & image $=\ker Q_a$, kernel simple & Thm.~\ref{l04:thm:scrtransfer}\\
\lean{superScreeningKernel} & super analogue & Leibniz, $Q\omega = 0$ \textbf{assumed} & Rem.~\ref{l04:rem:super}\\
\lean{IsSingularOfWeight}, \lean{IsStronglySingularOfWeight} & singular vectors & $s\ne0$, weight $>0$ & Def.~\ref{l04:def:singular}\\
\lean{not_}\allowbreak\lean{isSimple_}\allowbreak\lean{of_}\allowbreak\lean{stronglySingular_}\allowbreak\lean{FHL} & strongly singular $\Rightarrow$ not simple & $\lfourvac\ne0$, $L_1\lfourV_1 = 0$ & Thm.~\ref{l04:thm:notsimple}\\
\lean{not_}\allowbreak\lean{isSimple_}\allowbreak\lean{of_}\allowbreak\lean{shapovalov_}\allowbreak\lean{FHL} & singular + Shapovalov $\Rightarrow$ not simple & same + pairing condition & Thm.~\ref{l04:thm:notsimple}\\
\lean{tensorSVA} & tensor product of vertex superalgebras & \lean{ParityCompatible}, \lean{TensorGradingHyp}; \textbf{no $\omega$} & Thm.~\ref{l04:thm:tensor}, p.~\pageref{l04:warn:tensor}\\
\lean{BRSTComplex.}\allowbreak\lean{cohomology} & $\ker Q/\operatorname{im}Q$ is a vertex superalgebra & a \lean{BRSTComplex}; nontriviality not claimed & Thm.~\ref{l04:thm:brst}\\
\lean{cohomologyNearVOA} & conformal structure descends & $Q\omega = 0$ & Thm.~\ref{l04:thm:brstconf}\\
\lean{InvariantForm}, \lean{IsFullyInvariant} & Virasoro-adjoint form; FHL invariance & \textbf{two layers}; check which is assumed & Defs.~\ref{l04:def:invform}, \ref{l04:def:fullinv}\\
\lean{invariantForm_}\allowbreak\lean{unique} & normalised fully invariant forms coincide & none & Thm.~\ref{l04:thm:formbasics}\\
\lean{qpDecomp} & $\lfourV_p = \lfourQP_p+L_{-1}\lfourV_{p-1}$ & $L_1\lfourV_1 = 0$ & Thm.~\ref{l04:thm:qp}\\
\lean{canonicalInvariantForm} (+ \lean{_}\allowbreak\lean{isFullyInvariant}, \lean{_normalized}) & FHL/Li canonical form exists & $\lfourvac\ne0$, $L_1\lfourV_1 = 0$; symmetry \textbf{not proved} & Thm.~\ref{l04:thm:canonical}\\
\lean{invariantForm_}\allowbreak\lean{unique_}\allowbreak\lean{FHL} & uniqueness & same & Thm.~\ref{l04:thm:canonical}\\
\lean{virSlice_vacuum} & $L_1\lfourV_1 = 0$ is necessary & none & Thm.~\ref{l04:thm:virslice}\\
\lean{LiRadical.}\allowbreak\lean{leftRadical_}\allowbreak\lean{eq_}\allowbreak\lean{maxGradedIdeal} & radical $=\lfourm(\lfourV)$ & $B$ fully inv., normalised & Thm.~\ref{l04:thm:lirad}\\
\lean{LiRadical.}\allowbreak\lean{isSimple_}\allowbreak\lean{iff_}\allowbreak\lean{leftRadical_}\allowbreak\lean{eq_}\allowbreak\lean{bot} & simple $\iff$ nondegenerate & same & Thm.~\ref{l04:thm:lirad}\\
\lean{LiRadical.}\allowbreak\lean{mem_}\allowbreak\lean{maxGradedIdeal_}\allowbreak\lean{iff_}\allowbreak\lean{grading} & level-$n$ criterion & same & Thm.~\ref{l04:thm:lirad}\\
\lean{GramMat.}\allowbreak\lean{finrank_}\allowbreak\lean{maxQuotient_}\allowbreak\lean{grading} & $\dim L(\lfourV)_n = \operatorname{rank}\mathrm{canGram}_n$ & $\lfourvac\ne0$, $L_1\lfourV_1 = 0$, basis & Thm.~\ref{l04:thm:gram}\\
\lean{general_}\allowbreak\lean{fullyInvariant_}\allowbreak\lean{not_}\allowbreak\lean{a_}\allowbreak\lean{theorem} & general route false & \textbf{conditional on unconstructed witness} & Thm.~\ref{l04:thm:delim}, p.~\pageref{l04:warn:deadend}\\
\bottomrule
\end{longtable}
\endgroup

%======================================================================
\section{Notes and references}\label{l04:sec:notes}
%======================================================================

The invariant bilinear form and its mode formula are due to I.~B.~Frenkel, Y.-Z.~Huang and
J.~Lepowsky, \emph{On axiomatic approaches to vertex operator algebras and modules}, Memoirs of the
American Mathematical Society 104 (1993), no.~494 (Section~5.2 for contragredient modules and invariant
forms, Prop.~5.3.6 for symmetry, Section~3 for tensor products). The existence-and-uniqueness theorem, the
identification of the space of invariant forms with $(\lfourV_0/L_1\lfourV_1)^*$, and the corollary relating the
radical to the maximal ideal are in H.-S.~Li, \emph{Symmetric invariant bilinear forms on vertex operator
algebras}, Journal of Pure and Applied Algebra 96 (1994) 279--297. The library's docstrings cite these, and
the textbook account of T.~Gannon, \emph{Moonshine beyond the Monster}, Cambridge University Press (2006),
§5.1.3 and §5.2.1, for the conventions (equations (5.1.11)--(5.1.12), (5.2.1)--(5.2.4)).

General structure theory (ideals, quotients, simplicity, the $\lfourV_1$ Lie algebra, skew-symmetry) follows
J.~Lepowsky and H.~Li, \emph{Introduction to Vertex Operator Algebras and Their Representations},
Progress in Mathematics 227, Birkh\"auser (2004), Chapter~3; I.~Frenkel, J.~Lepowsky and A.~Meurman,
\emph{Vertex Operator Algebras and the Monster}, Academic Press (1988); and V.~Kac, \emph{Vertex
Algebras for Beginners}, 2nd ed., University Lecture Series 10, AMS (1998), §4.3 for tensor products. The
automorphism groups and inner automorphisms of VOAs, and the role of $\lfourV_1$, are discussed in C.~Dong and
G.~Mason, \emph{Rational vertex operator algebras and the effective central charge}, International
Mathematics Research Notices 2004, and in Dong--Mason's work on
the structure of $\lfourV_1$ (\emph{Integrability of $C_2$-cofinite vertex operator algebras}, IMRN 2006), where $\lfourV_1$ is shown to be reductive under rationality assumptions --- a result the
library does not formalise.

Screening operators and free-field realisations: B.~Feigin and D.~Fuchs, \emph{Representations of the
Virasoro algebra}, in Representations of Lie groups and related topics, Gordon and Breach (1990);
Vl.~S.~Dotsenko and V.~A.~Fateev, Nuclear Physics B 240 (1984) 312; D.~Adamovi\'c, \emph{Realizations of
simple affine vertex algebras and their modules: the cases $\widehat{sl(2)}$ and $\widehat{osp(1,2)}$},
Communications in Mathematical Physics 366 (2019) 1025--1067. BRST cohomology and the
Drinfeld--Sokolov reduction: B.~Feigin and E.~Frenkel, \emph{Quantization of the Drinfeld--Sokolov
reduction}, Physics Letters B 246 (1990) 75--81; V.~Kac, S.-S.~Roan and M.~Wakimoto, \emph{Quantum reduction
for affine superalgebras}, Communications in Mathematical Physics 241 (2003) 307--342
(arXiv:math-ph/0302015); E.~Frenkel and D.~Ben-Zvi, \emph{Vertex Algebras and Algebraic Curves}, 2nd ed.,
Mathematical Surveys and Monographs 88, AMS (2004), Chapter~15. The physical picture of the BPZ pairing
is that of A.~A.~Belavin, A.~M.~Polyakov and A.~B.~Zamolodchikov, Nuclear Physics B 241 (1984) 333--380,
and of P.~Di~Francesco, P.~Mathieu and D.~S\'en\'echal, \emph{Conformal Field Theory}, Springer (1997),
Chapters~6--7. The singular-vector criterion through invariant forms and the Gram-rank formula are the
library's own packaging of these classical facts; the Kac determinant itself is the subject of
Chapter~\ref{ch:L06}.

\endgroup
